% Les échanges du Cameroun avec l’extérieur — Géographie Tle A — Cours signé OnBuch+
% Programme officiel MINESEC. Compiler avec : tectonic N16-echanges-cameroun-avec-exterieur.tex
\newcommand{\DOCMATIERE}{Géographie}
\newcommand{\DOCNIVEAU}{Tle A}
\newcommand{\DOCMODULE}{Module 2 : La libéralisation des échanges (la mondialisation)}
\newcommand{\DOCLECON}{N16}
\newcommand{\DOCTITRE}{Les échanges du Cameroun avec l’extérieur}
% OnBuch+ — préambule « cours signés » ENRICHI (compilé avec tectonic / XeLaTeX)
% Système visuel « Pop » d'OnBuch+ : fond crème (jamais blanc), bordures encre
% épaisses, ombres « dures » décalées, titres Archivo Black en capitales.
% Version MATHÉMATIQUES : graphes (pgfplots/tikz), boîtes pédagogiques
% (définition, propriété, méthode, attention, le savais-tu, exercice résolu,
% exercice, TP, bilan), page de garde et signature OnBuch+ ancrée en bas.
%
% Macros attendues dans main.tex AVANT \input{preamble} :
%   \DOCMATIERE  \DOCNIVEAU  \DOCMODULE  \DOCLECON (numéro)  \DOCTITRE
\documentclass[11pt]{article}

\usepackage{fontspec}
\usepackage[french]{babel}
\usepackage[a4paper,margin=2cm,top=3.4cm,bottom=2.8cm,headheight=1.4cm,footskip=1.3cm]{geometry}
\usepackage{amsmath,amssymb,amsfonts}
\usepackage{mathtools} % \xmapsto, \coloneqq... souvent utilisés par les modèles
\usepackage{cancel} % simplifications barrées (bilans de forces)
% \abs{z} (module, valeur absolue) et \norm{u} (norme), utilisés par les modèles
\providecommand{\abs}{}\let\abs\relax\DeclarePairedDelimiter{\abs}{\lvert}{\rvert}
\providecommand{\norm}{}\let\norm\relax\DeclarePairedDelimiter{\norm}{\lVert}{\rVert}
\usepackage[table]{xcolor} % [table] : fournit aussi \rowcolors (alternance de lignes)
\usepackage{pagecolor}
\usepackage{tikz}
\usetikzlibrary{arrows.meta,positioning,calc,shapes.geometric,shapes.misc,shapes.arrows,decorations.pathmorphing,decorations.pathreplacing,decorations.markings,patterns,shadows,mindmap,trees,fit,backgrounds,matrix,intersections,angles,quotes}
\usepackage{pgfplots}
\usepackage[version=4]{mhchem} % équations nucléaires \ce{^{235}_{92}U}
\usepackage[european,siunitx]{circuitikz} % schémas électriques
\pgfplotsset{compat=1.18}
\usepgfplotslibrary{fillbetween,groupplots} % aires entre courbes (calcul intégral)
\usepackage{enumitem}
\usepackage{fancyhdr}
% [most] charge toutes les bibliothèques tcolorbox (listings, minted, raster,
% documentation...) et gonfle inutilement la mémoire TeX sur un document long
% (~300 boîtes) : on ne charge que ce qui est réellement utilisé.
\usepackage[skins,breakable]{tcolorbox}
\usepackage{graphicx}
\graphicspath{{/home/user/t-t/geo-onbuch/images/}}
\usepackage{colortbl}
\usepackage{booktabs}
\usepackage{tabularx}
\usepackage{multirow}
\usepackage{diagbox} % case d'en-tête diagonale des tableaux à double entrée
% Colonnes extensibles centrée / gauche / droite (C, L, R), souvent utilisées par les modèles
\newcolumntype{C}{>{\centering\arraybackslash}X}
\newcolumntype{L}{>{\raggedright\arraybackslash}X}
\newcolumntype{R}{>{\raggedleft\arraybackslash}X}
\usepackage{array}
\usepackage{multicol}
\usepackage{siunitx}
% parse-numbers=false : accepte aussi \SI{2,5 \times 10^{-3}}{...}
\sisetup{locale=FR,output-decimal-marker={,},group-minimum-digits=4,parse-numbers=false}
\DeclareSIUnit\ohm{\ensuremath{\symup{\Omega}}} % U+2126 absent de la police
\DeclareSIUnit\division{div} % réglages d'oscilloscope (ms/div, V/div)
\DeclareSIUnit\tr{tr} % tours (vitesse de rotation, ex. \SI{1500}{\tr\per\minute})
\DeclareSIUnit\molar{M} % concentration molaire (ex. \SI{3}{\molar}, \SI{1}{\micro\molar})
\DeclareSIUnit\an{an} % année (le modèle confond parfois avec \year, un compteur TeX) — ex. \SI{5}{\kilo\meter\per\an}
\DeclareSIUnit\FCFA{FCFA} % franc CFA (ex. \SI{500}{\FCFA\per\kilo\gram})
\DeclareSIUnit\degreeBrix{\textdegree Brix} % teneur en sucre d'un sirop/jus (réfractométrie)
\DeclareSIUnit\month{mois} % le modèle utilise parfois \month, un compteur TeX, comme unité de durée
\DeclareSIUnit\microliter{\micro\liter} % le modèle écrit parfois \microliter en un seul token
\DeclareSIUnit\million{millions} % dénombrement (ex. \SI{15}{\million\per\mL} spermatozoïdes)
\usepackage{qrcode}
% \degree : commande fréquemment utilisée par les modèles pour un angle ou une
% température en dehors de \SI{}{} (non fournie par les paquets chargés).
\providecommand{\degree}{^{\circ}}
% \qed (fin de démonstration), utilisé par les modèles sans amsthm
\providecommand{\qed}{\hfill$\square$}
% \barre : double barre des valeurs interdites dans un tableau de variations (tkz-tab)
\providecommand{\barre}{\ensuremath{\|}}
% environnement proof (amsthm non chargé) utilisé par les modèles
\ifdefined\proof\else\newenvironment{proof}[1][Démonstration]{\par\noindent\textit{#1.}\ }{\qed\par}\fi
\usepackage{lastpage}
\usepackage{hyperref}
\hypersetup{hidelinks}

% ============================================================
% Palette « Pop » OnBuch+ — mêmes hex que la classe P (plus_ui.dart)
% ============================================================
\definecolor{poporange}{HTML}{F59321}
\definecolor{popdark}{HTML}{E07A0C}
\definecolor{popcream}{HTML}{FFF6E8}
\definecolor{popink}{HTML}{1C1714}
\definecolor{poppurple}{HTML}{7A5AE0}
\definecolor{popgreen}{HTML}{2E9E6A}
\definecolor{poppink}{HTML}{E0457B}
\definecolor{popblue}{HTML}{2F7DE1}
\definecolor{popgold}{HTML}{C98A12}
\definecolor{popmuted}{HTML}{8A7F76}
\definecolor{popline}{HTML}{EADFD2}
\definecolor{popgray}{HTML}{8A7F76}
\colorlet{popgrey}{popgray}
% Alias tolérants (le modèle nomme parfois les couleurs différemment)
\colorlet{popwhite}{white}
\colorlet{popblack}{popink}
\colorlet{popred}{poppink}
\colorlet{popyellow}{popgold}
% Teintes claires (fonds de tableaux/encadrés) — le modèle nomme parfois
% les couleurs avec un suffixe L (ex. popblueL) sans les avoir définies.
\colorlet{poporangeL}{poporange!20}
\colorlet{popdarkL}{popdark!20}
\colorlet{poppurpleL}{poppurple!20}
\colorlet{popgreenL}{popgreen!20}
\colorlet{poppinkL}{poppink!20}
\colorlet{popblueL}{popblue!20}
\colorlet{popgoldL}{popgold!20}
\colorlet{popredL}{popred!20}
\colorlet{popgrayL}{popgray!20}
% Teintes claires pour fonds de tableaux / zones
\colorlet{popblueL}{popblue!10!white}
\colorlet{poporangeL}{poporange!14!white}
\colorlet{popgreenL}{popgreen!12!white}
\colorlet{poppurpleL}{poppurple!12!white}
\colorlet{poppinkL}{poppink!12!white}
\colorlet{popgoldL}{popgold!14!white}

\pagecolor{popcream}

% ============================================================
% Typographie
% ============================================================
\defaultfontfeatures{Ligatures=TeX}
\setmainfont{PlusJakartaSans-Medium.ttf}[Path=../../../fonts/,BoldFont=PlusJakartaSans-Bold.ttf]
% Maths en sans-serif (Fira Math) avec chiffres et lettres droites de Plus
% Jakarta, pour que les formules se fondent dans le texte.
\usepackage[math-style=ISO,bold-style=ISO]{unicode-math}
\setmathfont{FiraMath-Regular.otf}[Path=../../../fonts/]
\setmathfont{PlusJakartaSans-Medium.ttf}[Path=../../../fonts/,range={up/num,up/{Latin,latin},\mathup}]
\usepackage{icomma} % virgule décimale sans espace en mode math
% Caractères mathématiques Unicode absents des polices de texte (Plus Jakarta,
% Archivo Black) : rendus par leur équivalent mathématique, en texte comme en
% formule — sinon ils s'affichent en carré vide (ex. « Arithmétique dans ℤ »).
\usepackage{newunicodechar}
\newunicodechar{ℝ}{\ensuremath{\mathbb{R}}}
\newunicodechar{ℤ}{\ensuremath{\mathbb{Z}}}
\newunicodechar{ℂ}{\ensuremath{\mathbb{C}}}
\newunicodechar{ℕ}{\ensuremath{\mathbb{N}}}
\newunicodechar{ℚ}{\ensuremath{\mathbb{Q}}}
\newunicodechar{𝔻}{\ensuremath{\mathbb{D}}}
\newunicodechar{α}{\ensuremath{\alpha}}
\newunicodechar{β}{\ensuremath{\beta}}
\newunicodechar{γ}{\ensuremath{\gamma}}
\newunicodechar{δ}{\ensuremath{\delta}}
\newunicodechar{ε}{\ensuremath{\varepsilon}}
\newunicodechar{θ}{\ensuremath{\theta}}
\newunicodechar{λ}{\ensuremath{\lambda}}
\newunicodechar{μ}{\ensuremath{\mu}}
\newunicodechar{σ}{\ensuremath{\sigma}}
\newunicodechar{τ}{\ensuremath{\tau}}
\newunicodechar{φ}{\ensuremath{\varphi}}
\newunicodechar{ψ}{\ensuremath{\psi}}
\newunicodechar{ω}{\ensuremath{\omega}}
\newunicodechar{Σ}{\ensuremath{\Sigma}}
% majuscules produites par \MakeUppercase dans les titres (α-aminés -> Α)
\newunicodechar{Α}{\ensuremath{\alpha}}
\newunicodechar{Β}{\ensuremath{\beta}}
\newunicodechar{Γ}{\ensuremath{\Gamma}}
\newunicodechar{Δ}{\ensuremath{\Delta}}
\newunicodechar{Ε}{\ensuremath{\varepsilon}}
\newunicodechar{Θ}{\ensuremath{\Theta}}
\newunicodechar{Λ}{\ensuremath{\Lambda}}
\newunicodechar{Μ}{\ensuremath{\mu}}
\newunicodechar{Τ}{\ensuremath{\tau}}
\newunicodechar{Φ}{\ensuremath{\Phi}}
\newunicodechar{Ψ}{\ensuremath{\Psi}}
\newunicodechar{Ω}{\ensuremath{\Omega}}
\newunicodechar{↦}{\ensuremath{\mapsto}}
\newunicodechar{∈}{\ensuremath{\in}}
\newunicodechar{∉}{\ensuremath{\notin}}
\newunicodechar{∧}{\ensuremath{\wedge}}
\newunicodechar{⇔}{\ensuremath{\Leftrightarrow}}
\newunicodechar{⟺}{\ensuremath{\Longleftrightarrow}}
\newunicodechar{⇒}{\ensuremath{\Rightarrow}}
\newunicodechar{⟹}{\ensuremath{\Longrightarrow}}
\newunicodechar{∘}{\ensuremath{\circ}}
\newunicodechar{∪}{\ensuremath{\cup}}
\newunicodechar{∩}{\ensuremath{\cap}}
\newunicodechar{≡}{\ensuremath{\equiv}}
\newunicodechar{⊂}{\ensuremath{\subset}}
\newunicodechar{⊥}{\ensuremath{\perp}}
\newunicodechar{∃}{\ensuremath{\exists}}
\newunicodechar{∀}{\ensuremath{\forall}}
\newunicodechar{∛}{\ensuremath{\sqrt[3]{\,}}}
\newunicodechar{∜}{\ensuremath{\sqrt[4]{\,}}}
\newunicodechar{⃗}{\ensuremath{{}^{\to}}}
\newunicodechar{ⁿ}{\ifmmode^{n}\else\textsuperscript{\ensuremath{n}}\fi}
\newunicodechar{ˣ}{\ifmmode^{x}\else\textsuperscript{\ensuremath{x}}\fi}
\newunicodechar{ᵇ}{\ifmmode^{b}\else\textsuperscript{\ensuremath{b}}\fi}
\newunicodechar{ʳ}{\ifmmode^{r}\else\textsuperscript{\ensuremath{r}}\fi}
\newunicodechar{ᵘ}{\ifmmode^{u}\else\textsuperscript{\ensuremath{u}}\fi}
\newunicodechar{ʸ}{\ifmmode^{y}\else\textsuperscript{\ensuremath{y}}\fi}
\newunicodechar{ᵃ}{\ifmmode^{a}\else\textsuperscript{\ensuremath{a}}\fi}
\newunicodechar{ᵉ}{\ifmmode^{e}\else\textsuperscript{\ensuremath{e}}\fi}
\newunicodechar{ᵏ}{\ifmmode^{k}\else\textsuperscript{\ensuremath{k}}\fi}
\newunicodechar{ᵖ}{\ifmmode^{p}\else\textsuperscript{\ensuremath{p}}\fi}
\newunicodechar{ᴺ}{\ifmmode^{N}\else\textsuperscript{\ensuremath{N}}\fi}
\newunicodechar{ᶜ}{\ifmmode^{c}\else\textsuperscript{\ensuremath{c}}\fi}
\newunicodechar{ʰ}{\ifmmode^{h}\else\textsuperscript{\ensuremath{h}}\fi}
\newunicodechar{ᵠ}{\ifmmode^{q}\else\textsuperscript{\ensuremath{q}}\fi}
\newunicodechar{ₙ}{\ifmmode_{n}\else\textsubscript{\ensuremath{n}}\fi}
\newunicodechar{ₐ}{\ifmmode_{a}\else\textsubscript{\ensuremath{a}}\fi}
\newunicodechar{ᵢ}{\ifmmode_{i}\else\textsubscript{\ensuremath{i}}\fi}
\newunicodechar{ᵣ}{\ifmmode_{r}\else\textsubscript{\ensuremath{r}}\fi}
\newunicodechar{ₖ}{\ifmmode_{k}\else\textsubscript{\ensuremath{k}}\fi}
\newunicodechar{ᵦ}{\ifmmode_{\beta}\else\textsubscript{\ensuremath{\beta}}\fi}
\newunicodechar{ₓ}{\ifmmode_{x}\else\textsubscript{\ensuremath{x}}\fi}
\newfontfamily\popdisplay{ArchivoBlack-Regular.ttf}[Path=../../../fonts/]
\newfontfamily\poplogo{SpaceGrotesk-Bold.ttf}[Path=../../../fonts/]
\newfontfamily\popsemi{PlusJakartaSans-SemiBold.ttf}[Path=../../../fonts/]
\newfontfamily\popxbold{PlusJakartaSans-ExtraBold.ttf}[Path=../../../fonts/]
\newfontfamily\popxboldls{PlusJakartaSans-ExtraBold.ttf}[Path=../../../fonts/,LetterSpace=3.0]

\setlist[enumerate]{leftmargin=1.7em,itemsep=5pt,topsep=4pt}
\setlist[itemize]{leftmargin=1.7em,itemsep=4pt,topsep=4pt,label={\color{poporange}\textbullet}}
\setlength{\parindent}{0pt}
\setlength{\parskip}{5pt}
\renewcommand{\arraystretch}{1.25}

% pgfplots : style Pop par défaut
\pgfplotsset{
  every axis/.append style={axis line style={popink,line width=1pt},tick style={popink},
    grid style={popline,line width=0.5pt},label style={font=\small},tick label style={font=\footnotesize}},
  popcurve/.style={poporange,line width=1.8pt,no markers,smooth},
}
% Même style disponible dans un \draw TikZ simple (hors pgfplots)
\tikzset{popcurve/.style={poporange,line width=1.8pt}}
\tikzset{popgrid/.style={popline,very thin}}
\colorlet{popcurve}{poporange} % le nom du style est parfois utilisé comme couleur

% ============================================================
% Pill / squiggle
% ============================================================
\newcommand{\poppill}[3]{%
  \tikz[baseline=-0.55ex]\node[draw=popink,line width=1.2pt,rounded corners=2.6pt,%
    fill=#2,inner xsep=6.5pt,inner ysep=3pt]{%
    \popxboldls\fontsize{7}{7}\selectfont\color{#3}\MakeUppercase{#1}};%
}
\newcommand{\popsquiggle}[1]{%
  \tikz[baseline=-0.4ex]{
    \draw[#1,line width=2.6pt,line cap=round]
      plot [smooth] coordinates {(0,0.09) (0.34,0.01) (0.68,0.21) (1.02,0.01) (1.36,0.10)};
  }%
}
% Mot-clé surligné (marqueur orange sous le texte)
\newcommand{\cle}[1]{{\popxbold\color{popdark}#1}}

% ============================================================
% En-tête / pied
% ============================================================
\pagestyle{fancy}
\fancyhf{}
\renewcommand{\headrulewidth}{0pt}
\renewcommand{\footrulewidth}{0pt}
\lhead{\raisebox{-0.15ex}{\poplogo\fontsize{13}{13}\selectfont\color{popink}OnBuch\color{poporange}+}}
\rhead{\poppill{\DOCMATIERE\ \textbullet\ \DOCNIVEAU\ \textbullet\ Leçon \DOCLECON}{white}{popink}}
\lfoot{\popxbold\fontsize{7.6}{7.6}\selectfont\color{popmuted}\MakeUppercase{Cours signé OnBuch+}\ \textbullet\ {\popsemi\DOCTITRE}}
\rfoot{\tikz[baseline=-0.6ex]\node[rounded corners=8.5pt,fill=popink,minimum height=17pt,minimum width=17pt,inner xsep=5pt,inner sep=0]{\popxbold\fontsize{8}{8}\selectfont\color{popcream}\thepage\,/\,\pageref*{LastPage}};}

% ============================================================
% Titres
% ============================================================
\newcommand{\chaptitre}[2]{%
  \begin{tcolorbox}[enhanced,colback=poporange,colframe=popink,boxrule=2.6pt,arc=11pt,
      drop shadow=popink,left=15pt,right=15pt,top=12pt,bottom=11pt,before skip=3pt,after skip=18pt]
    \poppill{#2}{popink}{popcream}\par\vspace{9pt}
    {\raggedright\hyphenpenalty=10000\exhyphenpenalty=10000\popdisplay\fontsize{22}{24}\selectfont\color{popcream}\MakeUppercase{#1}\par}
  \end{tcolorbox}
}
% Section numérotée : pastille orange avec le numéro + titre Archivo
\newcounter{popsec}
\newcommand{\coursec}[1]{%
  \stepcounter{popsec}\setcounter{popsub}{0}%
  \par\vspace{16pt}\noindent
  \begin{minipage}{\linewidth}
  \tikz[baseline=-0.7ex]\node[draw=popink,line width=1.6pt,fill=poporange,rounded corners=4pt,minimum size=22pt,inner sep=0]{\popdisplay\fontsize{12}{12}\selectfont\color{popcream}\thepopsec};%
  \hspace{8pt}{\popdisplay\fontsize{15}{17}\selectfont\color{popink}\MakeUppercase{#1}}\par
  \vspace{4pt}\noindent{\color{poporange}\rule{36pt}{3.6pt}}
  \end{minipage}\par\nopagebreak\vspace{8pt}
}
\newcounter{popsub}
\newcommand{\courssub}[1]{%
  \stepcounter{popsub}%
  \par\vspace{9pt}\noindent{\popxbold\fontsize{11.5}{13}\selectfont\color{popink}{\color{poporange}\thepopsec.\thepopsub}\hspace{6pt}#1}\par\nopagebreak\vspace{4pt}
}

% ============================================================
% Boîtes pédagogiques — même recette : fond blanc, bordure épaisse colorée,
% ombre dure encre, badge plein. Titre optionnel [#1].
% ============================================================
\tcbset{popbase/.style={breakable,enhanced,colback=white,boxrule=2.3pt,arc=9pt,
  shadow={3pt}{-3pt}{0pt}{popink},left=14pt,right=14pt,top=11pt,bottom=11pt,before skip=11pt,after skip=15pt,
  fontupper=\popsemi\fontsize{10.6}{14.5}\selectfont\color{popink},
  fontlower=\popsemi\fontsize{10.6}{14.5}\selectfont\color{popink}}}
% Coche dessinée (le glyphe ✓ n'existe pas dans les polices du document)
\newcommand{\popcheck}[1]{\tikz[baseline=-0.5ex]\draw[#1,line width=1.6pt,line cap=round,line join=round](-0.13,0.0)--(-0.04,-0.09)--(0.13,0.1);}
\providecommand{\coche}{\popcheck{popgreen}}
\providecommand{\resultat}[1]{\par\smallskip\noindent\fcolorbox{popgreen}{popgreenL}{\parbox{\dimexpr\linewidth-2\fboxsep-2\fboxrule\relax}{\textbf{Résultat.} #1}}\par\smallskip}
\newcommand{\popboxhead}[3]{\poppill{#1}{#2}{white}\ifx\relax#3\relax\else\hspace{6pt}{\popxbold\fontsize{10.6}{12}\selectfont\color{popink}#3}\fi\par\vspace{7pt}}

\newtcolorbox{aretenir}[1][]{popbase,colframe=popgreen,before upper={\popboxhead{À retenir}{popgreen}{#1}}}
\newtcolorbox{exemplebox}[1][]{popbase,colframe=poporange,before upper={\popboxhead{Exemple}{poporange}{#1}}}
\newtcolorbox{definition}[1][]{popbase,colframe=popblue,before upper={\popboxhead{Définition}{popblue}{#1}}}
\newtcolorbox{propriete}[1][]{popbase,colframe=poppurple,before upper={\popboxhead{Propriété}{poppurple}{#1}}}
\newtcolorbox{methode}[1][]{popbase,colframe=popgold,before upper={\popboxhead{Méthode}{popgold}{#1}}}
\newtcolorbox{attention}[1][]{popbase,colframe=poppink,before upper={\popboxhead{Attention}{poppink}{#1}}}
\newtcolorbox{experience}[1][]{popbase,colframe=popblue,colback=popblueL,before upper={\popboxhead{Expérience}{popblue}{#1}}}
% « Le savais-tu ? » : carte violette pleine (contexte, culture, Cameroun)
\newtcolorbox{savaistu}[1][]{popbase,colback=poppurple,colframe=popink,
  fontupper=\popsemi\fontsize{10.4}{14.2}\selectfont\color{white},
  before upper={\poppill{Le savais-tu\,?}{white}{poppurple}\ifx\relax#1\relax\else\hspace{6pt}{\popxbold\color{white}#1}\fi\par\vspace{7pt}}}
% Exercice résolu : énoncé en haut, \tcblower puis la solution sur fond crème
\newcounter{popexo}\newcounter{popexr}
\newtcolorbox{exoresolu}[1][]{popbase,colframe=poporange,colbacklower=popcream,
  segmentation style={popink,line width=1.2pt,dashed},
  before upper={\stepcounter{popexr}\popboxhead{Exercice résolu \thepopexr}{poporange}{#1}},
  before lower={\poppill{Solution}{popink}{popcream}\par\vspace{6pt}}}
% Exercice (non résolu en place) — la correction est dans la partie Corrigés
\newtcolorbox{exercice}[1][]{popbase,colframe=popink,
  before upper={\stepcounter{popexo}\popboxhead{Exercice \thepopexo}{popink}{#1}}}
% Corrigé
\newtcolorbox{corrige}[1][]{popbase,colframe=popgreen,colback=popgreenL,
  before upper={\popboxhead{Corrigé}{popgreen}{#1}}}
% Objectifs / prérequis / bilan
\newtcolorbox{objectifs}{popbase,colframe=popink,colback=poporangeL,
  before upper={\popboxhead{Objectifs de la leçon}{popink}{}}}
\newtcolorbox{prerequis}{popbase,colframe=popink,colback=popblueL,
  before upper={\popboxhead{Prérequis}{popblue}{}}}
\newtcolorbox{bilan}{popbase,colframe=popink,colback=white,boxrule=2.8pt,
  before upper={\popboxhead{Fiche bilan}{poporange}{L'essentiel à connaître pour l'examen}}}
% Figure encadrée avec légende
\newtcolorbox{popfigure}[1][]{enhanced,breakable=false,colback=white,colframe=popline,boxrule=1.4pt,arc=8pt,
  left=10pt,right=10pt,top=10pt,bottom=8pt,before skip=10pt,after skip=12pt,halign=center,
  lower separated=false,halign lower=center,
  fontlower=\popsemi\fontsize{9}{11.5}\selectfont\color{popmuted},
  #1}
\newcommand{\legende}[1]{\par\vspace{6pt}{\popsemi\fontsize{9}{11.5}\selectfont\color{popmuted}#1}}

% ============================================================
% Signature OnBuch+
% ============================================================
\newcommand{\poplogoinline}[1]{{\poplogo\fontsize{#1}{#1}\selectfont\color{popink}OnBuch\color{poporange}+}}
\newcommand{\signatureonbuch}{%
  \begin{tcolorbox}[enhanced,colback=popink,colframe=popink,boxrule=2.2pt,arc=10pt,
      drop shadow=poporange,left=14pt,right=14pt,top=10pt,bottom=10pt,before skip=0pt,after skip=0pt]
    \tikz[baseline=-0.7ex]\node[circle,fill=popgreen,draw=popcream,line width=1.3pt,minimum size=22pt,inner sep=0]{\popcheck{white}};%
    \hspace{9pt}\begin{minipage}[c]{0.62\linewidth}
      {\popxbold\fontsize{12}{13}\selectfont\color{popcream}Signé\ \textbullet\ {\poplogo OnBuch\color{poporange}+}}\par\vspace{2pt}
      {\raggedright\popsemi\fontsize{8}{10.5}\selectfont\color{popline}\MakeUppercase{Équipe pédagogique OnBuch+}\\\MakeUppercase{Conforme au programme officiel MINESEC}\par}
    \end{minipage}\hfill
    {\poplogo\fontsize{22}{22}\selectfont\color{popcream}OnBuch\color{poporange}+}
  \end{tcolorbox}%
}

% ============================================================
% Page de garde
% #1 titre · #2 matière · #3 niveau · #4 module · #5 n° leçon · #6 durée ·
% #7 description / objectifs (texte ou itemize)
% ============================================================
\newcommand{\pagegarde}[7]{%
  \thispagestyle{empty}
  \newgeometry{a4paper,margin=1.7cm,top=1.6cm,bottom=1.4cm}%
  \begin{tikzpicture}[remember picture,overlay]
    \fill[poporange!18] ([xshift=-3.2cm,yshift=-3.6cm]current page.north east) circle (4.6cm);
    \fill[poppurple!14] ([xshift=2.4cm,yshift=7.6cm]current page.south west) circle (3.8cm);
  \end{tikzpicture}%
  {\poplogo\fontsize{30}{30}\selectfont\color{popink}OnBuch\color{poporange}+}\hfill
  \raisebox{0.6ex}{\poppill{Cours signé \textbullet\ Premium}{popink}{popcream}}\par
  \vspace{1pt}\popsquiggle{poppurple}\par
  \vspace{8pt}
  \begin{tcolorbox}[enhanced,colback=poporange,colframe=popink,boxrule=2.8pt,arc=13pt,
      drop shadow=popink,left=18pt,right=18pt,top=12pt,bottom=13pt,after skip=10pt]
    \poppill{#2 \textbullet\ #3}{popink}{popcream}\hspace{5pt}\poppill{#4}{white}{popink}\par\vspace{8pt}
    {\popdisplay\fontsize{14}{14}\selectfont\color{popink}LEÇON #5}\par\vspace{4pt}
    {\raggedright\hyphenpenalty=10000\exhyphenpenalty=10000\popdisplay\fontsize{25}{27}\selectfont\color{popcream}\MakeUppercase{#1}\par}
    \vspace{8pt}
    {\popxbold\fontsize{9.5}{9.5}\selectfont\color{popink}\MakeUppercase{Durée conseillée : #6}}
  \end{tcolorbox}
  \begin{tcolorbox}[enhanced,colback=white,colframe=popink,boxrule=2.2pt,arc=10pt,
      drop shadow=popink,left=16pt,right=16pt,top=10pt,bottom=10pt,after skip=8pt]
    \poppill{Dans cette leçon}{poppurple}{white}\par\vspace{5pt}
    \popsemi\fontsize{9.3}{12.6}\selectfont\color{popink}#7
  \end{tcolorbox}
  \noindent\begin{minipage}[t]{0.487\linewidth}
    \begin{tcolorbox}[enhanced,colback=popgreen,colframe=popink,boxrule=2.2pt,arc=10pt,
        drop shadow=popink,left=11pt,right=11pt,top=10pt,bottom=10pt,height=3.1cm,valign=center]
      \tcbox[colback=white,colframe=popink,boxrule=1.3pt,arc=5pt,boxsep=0pt,nobeforeafter,
        left=3pt,right=3pt,top=3pt,bottom=3pt,valign=center]{\qrcode[height=1.9cm]{https://play.google.com/store/apps/details?id=com.luvvix.onbuch&hl=fr}}%
      \hspace{8pt}\begin{minipage}[c]{0.5\linewidth}
        {\popdisplay\fontsize{11}{12.5}\selectfont\color{white}\MakeUppercase{Télécharge OnBuch}}\par\vspace{4pt}
        {\raggedright\popsemi\fontsize{8.4}{11}\selectfont\color{white}Gratuit \textbullet\ Google Play\\Tuteur IA, annales, concours, résultats\par}
      \end{minipage}
    \end{tcolorbox}
  \end{minipage}\hfill
  \begin{minipage}[t]{0.487\linewidth}
    \begin{tcolorbox}[enhanced,colback=poppurple,colframe=popink,boxrule=2.2pt,arc=10pt,
        drop shadow=popink,left=11pt,right=11pt,top=10pt,bottom=10pt,height=3.1cm,valign=center]
      \tikz[baseline=-0.7ex]\node[circle,fill=white,draw=popink,line width=1.3pt,minimum size=1.6cm,inner sep=0]{\popdisplay\fontsize{24}{24}\selectfont\color{poppurple}+};%
      \hspace{8pt}\begin{minipage}[c]{0.6\linewidth}
        {\popdisplay\fontsize{11}{12.5}\selectfont\color{white}\MakeUppercase{Passe à OnBuch+}}\par\vspace{4pt}
        {\raggedright\popsemi\fontsize{8.4}{11}\selectfont\color{white}Tous les cours signés, Tuteur IA illimité, fiches PDF — onglet OnBuch+ de l'app\par}
      \end{minipage}
    \end{tcolorbox}
  \end{minipage}
  \vspace{8pt}
  \popcontenu
  \vfill
  \signatureonbuch
  \restoregeometry
  \newpage
}

% Rangée « Ce cours contient » (page de garde)
\newcommand{\poptile}[3]{% #1 couleur #2 gros texte #3 légende
  \begin{tcolorbox}[enhanced,colback=#1,colframe=popink,boxrule=2pt,arc=9pt,shadow={2.5pt}{-2.5pt}{0pt}{popink},
      left=6pt,right=6pt,top=9pt,bottom=9pt,halign=center,valign=center,height=1.75cm,width=\linewidth,nobeforeafter]
    {\popdisplay\fontsize{12.5}{13}\selectfont\color{white}\MakeUppercase{#2}}\par\vspace{3pt}
    {\centering\popsemi\fontsize{7.8}{9.5}\selectfont\color{white}#3\par}
  \end{tcolorbox}}
\newcommand{\popcontenu}{%
  \par\noindent{\popxboldls\fontsize{7.5}{7.5}\selectfont\color{popmuted}CE COURS SIGNÉ CONTIENT}\par\vspace{7pt}
  \noindent\begin{minipage}[t]{0.235\linewidth}\poptile{popblue}{Cours}{complet, illustré, pas à pas}\end{minipage}\hfill
  \begin{minipage}[t]{0.235\linewidth}\poptile{poporange}{Méthodes}{et exercices résolus}\end{minipage}\hfill
  \begin{minipage}[t]{0.235\linewidth}\poptile{poppink}{Exercices}{type examen + corrigés}\end{minipage}\hfill
  \begin{minipage}[t]{0.235\linewidth}\poptile{popgreen}{Fiche bilan}{l'essentiel à retenir}\end{minipage}\par}

% Sommaire
\newtcolorbox{sommaire}{enhanced,colback=white,colframe=popink,boxrule=2.2pt,arc=10pt,
  drop shadow=popink,left=18pt,right=18pt,top=13pt,bottom=13pt,before skip=6pt,after skip=20pt,
  before upper={\poppill{Sommaire}{poppurple}{white}\par\vspace{9pt}\popsemi\fontsize{10.6}{14.5}\selectfont\color{popink}}}

% Signature de fin de cours — poussée tout en bas de la dernière page
\newcommand{\findecours}{%
  \par\vfill\nopagebreak
  \begin{center}\popsquiggle{poporange}\end{center}\vspace{4pt}
  \signatureonbuch
}
\AtBeginDocument{\def\llbracket{\mathopen{[\mkern-3.5mu[}}\def\rrbracket{\mathclose{]\mkern-3.5mu]}}} % intervalles d'entiers (glyphe ⟦ absent de la police)
% Glyphes absents de Fira Math : redéfinis à partir de glyphes disponibles
\AtBeginDocument{%
  \def\setminus{\mathbin{\backslash}}\let\smallsetminus\setminus
  \def\vdots{\mathord{\vbox{\baselineskip3.2pt\lineskiplimit0pt\kern4pt\hbox{.}\hbox{.}\hbox{.}}}}%
  \def\ddots{\mathinner{\mkern1mu\raise6.5pt\hbox{.}\mkern2mu\raise3.6pt\hbox{.}\mkern2mu\raise0.7pt\hbox{.}\mkern1mu}}%
  \def\triangle{\mathord{\scalebox{0.9}{$\symup{\Delta}$}}}%
  \def\nabla{\mathord{\raisebox{\depth}{\scalebox{1}[-1]{$\symup{\Delta}$}}}}%
  \def\checkmark{\popcheck{popdark}}%
  \def\star{\mathbin{\ast}}\def\bigstar{\mathord{\ast}}%
  \def\diamond{\mathbin{\scalebox{0.6}{$\Diamond$}}}%
  \def\bigcup{\mathop{\mathchoice{\raisebox{-0.3em}{\scalebox{1.7}{$\cup$}}}{\scalebox{1.25}{$\cup$}}{\cup}{\cup}}\displaylimits}%
  \def\bigcap{\mathop{\mathchoice{\raisebox{-0.3em}{\scalebox{1.7}{$\cap$}}}{\scalebox{1.25}{$\cap$}}{\cap}{\cap}}\displaylimits}%
  \def\lesseqqgtr{\mathrel{\lessgtr}}\def\bigoplus{\mathop{\oplus}\displaylimits}%
}
\providecommand{\ssi}{si et seulement si} % abréviation fréquente
\AtBeginDocument{\providecommand{\wideparen}{\overparen}} % arc de cercle

%% FIGFIX-BEGIN v3 — correctifs globaux des figures (docs/onbuch-generation/tools/figfix)
\usepackage{adjustbox}\usepackage{etoolbox}
% toute figure TikZ/pgfplots plus large que la ligne est réduite à la largeur disponible
\BeforeBeginEnvironment{tikzpicture}{\begin{adjustbox}{max width=\linewidth}}
\AfterEndEnvironment{tikzpicture}{\end{adjustbox}}
% légendes pgfplots lisibles par-dessus les courbes
\pgfplotsset{every axis/.append style={legend style={fill=white,fill opacity=0.92,text opacity=1,draw=black!15,inner sep=3pt}}}
% étiquettes d'arêtes (syntaxe edge["..."]) avec fond blanc
\tikzset{every edge quotes/.append style={fill=white,inner sep=1.2pt}}
% bulles de cartes mentales : texte à la ligne dans la bulle, bulle un peu plus grande
\tikzset{every concept/.append style={text width=2.0cm,align=center,minimum size=2.3cm,inner sep=2pt}}
%% FIGFIX-END
\begin{document}

\pagegarde{Les échanges du Cameroun avec l’extérieur}{Géographie}{Tle A}{Module 2 : La libéralisation des échanges (la mondialisation)}{N16}{2 h + dossier 5 (2 h)}{Cette leçon étudie les échanges extérieurs du Cameroun : ses partenaires africains (notamment dans la CEMAC) et mondiaux, les produits exportés et importés, ainsi que les problèmes structurels (déficit commercial, détérioration des termes de l'échange). Elle s'achève par l'étude guidée des marchés frontaliers, lieux concrets d'intégration régionale.
\vspace{4pt}
\begin{multicols}{2}\small
\begin{itemize}
\item À la fin de cette leçon, je sais définir balance commerciale, termes de l'échange et CEMAC
\item À la fin de cette leçon, je sais localiser sur une carte les principaux partenaires commerciaux du Cameroun en Afrique et dans le monde
\item À la fin de cette leçon, je sais inventorier et classer les principaux produits exportés et importés par le Cameroun
\item À la fin de cette leçon, je sais interpréter un tableau statistique ou un graphique de la balance commerciale camerounaise
\item À la fin de cette leçon, je sais construire un diagramme (histogrammes, secteurs) représentant la structure des échanges du Cameroun
\item À la fin de cette leçon, je sais légender un croquis ou une carte des flux commerciaux du Cameroun
\end{itemize}\end{multicols}}

\begin{sommaire}
\begin{multicols}{2}
\begin{enumerate}
\item Cadre général et notions clés des échanges extérieurs
\item Les produits échangés
\item Les échanges avec les partenaires africains et la CEMAC
\item Les échanges avec le reste du monde
\item Les problèmes des échanges extérieurs du Cameroun
\item Dossier 5 : Les marchés frontaliers du Cameroun
\item Activité d'intégration
\item Méthodes et exercices résolus
\item Exercices
\item Corrigés des exercices
\item Fiche bilan
\end{enumerate}
\end{multicols}
\end{sommaire}

\begin{objectifs}
\begin{itemize}
\item À la fin de cette leçon, je sais définir balance commerciale, termes de l'échange et CEMAC
\item À la fin de cette leçon, je sais localiser sur une carte les principaux partenaires commerciaux du Cameroun en Afrique et dans le monde
\item À la fin de cette leçon, je sais inventorier et classer les principaux produits exportés et importés par le Cameroun
\item À la fin de cette leçon, je sais interpréter un tableau statistique ou un graphique de la balance commerciale camerounaise
\item À la fin de cette leçon, je sais construire un diagramme (histogrammes, secteurs) représentant la structure des échanges du Cameroun
\item À la fin de cette leçon, je sais légender un croquis ou une carte des flux commerciaux du Cameroun
\item À la fin de cette leçon, je sais expliquer la détérioration des termes de l'échange et ses conséquences pour le Cameroun
\item À la fin de cette leçon, je sais décrire le fonctionnement et l'impact des marchés frontaliers (Banki, Kyé-Ossi, Abam-Minko, Ekondo-Titi…)
\end{itemize}
\end{objectifs}

\begin{prerequis}
\begin{itemize}
\item Notion de mondialisation et de libéralisation des échanges (Module 2)
\item Les grandes ressources du Cameroun (agriculture, pétrole, forêt) vues en classe de Première
\item Localisation des pays voisins du Cameroun et des organisations régionales africaines
\item Lecture de tableaux statistiques et construction de diagrammes simples
\end{itemize}
\end{prerequis}

\newpage

\coursec{Cadre général et notions clés des échanges extérieurs}

\courssub{Une situation d'accroche : deux visages du commerce camerounais}

Imagine-toi un matin de marché à \textbf{Banki}, dans l'Extrême-Nord, tout près de la frontière nigériane. Des femmes camerounaises étalent du maïs, du sel et des pagnes colorés ; des acheteurs nigérians chargent des sacs sur des motos. Dans l'autre sens, des motos neuves, des bidons d'huile de moteur et des téléphones portables entrent au Cameroun. C'est le \emph{commerce frontalier}, vif, informel pour une bonne part, qui fait vivre des milliers de familles.

Change de décor : le \textbf{port autonome de Douala}, premier port d'Afrique centrale. Ici, ce sont des grues géantes qui déchargent des conteneurs de riz, de voitures, de médicaments et de machines-outils. De l'autre côté du quai, des navires embarquent du cacao, des grumes et des sciages, du coton, de la banane ; et au large, les champs pétroliers alimentent des tankers qui partent vers la Chine, l'Europe ou l'Inde. C'est le \emph{commerce maritime international}, formel, encadré par des contrats, des assurances, des lettres de crédit.

Pourtant, derrière ces deux images, une même réalité chiffrée : \textbf{la valeur de ce qui entre (importations) dépasse presque chaque année la valeur de ce qui sort (exportations)}. Le Cameroun achète plus qu'il ne vend. Comment expliquer ce paradoxe ? Quels produits, quels partenaires, quels mécanismes ? C'est toute la problématique de cette leçon : \textbf{Comment le Cameroun s'insère-t-il dans la mondialisation par ses échanges extérieurs, et pourquoi sa balance commerciale reste-t-elle structurellement déficitaire ?}

\courssub{Définition et vocabulaire des échanges extérieurs}

Les \cle{échanges extérieurs} désignent l'ensemble des flux de marchandises (biens matériels) qui franchissent les frontières d'un pays. On distingue trois grands flux :

\begin{definition}[Exportation, importation, réexportation]
\begin{itemize}
    \item \textbf{Exportation} : sortie définitive d'un bien du territoire national vers l'étranger (vente à l'extérieur). Exemple : cacao camerounais expédié vers les Pays-Bas.
    \item \textbf{Importation} : entrée définitive d'un bien étranger sur le territoire national (achat à l'extérieur). Exemple : voitures d'occasion importées d'Europe ou d'Asie via le port de Douala.
    \item \textbf{Réexportation} : exportation de biens importés antérieurement sans transformation notable (ou après simple conditionnement). Le Cameroun réexporte une partie des produits pétroliers raffinés à la SONARA vers le Tchad ou la Centrafrique, ainsi que des marchandises de transit destinées aux pays enclavés.
\end{itemize}
\end{definition}

\begin{exemplebox}[Le circuit d'un sac de cacao et d'une voiture]
\emph{Exportation (cacao) :} Planteur de Ngomedzap (Centre) $\rightarrow$ piste rurale $\rightarrow$ collecteur/coopérative $\rightarrow$ camion vers Douala $\rightarrow$ entrepôt portuaire $\rightarrow$ navire $\rightarrow$ Rotterdam (Pays-Bas) $\rightarrow$ transformateur (Barry Callebaut, Cargill) $\rightarrow$ chocolat européen.

\emph{Importation (voiture d'occasion) :} Port européen ou asiatique $\rightarrow$ navire roulier (Ro-Ro) $\rightarrow$ Douala (dédouanement, TVA, droits de douane) $\rightarrow$ parc automobile (Bassa, Akwa) $\rightarrow$ acheteur à Yaoundé ou Bamenda.
\end{exemplebox}

\begin{popfigure}
\begin{tikzpicture}[scale=0.9, every node/.style={font=\small}]
    \tikzset{
        lieu/.style={rectangle, draw=popdark, fill=popblueL, rounded corners, minimum width=2.2cm, minimum height=1cm, align=center},
        port/.style={rectangle, draw=popdark, fill=poporangeL, rounded corners, minimum width=2.5cm, minimum height=1.2cm, align=center},
        mer/.style={rectangle, draw=popdark, fill=popgreenL, rounded corners, minimum width=2.5cm, minimum height=1cm, align=center},
        dest/.style={rectangle, draw=popdark, fill=poppinkL, rounded corners, minimum width=2.2cm, minimum height=1cm, align=center},
        fleche/.style={->, thick, popdark},
        flechemer/.style={->, very thick, popblue, decorate, decoration={snake, amplitude=1mm, segment length=4mm}}
    }
    \node[lieu, align=center] (planteur) at (0,4){Planteur cacao\\(Ngomedzap)};
    \node[lieu, align=center] (collecteur) at (4,4){Collecteur /\\Coopérative};
    \node[port, align=center] (doualaexp) at (8,4){Port de Douala\\(Embarquement)};
    \node[mer, align=center] (mer) at (12,4){Transport maritime\\(Atlantique)};
    \node[dest, align=center] (europe) at (16,4){Europe / Chine\\(Transformation)};
    \node[dest, align=center] (japon) at (16,0){Europe / Asie\\(Usine / Occasion)};
    \node[mer, align=center] (merimp) at (12,0){Transport maritime\\(Ro-Ro / Conteneur)};
    \node[port, align=center] (doualaimp) at (8,0){Port de Douala\\(Dédouanement)};
    \node[lieu, align=center] (parc) at (4,0){Parc auto\\(Douala / Yaoundé)};
    \node[lieu, align=center] (acheteur) at (0,0){Acheteur final\\(Yaoundé, Bafoussam)};
    \draw[fleche] (planteur) -- node[above, sloped] {Piste / Camion} (collecteur);
    \draw[fleche] (collecteur) -- node[above] {Camion} (doualaexp);
    \draw[flechemer] (doualaexp) -- node[above] {Navire} (mer);
    \draw[fleche] (mer) -- node[above] {Déchargement} (europe);
    \draw[fleche] (japon) -- node[above] {Navire} (merimp);
    \draw[flechemer] (merimp) -- node[above] {Arrivée} (doualaimp);
    \draw[fleche] (doualaimp) -- node[above] {Camion / Route} (parc);
    \draw[fleche] (parc) -- node[above] {Route} (acheteur);
    \node[align=center, font=\bfseries\small, popblue] at (8,5.4) {FLUX D'EXPORTATION (cacao, bois, pétrole, coton, banane...)};
    \node[align=center, font=\bfseries\small, poporange] at (8,-1.4) {FLUX D'IMPORTATION (produits manufacturés, carburants, blé, riz, machines...)};
\end{tikzpicture}
\legende{Schéma simplifié des circuits d'une exportation type (cacao) et d'une importation type (véhicule) via le port de Douala. Les flèches bleues ondulées symbolisent le transport maritime.}
\end{popfigure}

\courssub{Le Cameroun dans la mondialisation : ouverture et accords}

Le Cameroun est une \emph{économie ouverte} : son taux d'ouverture ((exportations + importations) / PIB) se situe autour de \num{40} à \num{50}\,\% selon les années (ordre de grandeur). Cette insertion repose sur des cadres juridiques multilatéraux et bilatéraux.

\begin{aretenir}[Cadres institutionnels des échanges du Cameroun]
\begin{itemize}
    \item \textbf{OMC (Organisation mondiale du commerce) :} le Cameroun est membre depuis 1995 (il était déjà membre du GATT). Il s'engage sur la non-discrimination (clause de la nation la plus favorisée), la transparence des barrières tarifaires et le règlement des différends.
    \item \textbf{APE (Accord de partenariat économique) avec l'Union européenne :} signé en 2009 (accord d'étape), ratifié et entré en application en 2014. Principe : libéralisation \textbf{asymétrique et progressive} — l'UE ouvre immédiatement son marché aux produits camerounais, tandis que le Cameroun démantèle ses droits de douane sur environ \num{80}\,\% de ses importations en provenance de l'UE sur une quinzaine d'années. Des mesures d'accompagnement et des clauses de sauvegarde protègent les productions naissantes.
    \item \textbf{CEMAC et ZLECAf :} la Communauté économique et monétaire de l'Afrique centrale (étudiée en section 3) organise le marché régional ; la Zone de libre-échange continentale africaine (ZLECAf, opérationnelle depuis 2021) vise à supprimer progressivement les droits de douane sur environ \num{90}\,\% des lignes tarifaires intra-africaines.
\end{itemize}
\end{aretenir}

\begin{attention}[Ne pas confondre « libéralisation » et « libre-échange total »]
Au Bac, on attend de vous que vous précisiez que la libéralisation est \textbf{progressive, sélective et asymétrique} (délais plus longs pour les pays en développement comme le Cameroun). Dire « le Cameroun a supprimé toutes ses barrières douanières » est \textbf{faux} et sanctionné. Citez l'APE UE-Cameroun et la ZLECAf comme cadres actuels.
\end{attention}

\courssub{La balance commerciale : définition, calcul et situation du Cameroun}

C'est l'indicateur synthétique des échanges de biens.

\begin{definition}[Balance commerciale (BC)]
La \cle{balance commerciale} est la différence entre la valeur des exportations (X) et la valeur des importations (M) de biens sur une période donnée (généralement l'année) :
\[
BC = X - M
\]
\begin{itemize}
    \item $BC > 0$ : \textbf{excédent} (le pays vend plus qu'il n'achète) ;
    \item $BC < 0$ : \textbf{déficit} (le pays achète plus qu'il ne vend) ;
    \item $BC = 0$ : \textbf{équilibre}.
\end{itemize}
\emph{Attention :} la balance commerciale ne porte que sur les \textbf{biens} (marchandises). La \cle{balance des paiements} inclut aussi les services (transport, tourisme, transferts) et les mouvements de capitaux.
\end{definition}

\begin{methode}[Calculer et interpréter un solde commercial à partir d'un tableau]
\begin{enumerate}
    \item \textbf{Identifier} les colonnes « Exportations (FOB) » et « Importations (CIF) ». \emph{Rappel :} FOB (\emph{free on board}) = valeur à la sortie du port exportateur ; CIF (\emph{cost, insurance, freight}) = valeur à l'entrée du port importateur (inclut fret et assurance). Les importations CIF sont donc mécaniquement plus élevées que si elles étaient évaluées FOB.
    \item \textbf{Additionner} les totaux annuels pour X et M (souvent en milliards de FCFA).
    \item \textbf{Calculer} $BC = X - M$.
    \item \textbf{Interpréter le signe :} déficit structurel ? excédent conjoncturel (pétrole) ? évolution sur 5-10 ans (amélioration ou aggravation) ?
    \item \textbf{Calculer le taux de couverture :} $\frac{X}{M} \times 100$. S'il est inférieur à 100, les exportations ne couvrent pas les importations.
\end{enumerate}
\end{methode}

\begin{exoresolu}[Calcul du solde commercial camerounais (ordres de grandeur)]
\textbf{Énoncé :} D'après les données de l'INS et de la BEAC, on retient les ordres de grandeur suivants pour les échanges de biens (en milliards de FCFA) :
\begin{center}
\begin{tabular}{|l|c|c|c|}\hline
\textbf{Année} & \textbf{Exportations (X)} & \textbf{Importations (M)} & \textbf{Balance (X-M)} \\ \hline
2020 & 2\,800 & 3\,500 & ? \\ \hline
2021 & 3\,200 & 3\,800 & ? \\ \hline
2022 & 4\,100 & 4\,500 & ? \\ \hline
2023 (est.) & 3\,900 & 4\,800 & ? \\ \hline
\end{tabular}
\end{center}
\textbf{1.} Complétez le tableau. \textbf{2.} Calculez le taux de couverture pour 2022. \textbf{3.} Interprétez.

\tcblower
\textbf{Solution détaillée :}

\textbf{1.} Tableau complété :
\begin{center}
\begin{tabular}{|l|c|c|c|}\hline
\textbf{Année} & \textbf{Exportations (X)} & \textbf{Importations (M)} & \textbf{Balance (X-M)} \\ \hline
2020 & 2\,800 & 3\,500 & \textbf{-700} \\ \hline
2021 & 3\,200 & 3\,800 & \textbf{-600} \\ \hline
2022 & 4\,100 & 4\,500 & \textbf{-400} \\ \hline
2023 (est.) & 3\,900 & 4\,800 & \textbf{-900} \\ \hline
\end{tabular}
\end{center}

\textbf{2.} Taux de couverture 2022 :
\[
\frac{X}{M} \times 100 = \frac{4\,100}{4\,500} \times 100 \approx 91{,}1\,\%
\]

\textbf{3.} \emph{Interprétation :} la balance est \textbf{structurellement déficitaire} sur toute la période. Le déficit se réduit entre 2020 et 2022 (reprise post-Covid, cours pétroliers et gaziers favorables), puis se creuse fortement en 2023 (hausse de la facture alimentaire et énergétique, recul de certains cours). Le taux de couverture inférieur à 100\,\% confirme que le pays doit financer l'écart par d'autres ressources : revenus de services, transferts de la diaspora, investissements directs étrangers et endettement extérieur.
\end{exoresolu}

\begin{popfigure}
\begin{tikzpicture}
\begin{axis}[
    width=\linewidth,
    height=9cm,
    xlabel={Année},
    ylabel={Valeur (milliards de FCFA)},
    legend pos=north west,
    grid=major,
    grid style={popmuted},
    tick label style={font=\small},
    label style={font=\small},
    xtick={2018,2019,2020,2021,2022,2023},
    ymin=2000, ymax=5500,
]
\addplot[popcurve, popblue, mark=*, mark size=2] coordinates {
    (2018, 2900) (2019, 3000) (2020, 2800) (2021, 3200) (2022, 4100) (2023, 3900)
};
\addlegendentry{Exportations (X)}
\addplot[popcurve, poporange, mark=*, mark size=2] coordinates {
    (2018, 3400) (2019, 3600) (2020, 3500) (2021, 3800) (2022, 4500) (2023, 4800)
};
\addlegendentry{Importations (M)}
\end{axis}
\end{tikzpicture}
\legende{Évolution des exportations et importations du Cameroun (milliards de FCFA, ordres de grandeur INS/BEAC). L'écart entre les deux courbes représente le déficit commercial. Le pic d'exportations de 2022 est lié aux cours élevés du pétrole et du gaz.}
\end{popfigure}

\courssub{Les termes de l'échange : un piège structurel pour les exportateurs de matières premières}

Au-delà des volumes, ce sont les \textbf{prix} qui pèsent lourd.

\begin{definition}[Termes de l'échange (TE)]
Les \cle{termes de l'échange} mesurent le pouvoir d'achat des exportations en biens importés. C'est le rapport entre l'indice des prix des exportations ($P_X$) et l'indice des prix des importations ($P_M$), base 100 sur une année de référence :
\[
TE = \frac{P_X}{P_M} \times 100
\]
\begin{itemize}
    \item $TE > 100$ : \textbf{amélioration} (une unité exportée achète plus d'importations qu'à l'année de base) ;
    \item $TE < 100$ : \textbf{détérioration} (il faut exporter davantage pour acheter la même quantité d'importations).
\end{itemize}
\end{definition}

\begin{propriete}[Détérioration tendancielle des termes de l'échange (thèse de Prebisch-Singer)]
Pour les pays exportateurs de \textbf{produits primaires} (pétrole, cacao, bois, coton, aluminium) et importateurs de \textbf{produits manufacturés} (machines, véhicules, médicaments, engrais), les termes de l'échange ont tendance à se détériorer sur le long terme.
\textbf{Raisons :}
\begin{enumerate}
    \item \textbf{Élasticité-revenu de la demande :} la demande de produits manufacturés croît plus vite que le revenu, alors que la demande de produits primaires (alimentaires, matières premières) croît moins vite.
    \item \textbf{Progrès technique asymétrique :} dans l'industrie (Nord), les gains de productivité se traduisent en hausses de revenus et prix soutenus ; dans l'agriculture et l'extraction (Sud), ils se traduisent plutôt en baisse des prix, sur des marchés très concurrentiels.
    \item \textbf{Volatilité :} les cours des matières premières (pétrole, cacao) flambent puis s'effondrent (chocs d'offre, spéculation), alors que les prix des produits manufacturés sont rigides à la baisse.
\end{enumerate}
\emph{Conséquence pour le Cameroun :} même si le volume des exportations augmente, leur valeur peut stagner ou baisser quand les cours s'effondrent (le prix du cacao au producteur a par exemple connu de fortes variations, de l'ordre de \SI{1000}{FCFA/kg} à plus de \SI{2000}{FCFA/kg} selon les campagnes), alors que la facture d'importation (riz, blé, machines, carburant) ne baisse pas symétriquement.
\end{propriete}

\begin{attention}[Piège rédactionnel : « termes de l'échange » $\neq$ « balance commerciale »]
\begin{itemize}
    \item \textbf{Balance commerciale} = différence de \textbf{valeurs} (volumes $\times$ prix) : $X - M$. Elle peut s'améliorer si les volumes exportés augmentent massivement (nouvelle mine, nouveau champ pétrolier).
    \item \textbf{Termes de l'échange} = rapport de \textbf{prix} uniquement ($P_X / P_M$). Ils peuvent se détériorer pendant que la balance s'améliore, si les volumes compensent la baisse des prix.
    \item \textbf{Au Bac :} si le sujet porte sur les \emph{termes de l'échange}, parlez de \textbf{prix relatifs}, de \textbf{pouvoir d'achat des exportations}, citez \textbf{Prebisch-Singer}. Ne vous contentez pas de dire « le déficit s'aggrave ».
\end{itemize}
\end{attention}

\courssub{Le port de Douala et le corridor Douala--N'Djamena : poumon et artère des échanges}

Le Cameroun dispose de deux ports principaux : le \textbf{Port autonome de Douala (PAD)}, port polyvalent historique, et le \textbf{Port autonome de Kribi (PAK)}, port en eau profonde entré en service en 2014 (conteneurs, vrac). Le PAD reste toutefois la porte d'entrée et de sortie numéro 1 du pays.

\begin{savaistu}[Le port de Douala : hub sous-régional]
\begin{itemize}
    \item Il traite \textbf{l'essentiel du commerce extérieur camerounais} (hors pétrole brut chargé en mer) et \textbf{la quasi-totalité du commerce extérieur du Tchad} (pays enclavé) ainsi qu'une part importante de celui de la Centrafrique.
    \item Trafic total : de l'ordre de \textbf{12 à 14 millions de tonnes par an} (ordre de grandeur), dominé par les importations, avec une part significative de transit vers le Tchad et la RCA.
    \item Infrastructures : quais conventionnels, terminal à conteneurs, quais à bois, à céréales et à hydrocarbures, silos.
    \item \textbf{Problèmes :} encombrement (délais de séjour des marchandises souvent supérieurs aux standards internationaux), coûts de manutention et de passage élevés, vétusté de certains équipements, multiplicité des contrôles (douane, police, gendarmerie) qui allongent les délais et renchérissent les coûts.
\end{itemize}
\end{savaistu}

\begin{experience}[Analyse du corridor Douala--N'Djamena (TP : construction d'un croquis)]
\textbf{Document support :} carte routière du Cameroun ; axe Douala $\rightarrow$ Yaoundé $\rightarrow$ Ngaoundéré $\rightarrow$ Maroua $\rightarrow$ Kousséri (frontière du Tchad) ; embranchement Bertoua $\rightarrow$ Garoua-Boulaï (frontière de la RCA).

\textbf{Consigne :} réalisez un croquis légendé du corridor Douala--N'Djamena en montrant :
\begin{enumerate}
    \item le port de Douala (porte d'entrée) ;
    \item les deux axes principaux : Douala--Yaoundé--Ngaoundéré--Maroua--Kousséri vers N'Djamena ; Yaoundé--Bertoua--Garoua-Boulaï vers Bangui ;
    \item les principaux goulots d'étranglement : traversée de Yaoundé (embouteillages), tronçons dégradés (notamment Bertoua--Garoua-Boulaï), postes frontaliers (Kousséri, Garoua-Boulaï) avec leurs files d'attente ;
    \item les flux : importations (conteneurs, carburant, blé, riz, matériaux) vers le Nord, le Tchad et la RCA ; exportations (coton, bétail, sésame, arachide, gomme arabique) du Nord, du Tchad et de la RCA vers Douala.
\end{enumerate}

\textbf{Interprétation attendue :} le corridor est \textbf{long (environ \num{1800} km entre Douala et N'Djamena), routier, saturé et coûteux}. Le temps de trajet réel atteint souvent \num{7} à \num{10} jours, et le coût du transport peut représenter une part très élevée (jusqu'à \num{30}-\num{40}\,\%) de la valeur des marchandises livrées au Tchad. \textbf{Perspectives :} meilleure exploitation du chemin de fer (Camrail, axe Douala--Yaoundé--Ngaoundéré), bitumage des tronçons manquants, désenclavement par le port de Kribi.
\end{experience}

\begin{aretenir}[Synthèse du cadre général]
\begin{itemize}
    \item Les échanges extérieurs du Cameroun reposent sur l'exportation de \textbf{produits primaires peu transformés} (pétrole, cacao, bois, coton, aluminium, banane) et l'importation de \textbf{produits manufacturés et de biens de consommation} (machines, véhicules, carburants raffinés, blé, riz, médicaments).
    \item La \textbf{balance commerciale est structurellement déficitaire} (taux de couverture souvent inférieur à 100\,\%), seulement atténuée les années de cours pétroliers élevés.
    \item La \textbf{détérioration des termes de l'échange} (baisse relative des prix des matières premières face aux prix des produits manufacturés) érode le pouvoir d'achat des exportations.
    \item Le \textbf{port de Douala} et le \textbf{corridor vers le Nord, le Tchad et la RCA} sont des infrastructures stratégiques mais congestionnées, qui freinent la compétitivité.
    \item Les \textbf{accords internationaux} (OMC, APE avec l'UE, CEMAC, ZLECAf) imposent une libéralisation progressive qui expose les filières locales à la concurrence.
\end{itemize}
\end{aretenir}

\coursec{Les produits échangés}

Dans la section précédente, nous avons posé le cadre général des échanges extérieurs du Cameroun : la balance commerciale, les termes de l'échange, les grands partenaires. Nous savons donc \emph{combien} le Cameroun échange et \emph{avec qui}. Il reste une question décisive : \textbf{qu'échange-t-il exactement ?} Car en géographie comme en économie, la nature des produits échangés dit presque tout de la place d'un pays dans la mondialisation. Pour le comprendre, partons d'une situation concrète.

\begin{exemplebox}[Deux camions à Douala]
Au port autonome de Douala, un camion quitte le terminal à bois chargé de grumes d'ayous destinées à la Chine ; un autre camion, venant d'un navire qui vient d'accoster, décharge des sacs de riz thaïlandais et des cartons de médicaments indiens. Le premier transporte une matière \emph{brute}, le second des produits \emph{transformés}. Cette scène quotidienne résume à elle seule la structure des échanges du Cameroun : le pays vend ce que la nature lui donne et achète ce que l'industrie des autres fabrique.
\end{exemplebox}

\courssub{Les produits d'exportation : le règne des matières premières}

\textbf{Observation et inventaire.} Interrogeons les statistiques douanières et les rapports de la Banque Mondiale et de l'INS (Institut National de la Statistique). Les principaux produits exportés par le Cameroun sont, par ordre d'importance approximatif :

\begin{itemize}
\item le \cle{pétrole brut}, extrait surtout au large de Rio del Rey (Sud-Ouest) et dans le bassin de Logone Birni (Extrême-Nord), longtemps premier produit d'exportation ;
\item le \cle{cacao}, dont le Cameroun est l'un des tout premiers producteurs mondiaux (quatrième ou cinquième rang mondial selon les années), cultivé dans le Centre, le Sud, le Sud-Ouest et l'Est ;
\item le \cle{bois} (grumes et sciages : ayous, sapelli, iroko), issu des forêts du Sud, de l'Est et du Littoral ;
\item le \cle{café} (robusta dans l'Ouest et le Littoral, arabica dans les hautes terres de l'Ouest et de l'Adamaoua) ;
\item le \cle{coton} (« or blanc » du Nord, produit sous l'égide de la SODECOTON et filé en partie par la CICAM) ;
\item la \cle{banane} (plantations de la PHP — Plantations du Haut Penja — et de la CDC dans le Littoral et le Sud-Ouest) ;
\item l'\cle{aluminium}, produit à Edéa par ALUCAM à partir d'alumine importée ;
\item le \cle{caoutchouc} (hévéas du Sud et du Littoral : HEVECAM, SUDCAM) ;
\item depuis 2018, le \cle{gaz naturel liquéfié} (GNL), exporté depuis l'unité flottante de Kribi (Hilli Episeyo) : une nouveauté majeure.
\end{itemize}

\begin{definition}[Produit d'exportation du Cameroun]
Les exportations camerounaises sont dominées par des \cle{produits primaires}, c'est-à-dire des matières premières agricoles, forestières ou minières, vendues \textbf{brutes ou très peu transformées}. Leur prix est fixé sur les marchés mondiaux (bourses de Londres, New York, Rotterdam), sans que le Cameroun puisse l'influencer : le pays est \emph{preneur de prix}.
\end{definition}

\begin{exemplebox}[Ordres de grandeur des exportations]
Selon les années (les cours mondiaux font fortement varier les parts), la structure des exportations camerounaises en valeur se présente approximativement ainsi : le pétrole brut (et désormais le gaz) représente environ \textbf{35 à 40 \%} des recettes d'exportation ; le cacao et ses produits environ \textbf{15 à 20 \%} ; le bois environ \textbf{10 à 15 \%} ; le coton, la banane, l'aluminium, le café et le caoutchouc se partagent le reste, à raison de 3 à 7 \% chacun. Ces chiffres sont des ordres de grandeur : une chute du prix du baril, comme en 2014-2015 ou en 2020, fait mécaniquement fondre la part du pétrole.
\end{exemplebox}

\begin{popfigure}
\begin{center}
\begin{tikzpicture}[scale=1.05]
% Secteurs : pétrole 38% (136.8°), cacao 18% (64.8°), bois 12% (43.2°), coton 7% (25.2°), banane 6% (21.6°), aluminium 5% (18°), autres 14% (50.4°)
\fill[popblue] (0,0) -- (90:2.6) arc (90:226.8:2.6) -- cycle;
\fill[poporange] (0,0) -- (226.8:2.6) arc (226.8:291.6:2.6) -- cycle;
\fill[popgreen] (0,0) -- (291.6:2.6) arc (291.6:334.8:2.6) -- cycle;
\fill[poppurple] (0,0) -- (334.8:2.6) arc (334.8:360:2.6) -- cycle;
\fill[poppink] (0,0) -- (0:2.6) arc (0:21.6:2.6) -- cycle;
\fill[popgold] (0,0) -- (21.6:2.6) arc (21.6:39.6:2.6) -- cycle;
\fill[popmuted] (0,0) -- (39.6:2.6) arc (39.6:90:2.6) -- cycle;
\draw[popdark] (0,0) circle (2.6);
% Légende
\begin{scope}[xshift=4.2cm, yshift=1.8cm]
\fill[popblue] (0,0) rectangle (0.35,0.25); \node[right, font=\small] at (0.45,0.125) {1. Pétrole brut et gaz : 38 \%};
\fill[poporange] (0,-0.55) rectangle (0.35,-0.30); \node[right, font=\small] at (0.45,-0.425) {2. Cacao : 18 \%};
\fill[popgreen] (0,-1.10) rectangle (0.35,-0.85); \node[right, font=\small] at (0.45,-0.975) {3. Bois : 12 \%};
\fill[poppurple] (0,-1.65) rectangle (0.35,-1.40); \node[right, font=\small] at (0.45,-1.525) {4. Coton : 7 \%};
\fill[poppink] (0,-2.20) rectangle (0.35,-1.95); \node[right, font=\small] at (0.45,-2.075) {5. Banane : 6 \%};
\fill[popgold] (0,-2.75) rectangle (0.35,-2.50); \node[right, font=\small] at (0.45,-2.625) {6. Aluminium : 5 \%};
\fill[popmuted] (0,-3.30) rectangle (0.35,-3.05); \node[right, font=\small] at (0.45,-3.175) {7. Autres (café, caoutchouc\dots) : 14 \%};
\end{scope}
\end{tikzpicture}
\end{center}
\legende{Structure approximative des exportations du Cameroun par produit, en \% de la valeur totale (ordres de grandeur, moyenne des années récentes ; sources : INS, Banque Mondiale). La prédominance du pétrole et des produits agricoles bruts est écrasante.}
\end{popfigure}

\begin{attention}[Brut ne veut pas dire « sans valeur »... mais presque]
Une matière première \cle{brute} ou \cle{peu transformée} se vend à un prix faible par rapport au produit fini qu'elle permet de fabriquer. Un kilo de fèves de cacao exporté rapporte au planteur et à l'État une fraction infime du prix de la tablette de chocolat fabriquée en Europe. On dit que le Cameroun capte une \textbf{faible valeur ajoutée} : c'est le cœur du problème de ses échanges extérieurs. Ne confonds pas « exporter beaucoup » et « gagner beaucoup ».
\end{attention}

\courssub{Les produits d'importation : l'industrie des autres}

\textbf{Observation.} Faisons maintenant l'inventaire des navires qui déchargent à Douala et à Kribi, et des camions qui franchissent les frontières. Le Cameroun importe essentiellement :

\begin{itemize}
\item des \cle{machines et équipements} (machines-outils, matériel de chantier, groupes électrogènes, équipements agricoles) ;
\item des \cle{véhicules} (voitures, camions, motos, pièces détachées), souvent d'occasion, venus d'Europe et d'Asie ;
\item des \cle{produits manufacturés} divers : textiles, chaussures, appareils électroménagers, téléphones, matériaux de construction (fer à béton, clinker pour les cimenteries) ;
\item des \cle{produits alimentaires} : riz (d'Asie), blé et farine (de France, de Russie), poisson congelé (sardines, maquereaux), huiles, sucre, lait en poudre ;
\item des \cle{hydrocarbures raffinés} : essence, gasoil, kérosène, carburéacteur ;
\item des \cle{médicaments} et produits pharmaceutiques (Inde, France, Chine) ;
\item des produits chimiques (engrais, plastiques).
\end{itemize}

\begin{definition}[Produit d'importation du Cameroun]
Les importations camerounaises sont dominées par des \cle{biens manufacturés} et des \cle{biens d'équipement}, c'est-à-dire des produits à forte \cle{valeur ajoutée}, issus de l'industrie des pays développés et émergents (Chine, France, Inde, Belgique, Turquie). S'y ajoutent des importations alimentaires croissantes, signe d'une agriculture vivrière insuffisante face à la croissance démographique et urbaine.
\end{definition}

\begin{attention}[Le paradoxe pétrolier camerounais]
Le Cameroun \textbf{exporte du pétrole brut} mais \textbf{importe de l'essence et du gasoil raffinés}. Comment est-ce possible ? La seule raffinerie du pays, la SONARA de Limbé (capacité nominale d'environ 2 millions de tonnes par an), est loin de couvrir la demande nationale, a été durablement arrêtée après l'incendie de mai 2019, et son outil industriel est vieillissant. Le pays vend donc son brut à bas prix relatif et rachète cher les produits raffinés. Ce paradoxe illustre parfaitement le drame de la non-transformation locale : la valeur ajoutée du raffinage est captée par l'étranger. À retenir pour tout commentaire de document sur la détérioration des termes de l'échange.
\end{attention}

\courssub{Des échanges de type Nord-Sud inégaux}

\textbf{Raisonnement.} Croisons les deux inventaires précédents. À l'exportation : des matières premières dont les prix sont volatils et tendanciellement faibles. À l'importation : des produits manufacturés dont les prix sont stables et élevés. La conclusion s'impose.

\begin{propriete}[Structure inégale des échanges du Cameroun]
Le Cameroun entretient avec le reste du monde des échanges de type \cle{Nord-Sud} : il exporte des produits primaires bruts (faible valeur ajoutée) et importe des produits transformés (forte valeur ajoutée). Cette structure explique la \cle{détérioration des termes de l'échange} : à quantité exportée égale, le pays peut acheter de moins en moins de produits importés au fil du temps, ce qui pèse durablement sur la balance commerciale.
\end{propriete}

\begin{popfigure}
\begin{center}
\begin{tikzpicture}
\begin{axis}[
ybar, bar width=0.9cm, width=0.85\linewidth, height=6.5cm,
ylabel={Part approximative (\%)},
symbolic x coords={Exportations, Importations},
xtick=data, ymin=0, ymax=100,
nodes near coords, every node near coord/.append style={font=\small},
legend style={at={(0.5,-0.18)}, anchor=north, legend columns=2},
]
\addplot[fill=popgreen] coordinates {(Exportations,85) (Importations,15)};
\addplot[fill=popblue] coordinates {(Exportations,15) (Importations,85)};
\legend{Produits primaires (bruts), Produits manufacturés (transformés)}
\end{axis}
\end{tikzpicture}
\end{center}
\legende{Le miroir inégal des échanges camerounais : environ 85 \% des exportations sont des produits primaires, tandis qu'environ 85 \% des importations sont des produits manufacturés ou transformés (ordres de grandeur ; sources : INS, Banque Mondiale).}
\end{popfigure}

\begin{savaistu}[Le cacao camerounais : champion du volume, nain de la transformation]
Le Cameroun figure parmi les cinq premiers producteurs mondiaux de cacao, avec une récolte souvent comprise entre 250 000 et 300 000 tonnes par an (ordre de grandeur). Pourtant, la très grande majorité des fèves part à l'état brut vers les chocolateries des Pays-Bas, de Belgique ou de Malaisie. Les unités locales de transformation (broyage, beurre de cacao) — comme les usines de la SIC Cacaos (groupe Barry Callebaut) à Douala — ne traitent qu'une part minoritaire de la récolte. Transformer davantage sur place créerait des emplois et multiplierait la valeur exportée : c'est tout l'enjeu de l'« industrialisation par la transformation des matières premières » prônée dans la stratégie de développement du pays.
\end{savaistu}

\courssub{Évolution récente : gaz, diversification timide et dépendance alimentaire}

La structure des échanges n'est pas figée. Trois tendances récentes méritent l'attention du candidat au Baccalauréat :

\begin{enumerate}
\item \textbf{L'arrivée du gaz naturel.} Depuis 2018, le Cameroun exporte du gaz naturel liquéfié depuis Kribi. Le gaz vient s'ajouter au pétrole brut et compenser partiellement le déclin de la production pétrolière nationale (passée d'environ 8 à 9 millions de tonnes par an dans les années 1980 à moins de 4 millions aujourd'hui, ordre de grandeur).
\item \textbf{Une diversification encore timide.} De nouveaux produits apparaissent : huile de palme, caoutchouc transformé, ciment exporté vers le Tchad et la Centrafrique, produits agricoles vers la CEMAC. Mais leur poids reste faible : la diversification est un objectif plus qu'une réalité.
\item \textbf{La montée inquiétante des importations alimentaires.} Riz, blé, poisson congelé : la facture alimentaire du Cameroun dépasse couramment plusieurs centaines de milliards de FCFA par an (ordre de grandeur : 500 à 900 milliards selon les années). Pays agricole par excellence, le Cameroun importe pourtant une part croissante de sa nourriture, sous l'effet de l'urbanisation rapide, des préférences alimentaires nouvelles (pain, riz) et de l'insuffisante modernisation de l'agriculture vivrière.
\end{enumerate}

\begin{methode}[Construire et légender un diagramme circulaire à partir d'un tableau statistique]
Au Baccalauréat, on peut te demander de représenter la structure des exportations par un diagramme en secteurs (camembert). Voici la démarche :
\begin{enumerate}
\item \textbf{Vérifier le total} : les parts doivent être exprimées en pourcentages et totaliser 100 \% (sinon, calcule chaque part : $\text{part} = \dfrac{\text{valeur du produit}}{\text{valeur totale}} \times 100$).
\item \textbf{Convertir en angles} : chaque secteur vaut $\text{angle} = \text{part} \times 3{,}6$ (car $100\,\% \leftrightarrow 360°$). Exemple : le pétrole à 38 \% donne $38 \times 3{,}6 \approx 137°$.
\item \textbf{Tracer} les secteurs du plus grand au plus petit, dans le sens des aiguilles d'une montre, en partant de midi.
\item \textbf{Choisir des couleurs} distinctes et les reporter dans la légende.
\item \textbf{Rédiger la légende} sous le graphique, en commençant par le titre (« Structure des exportations du Cameroun en \dots »), puis les catégories numérotées, puis la source et l'unité (\%).
\item \textbf{Interpréter} en une phrase : dégager le trait dominant (ici : prédominance des produits primaires).
\end{enumerate}
\end{methode}

\begin{exoresolu}[Calculer et interpréter la part des trois grands produits d'exportation]
\textbf{Énoncé.} Une année donnée, les exportations camerounaises totalisent environ 2 500 milliards de FCFA. Le pétrole brut en rapporte 950 milliards, le cacao 450 milliards et le bois 300 milliards. 1) Calcule la part de chacun de ces trois produits dans les exportations totales. 2) Calcule l'angle du secteur correspondant au pétrole dans un diagramme circulaire. 3) Tire une conclusion géographique.
\tcblower
\textbf{Solution détaillée.}
\begin{enumerate}
\item Parts respectives :
\begin{align*}
\text{Pétrole} &= \dfrac{950}{2500} \times 100 = 38\,\% \\
\text{Cacao} &= \dfrac{450}{2500} \times 100 = 18\,\% \\
\text{Bois} &= \dfrac{300}{2500} \times 100 = 12\,\%
\end{align*}
À eux trois, ces produits représentent $38 + 18 + 12 = 68\,\%$ des exportations.
\item Angle du secteur pétrolier : $38 \times 3{,}6 = 136{,}8°$, soit environ $137°$.
\item \textbf{Conclusion.} Près des deux tiers des recettes d'exportation du Cameroun proviennent de trois matières premières vendues brutes. Cette concentration rend l'économie du pays très vulnérable aux fluctuations des cours mondiaux : une baisse du prix du baril ou de la fève suffit à dégrader fortement la balance commerciale. D'où la nécessité de transformer localement et de diversifier.
\end{enumerate}
\end{exoresolu}

\courssub{Travaux pratiques guidés : classer les produits échangés}

\begin{experience}[TP — Construire le tableau de classement des produits échangés]
\textbf{Objectif.} Inventorier, classer et interpréter les produits des échanges extérieurs du Cameroun (savoir-faire officiel).

\textbf{Matériel.} La liste suivante de 14 produits, extraite des statistiques douanières : pétrole brut ; cacao en fèves ; téléphones portables ; bois en grumes ; riz ; coton fibre ; essence raffinée ; bananes ; camions ; café ; médicaments ; aluminium ; machines agricoles ; poisson congelé.

\textbf{Protocole.}
\begin{enumerate}
\item Construis un tableau à quatre colonnes : \emph{Exportations primaires}, \emph{Exportations manufacturées}, \emph{Importations primaires}, \emph{Importations manufacturées}.
\item Range chacun des 14 produits dans la bonne colonne.
\item Compte les produits de chaque colonne et compare.
\item Rédige un paragraphe de cinq lignes intitulé « Ce que révèle ce classement ».
\end{enumerate}

\textbf{Observations attendues.} Les exportations primaires comptent 6 produits (pétrole brut, cacao, bois, coton, bananes, café) contre 1 seul manufacturé (aluminium, et encore produit à partir d'alumine importée). Les importations manufacturées comptent 5 produits (téléphones, camions, médicaments, machines, essence raffinée) contre 2 primaires (riz, poisson).

\textbf{Interprétation.} Le tableau met en évidence la structure Nord-Sud des échanges : le Cameroun vend du brut et achète du transformé. L'essence raffinée, importée alors que le brut est exporté, illustre le déficit de transformation locale.
\end{experience}

\begin{exercice}[À toi de jouer]
À partir du tableau ci-dessous (valeurs fictives mais réalistes, en milliards de FCFA) :
\begin{center}
\begin{tabularx}{\linewidth}{|l|X|X|}\hline
\rowcolor{popblueL} \textbf{Produit} & \textbf{Exportations} & \textbf{Importations} \\ \hline
Pétrole brut et gaz & 950 & 0 \\ \hline
Produits pétroliers raffinés & 0 & 600 \\ \hline
Cacao & 450 & 0 \\ \hline
Produits alimentaires & 50 & 700 \\ \hline
Machines et véhicules & 0 & 900 \\ \hline
\end{tabularx}
\end{center}
1) Calcule le solde (exportations $-$ importations) pour chaque ligne. 2) Quelle ligne présente le déficit le plus lourd ? 3) Propose deux solutions pour réduire ce déficit.
\end{exercice}

\begin{corrige}[Exercice — Soldes par produit]
1) Soldes : pétrole brut et gaz : $+950$ ; produits raffinés : $-600$ ; cacao : $+450$ ; produits alimentaires : $-650$ ; machines et véhicules : $-900$ (milliards de FCFA). 2) Le déficit le plus lourd concerne les \textbf{machines et véhicules} ($-900$), suivi des produits alimentaires ($-650$). 3) Solutions : développer la transformation locale (raffinage, broyage du cacao, montage industriel) pour substituer la production nationale aux importations ; moderniser l'agriculture vivrière et les filières riz, maïs et pêche pour réduire la facture alimentaire. On peut aussi mentionner la promotion du « consommons local » et l'intégration régionale CEMAC.
\end{corrige}

\begin{aretenir}[L'essentiel de la section]
\begin{itemize}
\item Le Cameroun exporte des \cle{matières premières brutes} : pétrole et gaz, cacao, bois, coton, café, banane, aluminium, caoutchouc.
\item Il importe des \cle{produits manufacturés} (machines, véhicules, médicaments), des hydrocarbures raffinés et de plus en plus de \cle{produits alimentaires}.
\item Cette structure traduit des échanges \cle{Nord-Sud inégaux} : faible valeur ajoutée vendue, forte valeur ajoutée achetée — source de la détérioration des termes de l'échange.
\item Le paradoxe pétrolier (exportation de brut, importation de raffiné) et le cacao non transformé illustrent le déficit d'industrialisation.
\item L'évolution récente (gaz de Kribi, diversification timide, explosion de la facture alimentaire) montre que la structure bouge, mais lentement.
\end{itemize}
\end{aretenir}

Maintenant que nous savons \emph{quoi} le Cameroun échange, demandons-nous \emph{avec qui} ces échanges se font à l'échelle du continent : c'est l'objet de la section suivante, consacrée aux partenaires africains et à la CEMAC.

\coursec{Les échanges avec les partenaires africains et la CEMAC}

Dans la section précédente, nous avons dressé l'inventaire des produits que le Cameroun exporte (pétrole brut, cacao, café, coton, bois, bananes, aluminium) et importe (carburants raffinés, machines, produits pharmaceutiques, denrées comme le riz et le poisson congelé). Mais avec \emph{qui} le Cameroun échange-t-il ? Une partie de la réponse se trouve à nos portes : le continent africain, et d'abord la sous-région d'Afrique centrale. Imagine un camion de tomates et d'oignons quittant Bafoussam pour N'Djamena, ou un piroguier chargeant du savon et des tissus à Limbé vers Malabo : ce sont des échanges bien réels, mais qui pèsent étonnamment peu dans les statistiques officielles. C'est ce paradoxe que nous allons comprendre.

\courssub{La CEMAC : une communauté économique et monétaire}

\begin{definition}[La CEMAC et l'intégration régionale]
La \cle{CEMAC} (Communauté Économique et Monétaire de l'Afrique Centrale) est une organisation d'intégration régionale créée en 1994 (traité de N'Djamena), entrée en vigueur en 1999, qui succède à l'UDEAC. Elle regroupe six États : le \cle{Cameroun}, le \cle{Tchad}, la \cle{Centrafrique} (RCA), le \cle{Gabon}, le \cle{Congo} (Brazzaville) et la \cle{Guinée équatoriale}. Ses objectifs principaux sont la création d'un \cle{marché commun} (libre circulation des personnes, des biens et des capitaux), l'harmonisation des politiques économiques et le maintien d'une \cle{monnaie commune}, le \cle{franc CFA} (FCFA), gérée par la \cle{BEAC} (Banque des États de l'Afrique Centrale), dont le siège est à Yaoundé. On parle d'\cle{intégration régionale} lorsque des États voisins décident de coordonner leurs politiques économiques et de supprimer progressivement les barrières qui freinent les échanges entre eux.
\end{definition}

L'intuition est simple : six pays, souvent petits démographiquement (la Centrafrique compte environ 5 millions d'habitants, le Gabon à peine plus de 2 millions), ont tout intérêt à s'unir pour peser davantage dans le commerce mondial, comme des commerçants d'un même marché qui s'associent pour acheter en gros. La monnaie commune évite en outre les coûts de change entre eux : un commerçant camerounais qui vend du savon à Libreville est payé dans la même monnaie qu'à Douala.

\begin{attention}[Ne pas confondre CEMAC et CEEAC]
Piège classique du Baccalauréat : la \cle{CEMAC} est une organisation \emph{économique et monétaire} à \textbf{six membres}. La \cle{CEEAC} (Communauté Économique des États de l'Afrique Centrale) est une organisation \emph{politique} plus large (onze membres, dont l'Angola, le Rwanda, le Burundi, la RDC, São Tomé-et-Principe). Le Nigeria, lui, n'appartient à \textbf{aucune des deux} : il est membre de la CEDEAO. Dans une copie, écrire « le Nigeria, membre de la CEMAC » est une erreur éliminatoire.
\end{attention}

\courssub{Le Cameroun, première économie et « grenier » de la sous-région}

Au sein de la CEMAC, le Cameroun occupe une position dominante : il représente, à lui seul, environ \textbf{40 \% du PIB} de la communauté et près de la moitié de sa population (environ 28 millions d'habitants sur un total proche de 60 millions ; ce sont des ordres de grandeur, les chiffres variant selon les sources et les années). C'est la \cle{première économie} de la sous-région.

Cette puissance relative s'explique par la \emph{diversité} de l'économie camerounaise : alors que le Gabon, le Congo, le Tchad et la Guinée équatoriale dépendent massivement du pétrole, le Cameroun dispose d'une agriculture vivrière dynamique, d'un secteur industriel léger et d'un port (Douala) qui sert de débouché maritime aux pays enclavés (Tchad, Centrafrique). Le Cameroun joue ainsi le rôle de \cle{fournisseur de la sous-région} :

\begin{itemize}
\item \textbf{Produits alimentaires} : tomates, oignons, pommes de terre, maïs, bananes plantains, huile de palme, poissons fumés, sucre (SOSUCAM), boissons (brasseries, eaux minérales) partent vers le Tchad, la Centrafrique, le Gabon et la Guinée équatoriale ;
\item \textbf{Produits manufacturés légers} : savons, tissus, chaussures, matériaux de construction (ciment de CIMENCAM et Dangote), produits plastiques, cosmétiques ;
\item \textbf{Services de transit} : le corridor Douala--N'Djamena et l'axe Douala--Bangui acheminent les importations des pays enclavés.
\end{itemize}

En sens inverse, le Cameroun importe peu de la CEMAC : essentiellement du bétail et de la viande séchée (Tchad), du poisson de lac, et quelques produits artisanaux. Les échanges sont donc \emph{déséquilibrés en faveur du Cameroun} à l'intérieur de la communauté.

\begin{popfigure}
\includegraphics[width=\linewidth,height=9cm,keepaspectratio]{N16/cemac.png}
\legende{La CEMAC (Communauté économique et monétaire de l'Afrique centrale) sur la carte de l'Afrique : les six États membres sont en vert (Tchad, République centrafricaine, Cameroun, Guinée équatoriale, Gabon, Congo) ; les autres pays africains sont en gris, dont le Nigeria à l'ouest, voisin du Cameroun. Le Cameroun est le pays vert situé à l'ouest, ouvert sur le golfe de Guinée, par où partent les flux vers le Tchad et la RCA (corridor Douala--N'Djamena, le plus chargé) et vers le Gabon et le Congo. La carte ne représente ni les flux ni leur importance ; l'élève les ajoute sur son croquis par des flèches d'épaisseur proportionnelle. \\ Source : Wikimedia Commons, Sven-Steffen Arndt, AnzeeTM et Hoshie, domaine public.}
\end{popfigure}

\courssub{Le Nigeria, partenaire majeur hors CEMAC}

Le \cle{Nigeria}, voisin occidental du Cameroun, n'est pas membre de la CEMAC mais de la \cle{CEDEAO} (Communauté Économique des États de l'Afrique de l'Ouest). C'est pourtant l'un des tout premiers partenaires africains du Cameroun, et de loin le plus important marché frontalier : avec plus de 200 millions d'habitants, le Nigeria est un « ventre » immense à quelques kilomètres de nos frontières.

Les échanges avec le Nigeria prennent deux formes :

\begin{itemize}
\item \textbf{Échanges officiels} : le Cameroun importe du Nigeria des produits manufacturés (plastiques, matériaux, médicaments, pièces détachées) et des produits pétroliers raffinés ; il exporte vers le Nigeria des produits agricoles (maïs, haricots niébé, oignons, bananes) et du bétail ;
\item \textbf{Échanges informels} : une part considérable du commerce bilatéral échappe aux statistiques douanières. Carburant de contrebande (« \emph{funta} » ou essence nigériane revendue au Cameroun), friperie, denrées alimentaires, motos et appareils électroniques circulent par les marchés frontaliers et les pistes, surtout dans le Nord-Ouest, le Sud-Ouest et l'Extrême-Nord (nous y reviendrons en détail dans le dossier 5).
\end{itemize}

\begin{savaistu}[Un commerce intra-africain structurellement faible]
Sur l'ensemble du continent africain, les échanges \emph{entre pays africains} ne représentent souvent qu'environ \textbf{15 \% du commerce total} de l'Afrique (ordre de grandeur donné par la CNUCED ; la proportion exacte varie selon les années), contre près de 60 \% en Europe et environ 40 \% en Asie. C'est l'une des raisons d'être de la \cle{ZLECAF} (Zone de Libre-Échange Continentale Africaine), lancée en 2019 et opérationnelle depuis 2021 : créer un marché unique de plus d'un milliard de consommateurs en supprimant progressivement les droits de douane entre pays africains.
\end{savaistu}

\courssub{La faiblesse des échanges intra-CEMAC : un paradoxe à expliquer}

Voici le paradoxe central de cette leçon : malgré une monnaie commune, des frontières théoriquement ouvertes et des besoins complémentaires, le commerce \emph{intra-CEMAC} ne représente qu'une part très faible du commerce total des pays membres — de l'ordre de \textbf{3 à 5 \%} selon les estimations (ordre de grandeur à retenir, les chiffres variant selon les sources). Autrement dit, le Cameroun échange beaucoup plus avec la Chine ou l'Union européenne qu'avec le Gabon ou le Tchad. Comment l'expliquer ?

\begin{enumerate}
\item \textbf{Des productions similaires et peu complémentaires} : les économies de la sous-région exportent toutes des matières premières brutes (pétrole pour cinq des six membres, bois, cacao). Un pays pétrolier n'a pas besoin du pétrole de son voisin : il n'y a donc presque rien à s'échanger entre eux.
\item \textbf{L'état déplorable des infrastructures} : les routes inter-États sont peu nombreuses et souvent dégradées (l'axe Douala--Bangui reste en partie non bitumé ; le corridor vers N'Djamena est long et coûteux). Il n'existe aucune liaison ferroviaire entre les pays de la CEMAC.
\item \textbf{Les lenteurs et tracasseries douanières} : malgré la libre circulation proclamée, les contrôles multiples, les taxes informelles aux frontières et les barrages routiers allongent les délais et renchérissent les coûts de transaction.
\item \textbf{L'orientation historique des échanges vers l'extérieur} : héritage de la colonisation, les flux (ports, routes, monnaie adossée à l'euro) restent tournés vers l'Europe, puis vers l'Asie.
\item \textbf{Le poids du secteur informel} : une partie réelle des échanges existe mais n'apparaît pas dans les statistiques, ce qui sous-estime le commerce régional.
\end{enumerate}

\begin{popfigure}
\begin{center}
\begin{tikzpicture}
\begin{axis}[
  ybar, bar width=1.6cm,
  width=0.85\linewidth, height=6.5cm,
  ymin=0, ymax=100,
  ylabel={Part du commerce total du Cameroun (\%)\%},
  symbolic x coords={Afrique, Reste du monde},
  xtick=data,
  ytick={0,20,40,60,80,100},
  nodes near coords, nodes near coords align={vertical},
  every node near coord/.append style={font=\small\bfseries},
  axis line style={popdark},
  tick label style={font=\small},
]
\addplot[fill=popgreen, draw=popdark] coordinates {(Afrique,10)};
\addplot[fill=popblue, draw=popdark] coordinates {(Reste du monde,90)};
\end{axis}
\end{tikzpicture}
\end{center}
\legende{Histogramme : part approximative du commerce du Cameroun avec l'Afrique (environ 10 \%, dont seulement 3 à 5 \% avec la CEMAC) et avec le reste du monde (environ 90 \%). Ordres de grandeur indicatifs, variables selon les années et les sources (INS, Banque mondiale).}
\end{popfigure}

\courssub{Les autres partenaires africains : CEDEAO et ZLECAF}

Au-delà du Nigeria et de la CEMAC, le Cameroun entretient des relations commerciales avec d'autres ensembles africains :

\begin{itemize}
\item la \cle{CEDEAO} (quinze membres ouest-africains) : outre le Nigeria, le Cameroun échange avec le Ghana, la Côte d'Ivoire ou le Sénégal (produits manufacturés, cacao transformé, services) ;
\item la \cle{ZLECAF} : ratifiée par le Cameroun, elle promet à terme la suppression des droits de douane sur 90 \% des produits échangés entre pays africains. Pour le Cameroun, c'est une opportunité (écouler ses produits manufacturés légers sur un marché continental) mais aussi un risque (concurrence des industries nigérianes, égyptiennes ou sud-africaines plus compétitives).
\end{itemize}

\begin{methode}[Légender une carte de flux (savoir-faire exigible)]
Pour construire ou légender un croquis de flux commerciaux au Baccalauréat :
\begin{enumerate}
\item \textbf{Titrer} la carte de manière complète (quoi ? où ? quand ?) : « Les flux commerciaux du Cameroun vers la CEMAC » ;
\item \textbf{Orienter} : placer une flèche du nord ;
\item \textbf{Localiser} : nommer les pays, les principales villes (ports, capitales) et les axes (corridors) ;
\item \textbf{Choisir des figurés adaptés} : des \emph{flèches d'épaisseur proportionnelle} à l'importance du flux (la plus épaisse pour le flux dominant, ici Douala--N'Djamena) ; des points pour les villes, des teintes différentes pour les pays ;
\item \textbf{Rédiger la légende} dans un cartouche : chaque figuré utilisé doit y être expliqué, dans l'ordre logique (d'abord les flux, puis les villes, puis les pays) ;
\item \textbf{Vérifier la hiérarchie} : au moins trois épaisseurs de flèches pour montrer les flux majeurs, moyens et faibles.
\end{enumerate}
Erreur à éviter : dessiner des flèches toutes identiques — le correcteur attend précisément la \emph{proportionnalité}.
\end{methode}

\begin{exoresolu}[Expliquer la faiblesse des échanges intra-CEMAC]
Énoncé : « Les échanges intra-CEMAC ne représentent que 3 à 5 \% du commerce total des pays membres. À l'aide de vos connaissances, proposez trois explications à cette faiblesse, en illustrant chacune d'un exemple précis. »
\tcblower
Solution détaillée :
\begin{enumerate}
\item \textbf{Similitude des productions} : les économies membres exportent toutes des matières premières brutes. Exemple : le Gabon et le Congo, producteurs de pétrole, n'ont aucun besoin du pétrole tchadien ; la complémentarité commerciale est donc quasi nulle.
\item \textbf{Insuffisance des infrastructures de transport} : les axes inter-États sont rares et dégradés. Exemple : la route Douala--Bangui (environ 1 500 km) comporte de longs tronçons non bitumés, ce qui renchérit le coût des marchandises camerounaises livrées en Centrafrique.
\item \textbf{Tracasseries douanières et informel} : malgré la libre circulation proclamée par la CEMAC, les contrôles et taxes informelles aux frontières freinent les flux officiels. Exemple : une partie du commerce entre le Cameroun et le Gabon transite par les marchés frontaliers (Kyé-Ossi) et échappe aux statistiques, ce qui minore mécaniquement la part mesurée des échanges.
\end{enumerate}
Conclusion attendue : la faiblesse des échanges intra-CEMAC tient moins à l'absence de volonté politique qu'à des économies peu complémentaires et à des obstacles matériels persistants.
\end{exoresolu}

\begin{experience}[TP : localisation des pays CEMAC et des axes d'échanges sur carte muette]
\textbf{Matériel} : carte muette de l'Afrique centrale, crayons de couleur, légende vierge, atlas.
\textbf{Protocole guidé} :
\begin{enumerate}
\item Repérer et colorier différemment les six pays de la CEMAC ; hachurer le Nigeria (hors CEMAC) d'une couleur distincte ;
\item Placer et nommer : Douala, Yaoundé, N'Djamena, Bangui, Libreville, Brazzaville, Malabo, Lagos ;
\item Tracer les corridors : Douala--N'Djamena (route + voie ferrée jusqu'à Ngaoundéré), Douala--Bangui, l'axe côtier vers Bata et Libreville ;
\item Dessiner trois flèches d'épaisseur décroissante figurant les flux camerounais vers le Tchad, la RCA et le Gabon ;
\item Rédiger la légende et le titre, placer le nord.
\end{enumerate}
\textbf{Observations attendues} : les flux convergent vers Douala, seul grand port de la sous-région ; les pays enclavés (Tchad, RCA) dépendent du territoire camerounais pour leurs importations.
\textbf{Interprétation} : la position centrale et littorale du Cameroun explique son rôle de plaque tournante commerciale de l'Afrique centrale.
\end{experience}

\begin{exercice}[Pour t'entraîner]
1. Définis la CEMAC et cite ses six membres de mémoire.\par
2. Explique pourquoi le Nigeria est un partenaire majeur du Cameroun sans être membre de la CEMAC.\par
3. Sur un croquis, représente par des flèches proportionnelles les flux du Cameroun vers le Tchad, la Centrafrique et le Gabon, avec titre, nord et légende.\par
4. Vrai ou faux (justifie) : « La CEEAC est la zone monétaire du franc CFA. »
\end{exercice}

\begin{corrige}[Exercice]
1. La CEMAC est la Communauté Économique et Monétaire de l'Afrique Centrale, organisation d'intégration régionale visant un marché commun et dotée d'une monnaie commune (le FCFA). Membres : Cameroun, Tchad, Centrafrique, Gabon, Congo, Guinée équatoriale.\par
2. Le Nigeria est membre de la CEDEAO (Afrique de l'Ouest), pas de la CEMAC. Mais sa proximité géographique et son immense marché (plus de 200 millions d'habitants) en font un partenaire majeur : le Cameroun y exporte des produits agricoles et en importe des manufacturés et du carburant, une grande partie des flux étant informelle.\par
3. Le croquis doit comporter trois flèches d'épaisseurs décroissantes (Tchad $>$ RCA $>$ Gabon), partant de Douala, un titre complet, la flèche du nord et une légende ordonnée.\par
4. Faux : la zone monétaire du franc CFA d'Afrique centrale correspond à la CEMAC (avec la BEAC) ; la CEEAC est une organisation politique élargie à onze États, sans monnaie commune.
\end{corrige}

\begin{aretenir}[L'essentiel de la section]
La CEMAC regroupe six pays d'Afrique centrale autour d'un marché commun et du franc CFA. Le Cameroun, première économie de la communauté, fournit ses voisins en produits alimentaires et manufacturés légers et leur offre le port de Douala. Le Nigeria (CEDEAO) est le principal partenaire africain hors CEMAC, avec d'importants échanges informels. Pourtant, le commerce intra-CEMAC reste très faible (3 à 5 \% du commerce total) en raison de productions similaires, d'infrastructures défaillantes et de lenteurs douanières — un paradoxe que la ZLECAF ambitionne de corriger à l'échelle continentale.
\end{aretenir}

\coursec{Les échanges avec le reste du monde}

Nous venons d'analyser les échanges du Cameroun avec ses partenaires africains, notamment au sein de la CEMAC : ils restent limités en valeur et portent surtout sur des produits de première nécessité et des hydrocarbures. Or, si tu regardes l'étiquette d'un téléphone portable vendu à Yaoundé, l'origine du cacao expédié depuis le port de Douala, ou la facture d'une machine importée pour une usine d'Ebolowa, tu constates une évidence : l'essentiel du commerce extérieur du Cameroun se joue bien au-delà du continent africain. C'est ce commerce avec le \cle{reste du monde} — Europe, Asie, Amérique — que nous étudions maintenant.

\courssub{Un commerce tourné vers l'extérieur : le caractère extraverti des échanges}

\begin{definition}[Extraversion des échanges]
On dit que les échanges d'un pays sont \cle{extravertis} lorsqu'ils sont orientés essentiellement vers l'extérieur du pays et de sa région d'appartenance, au détriment des échanges internes et régionaux. Le commerce extérieur camerounais est structurellement extraverti : plus de 80 \% de sa valeur se réalise avec des partenaires non africains, héritage direct de l'économie coloniale de traite organisée pour approvisionner l'Europe en matières premières.
\end{definition}

Ce caractère extraverti s'accompagne d'une \cle{asymétrie} fondamentale : le Cameroun \textbf{vend des matières premières} (pétrole brut, cacao, bois, coton, banane, café, aluminium) et \textbf{achète des produits finis ou semi-finis} (machines, véhicules, médicaments, produits alimentaires transformés, textiles, équipements électriques et électroniques). La valeur ajoutée — c'est-à-dire la richesse créée par la transformation — est donc captée à l'étranger. Un kilogramme de cacao exporté brut rapporte infiniment moins qu'une tablette de chocolat importée.

\begin{aretenir}[L'essentiel de la structure des échanges]
Le Cameroun exporte des produits primaires peu transformés vers l'Europe et l'Asie, et importe des biens manufacturés, des équipements et des produits alimentaires. Cette structure, héritée de la colonisation, explique la fragilité de sa balance commerciale et la détérioration de ses termes de l'échange (notion étudiée en section 5).
\end{aretenir}

\courssub{Les partenaires européens : l'Union européenne, débouché historique}

L'\cle{Union européenne} (UE) est le partenaire commercial historique du Cameroun, héritage des liens tissés depuis la période coloniale et consolidés par les conventions de Yaoundé (1963, 1969) puis de Lomé (1975-2000). Elle reste, prise comme bloc, le \textbf{premier client} du Cameroun : elle absorbe, selon les années, entre 35 \% et 45 \% des exportations camerounaises (ordre de grandeur basé sur les données Douanes/INS 2018-2022, parts variant avec les cours du pétrole).

Les principaux partenaires européens sont :

\begin{itemize}
\item la \textbf{France}, ancienne puissance administratrice : elle est à la fois un client (pétrole, bois, cacao) et surtout un \textbf{fournisseur} majeur (machines, véhicules, produits pharmaceutiques, farine de blé, produits agroalimentaires) ;
\item les \textbf{Pays-Bas}, premier débouché mondial du \cle{cacao} camerounais : le port d'Amsterdam est la première place cacaoyère du monde, où aboutissent les fèves embarquées à Douala ;
\item l'\textbf{Espagne} et l'\textbf{Italie}, gros importateurs de \cle{bois} (grumes et sciages) et de pétrole brut ;
\item le \textbf{Portugal} et la \textbf{Belgique}, débouchés secondaires du bois et du cacao.
\end{itemize}

L'Europe est ainsi le grand débouché des trois produits d'exportation emblématiques du Cameroun : le \textbf{cacao} (vers les Pays-Bas), la \textbf{banane} (vers la France et l'Espagne, via la filière PHP — Plantations du Haut Penja — et la marque commercialisée par Compagnie fruitière), et le \textbf{bois} (vers l'Espagne, l'Italie et la France).

\begin{savaistu}[Des conventions de Lomé aux APE]
Depuis les accords de Lomé (1975), les pays ACP (Afrique, Caraïbes, Pacifique) exportaient vers l'Europe sans droits de douane, sans réciprocité. L'Organisation mondiale du commerce (OMC) ayant jugé ce régime discriminatoire, il a été remplacé par les \cle{Accords de partenariat économique} (APE), fondés sur la réciprocité progressive. Le Cameroun a signé un APE « intérimaire » avec l'UE (entrée en application provisoire en août 2014), qui garantit l'accès libre et sans contingent de ses exportations au marché européen, tout en l'engageant à ouvrir graduellement (sur 15 à 25 ans) son propre marché aux produits européens — ce qui inquiète les industriels locaux, peu compétitifs face aux importations subventionnées.
\end{savaistu}

\courssub{La montée en puissance de la Chine}

Le fait majeur des vingt dernières années est l'irruption de la \cle{Chine}, devenue, selon les années et selon que l'on compte l'UE comme bloc ou pays par pays, le \textbf{premier ou le deuxième partenaire commercial} du Cameroun. La Chine est désormais le \textbf{premier fournisseur} du Cameroun pris pays par pays (environ 15 \% à 20 \% des importations, ordre de grandeur 2018-2022) : elle inonde les marchés de Douala, Yaoundé et Bafoussam de produits manufacturés bon marché — textiles, chaussures, appareils électroménagers, motos, téléphones, matériaux de construction, machines-outils.

En sens inverse, la Chine est devenue un client important : elle achète du \textbf{pétrole brut}, du \textbf{bois} (notamment des grumes, malgré les restrictions à l'exportation de grumes non transformées entrées en vigueur en 2017), du \textbf{coton} fibre et du caoutchouc. Les entreprises chinoises (Sinohydro, CMEC, etc.) sont en outre très présentes dans les grands travaux d'infrastructures (port en eau profonde de Kribi, autoroutes, stades, barrages), souvent financés par des prêts concessionnels chinois remboursés en partie par des livraisons de matières premières (mécanisme « ressources contre infrastructures »).

\begin{exemplebox}[Chine et Union européenne : poids comparés (ordres de grandeur 2018-2022)]
Pour une année récente typique, la structure du commerce camerounais peut se résumer ainsi (source : Douanes Cameroun, INS, UN Comtrade) :
\begin{itemize}
\item \textbf{Exportations (clients)} : UE environ 40 \% (dont Pays-Bas ~15 \%, Espagne ~8 \%, Italie ~8 \%, France ~5 \%), Chine ~10 \%, Inde ~8 \%, États-Unis ~5 \%, Autres ~37 \% ;
\item \textbf{Importations (fournisseurs)} : Chine ~18 \%, UE ~30 \% (dont France ~10 \%), Inde ~6 \%, États-Unis ~4 \%, Autres ~42 \%.
\end{itemize}
Ces chiffres sont des \textbf{ordres de grandeur} : ils fluctuent fortement avec le prix du pétrole, premier produit exporté en valeur. Ce qui est structurel, en revanche, c'est la progression continue de la part chinoise dans les importations depuis les années 2000.
\end{exemplebox}

\begin{savaistu}[La Chine, nouveau partenaire majeur de l'Afrique centrale]
En moins de vingt ans, la Chine est devenue le premier partenaire commercial de l'ensemble de l'Afrique, détrônant les anciennes puissances coloniales. En Afrique centrale, elle achète pétrole, bois et minerais, et vend des manufacturés à bas prix. Ce « partenariat Sud-Sud », présenté comme gagnant-gagnant, reproduit néanmoins la même asymétrie : l'Afrique exporte du brut et importe du transformé. Les produits chinois bon marché ont aussi mis en difficulté les industries textiles locales (ex. CICAM à Douala) et les petits commerçants de détail.
\end{savaistu}

\courssub{Les autres partenaires : Inde, États-Unis, Asie, Turquie}

Au-delà du duo Europe-Chine, d'autres partenaires comptent :

\begin{itemize}
\item l'\textbf{Inde}, cliente du pétrole brut, du bois et du coton, et fournisseuse de produits pharmaceutiques (génériques), de riz et de machines ;
\item les \textbf{États-Unis}, qui achètent du pétrole et du cacao et exportent vers le Cameroun des véhicules, des machines et des produits agroalimentaires. Le Cameroun bénéficie de l'\cle{AGOA} (\emph{African Growth and Opportunity Act}, loi américaine de 2000), qui ouvre le marché états-unien, en franchise de droits, à des milliers de produits africains. Le Cameroun en a surtout profité pour ses exportations pétrolières ; son éligibilité a été suspendue en janvier 2020 pour des raisons de droits de l'homme, illustrant la fragilité et la conditionnalité politique de ces préférences unilatérales ;
\item les autres pays asiatiques : la \textbf{Corée du Sud}, la \textbf{Malaisie} et l'\textbf{Indonésie} (achats de bois et d'huile de palme, ventes de manufacturés et d'huile de palme raffinée), le \textbf{Japon} (véhicules et équipements) ;
\item la \textbf{Turquie} et les pays du \textbf{Golfe} (Émirats, Arabie saoudite, Qatar), partenaires émergents (céréales, produits sidérurgiques, ciment, investissements portuaires et hôteliers).
\end{itemize}

\begin{attention}[Client ou fournisseur : ne pas confondre !]
Un \cle{client} est un pays \textbf{à qui le Cameroun exporte} (il achète nos produits) ; un \cle{fournisseur} est un pays \textbf{de qui le Cameroun importe} (il nous vend ses produits). Un même pays peut être les deux à la fois : la France est un fournisseur majeur et un client notable ; la Chine est le premier fournisseur (pays par pays) et un client en progression. Au Bac, quand on te demande de classer les partenaires, précise toujours si tu raisonnes sur les exportations (clients) ou sur les importations (fournisseurs) : mélanger les deux colonnes d'un tableau statistique est l'erreur la plus fréquente.
\end{attention}

\courssub{Cartographier les flux : routes maritimes mondiales et diagramme des partenaires}

La carte ci-dessous montre les grandes routes maritimes mondiales sur lesquelles circulent les marchandises du Cameroun avec le reste du monde ; le diagramme qui suit en précise les partenaires. Pour construire ta propre carte de flux, retiens ce \textbf{code couleur harmonisé : flèches bleues (exportations), flèches oranges (importations).}

\begin{popfigure}
\includegraphics[width=\linewidth,height=8cm,keepaspectratio]{N16/shipping_routes.jpg}
\legende{Le trafic maritime commercial mondial (les continents sont en noir ; plus le bleu est foncé, plus le trafic est dense). Les exportations camerounaises (pétrole, cacao, bois, coton, banane) et les importations (produits manufacturés, machines, véhicules, produits alimentaires) empruntent ces routes depuis le golfe de Guinée, sur la façade atlantique de l'Afrique : vers l'Europe et l'Amérique du Nord par l'Atlantique, vers l'Asie par le cap de Bonne-Espérance ou par Suez. Les partenaires du Cameroun (Union européenne, Chine, États-Unis) et les parts de chaque flux ne figurent pas sur la carte : ils proviennent du tableau et du diagramme de la leçon (ordres de grandeur 2018-2022). \\ Source : Wikimedia Commons, T. Hengl, CC BY-SA 3.0.}
\end{popfigure}

\begin{popfigure}
\begin{center}
\begin{tikzpicture}
\begin{axis}[
  ybar, bar width=9pt,
  width=0.95\linewidth, height=6.2cm,
  ymin=0, ymax=45,
  ylabel={Part (\%)},
  symbolic x coords={UE, Chine, Inde, États-Unis, Autres},
  xtick=data,
  legend style={at={(0.98,0.95)},anchor=north east,font=\footnotesize},
  ymajorgrids=true, grid style={popline!50},
  nodes near coords, every node near coord/.append style={font=\scriptsize}
]
\addplot[fill=popblue,draw=popdark] coordinates {(UE,40) (Chine,10) (Inde,8) (États-Unis,5) (Autres,37)};
\addplot[fill=poporange,draw=popdark] coordinates {(UE,30) (Chine,18) (Inde,6) (États-Unis,4) (Autres,42)};
\legend{Exportations (clients), Importations (fournisseurs)}
\end{axis}
\end{tikzpicture}
\end{center}
\legende{Principaux clients et fournisseurs du Cameroun (parts en \% de la valeur des échanges, ordres de grandeur indicatifs 2018-2022 ; l'UE est comptée comme bloc). Lecture : l'UE est le premier client, la Chine le premier fournisseur pris pays par pays. Source : Douanes Cameroun, INS, UN Comtrade.}
\end{popfigure}

\begin{methode}[Construire une carte de flux à partir d'un tableau de partenaires]
\begin{enumerate}
\item \textbf{Lis le tableau} en distinguant les deux colonnes : exportations (clients) et importations (fournisseurs). Repère les 3 à 5 partenaires dominants de chaque colonne.
\item \textbf{Choisis deux couleurs} : une pour les exportations (ex. bleu), une pour les importations (ex. orange). \textbf{Garde-les constantes} sur toute la carte et le diagramme associé.
\item \textbf{Localise} le Cameroun et chaque partenaire sur un fond de carte (même très simplifié) ; place une flèche du nord et une échelle indicative.
\item \textbf{Trace les flux} par des flèches dont l'épaisseur est proportionnelle à la part du partenaire (par exemple 1 mm pour 10 \%). Oriente les flèches : du Cameroun vers le client (exportation), du fournisseur vers le Cameroun (importation).
\item \textbf{Rédige la légende} : elle doit permettre de lire la carte sans le texte (symboles, couleurs, échelle des épaisseurs, unité, date).
\item \textbf{Donne un titre} complet : phénomène, espace, date.
\end{enumerate}
\end{methode}

\begin{exoresolu}[Lecture d'un tableau de partenaires commerciaux]
\textbf{Énoncé.} On donne pour une année récente : exportations du Cameroun — UE 40 \%, Chine 10 \%, Inde 8 \% ; importations — Chine 18 \%, UE 30 \%, Inde 6 \%. 1) Qui est le premier client du Cameroun ? Son premier fournisseur ? 2) Calcule le solde apparent des échanges avec la Chine en points de pourcentage et commente.
\tcblower
\textbf{Solution.} 1) Le premier \textbf{client} est l'Union européenne (40 \% des exportations, portée par les Pays-Bas pour le cacao, l'Espagne et l'Italie pour le bois). Le premier \textbf{fournisseur} pris pays par pays est la Chine (18 \% des importations), devant la France ; si l'on compte l'UE comme bloc, c'est l'UE (30 \%) qui est le premier fournisseur. 2) Avec la Chine, le Cameroun exporte 10 \% de ses ventes mais y puise 18 \% de ses achats : l'écart est de $10 - 18 = -8$ points. Le Cameroun achète donc à la Chine beaucoup plus qu'il ne lui vend : la relation est structurellement déficitaire, ce qui illustre l'asymétrie des échanges — manufacturés chinois diversifiés contre matières premières camerounaises en volume et valeur insuffisants.
\end{exoresolu}

\begin{attention}[Pièges de rédaction au Bac]
Évite trois erreurs classiques : (1) affirmer que la France est « le » partenaire du Cameroun sans nuance — c'était vrai il y a trente ans, mais la Chine et l'UE prise comme bloc l'ont dépassée ; (2) confondre la valeur et le volume des échanges (le pétrole pèse lourd en valeur, peu en volume) ; (3) oublier que les chiffres varient fortement d'une année à l'autre avec le cours du pétrole : cite toujours des \textbf{ordres de grandeur} et \textbf{date tes données} (ex. « moyenne 2018-2022 »).
\end{attention}

En définitive, le commerce du Cameroun avec le reste du monde est dominé par deux pôles — l'Union européenne, débouché historique de ses matières premières, et la Chine, pourvoyeuse de manufacturés et financeur d'infrastructures — et il reste profondément asymétrique : le pays vend brut et achète transformé. Cette structure est à l'origine des problèmes que nous analyserons dans la prochaine section : déficit structurel de la balance commerciale et détérioration des termes de l'échange.

\coursec{Les problèmes des échanges extérieurs du Cameroun}

Dans la section précédente, nous avons vu que le Cameroun échange avec le reste du monde : il exporte du pétrole brut, du cacao, du bois, du coton, de l'aluminium, et il importe des biens d'équipement, des véhicules, des produits pharmaceutiques, du riz, du poisson congelé. À première vue, tout semble fonctionner : les navires entrent et sortent du port de Douala, les marchés sont approvisionnés. Pourtant, derrière cette apparente fluidité se cachent de graves déséquilibres. Imagine un commerçant du marché Mokolo qui vend chaque année ses tomates de moins en moins cher, tout en achetant ses brouettes et ses arrosoirs de plus en plus cher : même s'il vend beaucoup, il s'appauvrit. C'est exactement la situation du Cameroun dans le commerce mondial. Cette section analyse les \cle{problèmes} qui fragilisent les échanges extérieurs du pays, puis esquisse les pistes de solutions.

\courssub{Le déficit chronique de la balance commerciale}

\textbf{Observation de départ.} Au port de Douala, les grues déchargent des milliers de conteneurs : engins de chantier, voitures d'occasion, sacs de riz, médicaments, téléviseurs. En sens inverse, on charge surtout des produits bruts : grumes de bois, fèves de cacao, pétrole. La valeur de ce qui entre dépasse souvent la valeur de ce qui sort.

\begin{definition}[Balance commerciale et déficit]
La \cle{balance commerciale} est la différence entre la valeur des exportations et la valeur des importations de biens d'un pays sur une période donnée (généralement une année). Elle est \textbf{excédentaire} si les exportations dépassent les importations, \textbf{déficitaire} dans le cas contraire. Le Cameroun connaît un déficit \emph{chronique}, c'est-à-dire récurrent d'année en année (de l'ordre de plusieurs centaines de milliards de FCFA, voire plus de 1 000 milliards certaines années comme en 2022 ; l'ordre de grandeur varie fortement selon le cours du pétrole).
\end{definition}

\textbf{Les causes du déficit.} Deux postes d'importation pèsent lourd :

\begin{itemize}
\item Les \cle{biens d'équipement} : machines, engins, véhicules, matériel électrique, indispensables aux chantiers (routes, barrages, stades) et aux entreprises, mais que le Cameroun ne fabrique pas.
\item Les \cle{produits alimentaires} : riz, blé (farine), poisson congelé, huile, sucre, lait. Le pays importe une part importante de sa consommation alimentaire, alors qu'il possède d'immenses terres arables et deux façades maritimes poissonneuses.
\end{itemize}

\begin{savaistu}[Importer du riz et du poisson au pays des terres fertiles]
Le Cameroun dispose d'environ 7 millions d'hectares de terres arables et d'un potentiel halieutique important (côte atlantique, fleuves, lac Tchad). Pourtant, il importe chaque année 600 000 à 900 000 tonnes de riz (en provenance surtout d'Asie : Thaïlande, Inde, Vietnam) et des quantités massives de poisson congelé (capitaine, maquereau, souvent pêché au large de l'Afrique de l'Ouest par des flottes étrangères). La facture alimentaire atteint 300 à 500 milliards de FCFA par an (ordre de grandeur). C'est le paradoxe d'un pays agricole qui ne parvient pas encore à nourrir sa population avec sa propre production.
\end{savaistu}

\courssub{La détérioration des termes de l'échange}

\textbf{Intuition.} En 1980, un planteur de Ngoumou vendait un sac de cacao et pouvait acheter, avec la recette, une bicyclette neuve. Aujourd'hui, avec le même sac, il ne paie même plus la moitié d'une moto d'occasion. Le cacao n'a pas disparu, mais son \emph{pouvoir d'achat} sur les produits manufacturés a fondu. C'est cela, la détérioration des termes de l'échange.

\begin{definition}[Détérioration des termes de l'échange]
Les \cle{termes de l'échange} mesurent le rapport entre le niveau des prix des exportations et celui des prix des importations :
\[ \text{Termes de l'échange} = \frac{\text{indice des prix à l'exportation}}{\text{indice des prix à l'importation}} \times 100 \]
On parle de \cle{détérioration des termes de l'échange} lorsque ce rapport diminue : les prix des produits exportés (matières premières) baissent ou stagnent, tandis que les prix des produits importés (biens manufacturés) augmentent. Le pays doit alors exporter \emph{de plus en plus} pour importer \emph{la même quantité}.
\end{definition}

\begin{exemplebox}[Le sac de cacao contre le tracteur]
Supposons qu'en l'an 2000, un tracteur coûte 10 millions de FCFA et un sac de cacao de 65 kg vaut 50 000 FCFA : il faut 200 sacs pour acheter un tracteur. En 2020, le tracteur coûte 25 millions de FCFA (hausse des prix manufacturés) et le sac ne vaut plus que 62 500 FCFA : il faut désormais 400 sacs. En vingt ans, l'effort d'exportation nécessaire a \textbf{doublé}, alors que le planteur n'a rien changé à son travail.
\end{exemplebox}

\begin{popfigure}
\begin{center}
\begin{tikzpicture}
\begin{axis}[
width=0.92\linewidth, height=7cm,
xlabel={Années}, ylabel={Indice des prix (base 100 en 2000)},
xmin=2000, xmax=2024, ymin=60, ymax=220,
xtick={2000,2006,2012,2018,2024},
ytick={60,100,140,180,220},
legend style={at={(0.02,0.98)}, anchor=north west, font=\small},
grid=major, grid style={popline!40},
]
\addplot[popcurve, color=popgreen, very thick, smooth] coordinates {
(2000,100) (2004,92) (2008,105) (2012,88) (2016,95) (2020,85) (2024,110)
};
\addplot[popcurve, color=poporange, very thick, smooth] coordinates {
(2000,100) (2004,112) (2008,128) (2012,145) (2016,162) (2020,178) (2024,195)
};
\legend{Cacao exporté (indice), Biens manufacturés importés (indice)}
\end{axis}
\end{tikzpicture}
\end{center}
\legende{Évolution comparée (schématique, ordres de grandeur) du prix du cacao exporté et du prix des biens manufacturés importés, base 100 en 2000. L'écart qui se creuse entre les deux courbes matérialise la détérioration des termes de l'échange : le Cameroun doit vendre toujours plus de cacao pour acheter la même machine.}
\end{popfigure}

\begin{methode}[Analyser un graphique de termes de l'échange]
\begin{enumerate}
\item \textbf{Lire les axes} : identifier les grandeurs (indices de prix, années) et l'unité (base 100 à quelle date ?).
\item \textbf{Décrire les tendances générales} : chaque courbe monte-t-elle, descend-elle, stagne-t-elle sur l'ensemble de la période ?
\item \textbf{Repérer les croisements} : le moment où la courbe des prix importés dépasse durablement celle des prix exportés marque le basculement vers la détérioration.
\item \textbf{Repérer les ruptures} : chutes brutales (krach des cours, crise de 2008, chute pétrolière de 2014-2015) et rebonds ponctuels.
\item \textbf{Calculer l'écart} : comparer les indices à la fin de la période (ici environ 110 contre 195) et conclure sur l'effort d'exportation supplémentaire exigé.
\item \textbf{Conclure avec prudence} : relier le graphique aux conséquences concrètes (déficit, endettement, appauvrissement des producteurs).
\end{enumerate}
\end{methode}

\courssub{La dépendance à quelques produits d'exportation}

Le Cameroun tire l'essentiel de ses devises d'un petit nombre de produits : le \cle{pétrole} (première source de devises, environ 40 à 50 \% des recettes d'exportation selon les années, ordre de grandeur), le \cle{cacao}, le \cle{bois}, puis le coton, l'aluminium, le café, la banane. Or les cours de ces matières premières sont fixés sur les marchés mondiaux (Londres, New York) et fluctuent fortement, sans que le Cameroun puisse les influencer.

\begin{exoresolu}[Effet d'une chute des cours du cacao sur les recettes d'exportation]
\textbf{Énoncé.} Le Cameroun exporte 250 000 tonnes de cacao par an. Le cours mondial passe de 2 400 à 1 500 dollars la tonne (chute de 37,5 \%, comme cela s'est déjà produit). Calcule la perte de recettes d'exportation en dollars, puis commente.
\tcblower
\textbf{Solution.}
\begin{align*}
\text{Recettes avant} &= 250\,000 \times 2\,400 = 600\ \text{millions de dollars}\\
\text{Recettes après} &= 250\,000 \times 1\,500 = 375\ \text{millions de dollars}\\
\text{Perte} &= 600 - 375 = 225\ \text{millions de dollars}
\end{align*}
Soit une perte de 225 millions de dollars (environ 135 milliards de FCFA), sans que le pays ait exporté un gramme de moins. Cette perte se répercute sur le budget de l'État (taxes à l'exportation), sur le revenu des planteurs et sur la balance commerciale. C'est toute la \cle{vulnérabilité} d'une économie dépendante des cours mondiaux.
\end{exoresolu}

\begin{popfigure}
\begin{center}
\begin{tikzpicture}[node distance=1.1cm, every node/.style={font=\small}]
\node[draw=popdark, fill=poporangeL, rounded corners, align=center, text width=4.2cm] (a) {Dépendance à quelques matières premières brutes};
\node[draw=popdark, fill=poppinkL, rounded corners, align=center, text width=4.2cm, right=3.2cm of a] (b) {Vulnérabilité aux fluctuations des cours mondiaux};
\node[draw=popdark, fill=poppurpleL, rounded corners, align=center, text width=4.2cm, below=1.4cm of b] (c) {Recettes instables : déficit commercial et budgétaire};
\node[draw=popdark, fill=popblueL, rounded corners, align=center, text width=4.2cm, left=3.2cm of c] (d) {Peu de moyens pour industrialiser et transformer};
\draw[-{Stealth[length=3mm]}, very thick, popdark] (a) -- (b);
\draw[-{Stealth[length=3mm]}, very thick, popdark] (b) -- (c);
\draw[-{Stealth[length=3mm]}, very thick, popdark] (c) -- (d);
\draw[-{Stealth[length=3mm]}, very thick, popdark] (d) -- (a);
\node[font=\bfseries\small, popdark] at ($(a)!0.5!(c)$) {Cercle vicieux};
\end{tikzpicture}
\end{center}
\legende{Le cercle vicieux de la dépendance aux matières premières : 1 dépendance à quelques produits bruts ; 2 exposition aux chocs des cours ; 3 instabilité des recettes ; 4 sous-investissement dans l'industrie de transformation, ce qui reconduit la dépendance.}
\end{popfigure}

\courssub{La faible transformation locale des matières premières}

Le Cameroun exporte l'essentiel de ses produits \emph{bruts} : grumes de bois plutôt que meubles, fèves de cacao plutôt que chocolat, coton-fibre plutôt que tissus, bauxite importée transformée en aluminium mais sans filière locale complète. Or c'est la \cle{transformation} qui crée la valeur ajoutée, les emplois industriels et les devises.

\textbf{Pourquoi si peu d'industries de transformation ?}
\begin{itemize}
\item Le \textbf{coût et l'irrégularité de l'énergie} : les délestages électriques pénalisent les usines (une chocolaterie ou une scierie moderne exige un courant stable).
\item Le \textbf{manque de capitaux} et de financement à long terme pour les entrepreneurs locaux.
\item La \textbf{concurrence} des produits manufacturés importés, souvent moins chers.
\item Des \textbf{compétences techniques} encore insuffisantes et un marché intérieur étroit.
\end{itemize}

\begin{exemplebox}[Le paradoxe du chocolat]
Le Cameroun figure parmi les premiers producteurs mondiaux de cacao (environ 250 000 à 300 000 tonnes par an, ordre de grandeur), mais l'essentiel des fèves part brut vers l'Europe et l'Asie. Une tablette de chocolat vendue à Douala peut ainsi avoir été fabriquée en Suisse ou en Belgique... avec du cacao camerounais. La valeur ajoutée (torréfaction, broyage, conchage, emballage, marque) est captée à l'étranger : le prix final du chocolat peut atteindre dix fois celui de la fève.
\end{exemplebox}

\courssub{Les insuffisances des infrastructures et le poids de l'informel}

\textbf{Des infrastructures insuffisantes et coûteuses.}
\begin{itemize}
\item La \cle{congestion du port de Douala}, par lequel transite l'essentiel du commerce extérieur du pays (et celui du Tchad et de la Centrafrique) : faible tirant d'eau (environ 7-8 m) qui limite les grands navires, lenteurs des opérations et des formalités, coûts de passage élevés. Le port de Kribi, en eau profonde (16-18 m), opérationnel depuis 2014, commence à soulager ce goulot d'étranglement.
\item Des \textbf{routes dégradées} : l'axe Douala--Yaoundé, les pistes des zones de production (cacao du Centre, coton du Nord) augmentent les délais et cassent les marchandises.
\item Des \textbf{coûts de transport élevés} qui renchérissent les exportations et réduisent leur compétitivité : exporter un conteneur depuis l'intérieur du Cameroun coûte nettement plus cher que depuis un pays doté d'infrastructures modernes.
\end{itemize}

\textbf{L'informel et la fraude aux frontières.} Une part importante des échanges échappe aux statistiques et aux douanes : contrebande de carburant vers le Nigeria (profitant de la différence de prix subventionné), importations frauduleuses de riz et de produits manufacturés par les corridors de l'Adamaoua, sous-facturation, passage par les pistes non contrôlées. Conséquences : pertes de recettes douanières pour l'État, concurrence déloyale pour les producteurs et commerçants réguliers, et fausse mesure des échanges réels (la balance commerciale officielle sous-estime les flux).

\begin{attention}[Un excédent ponctuel ne règle pas les problèmes structurels]
Lorsque le cours du pétrole flambe (comme en 2008 ou en 2022), la balance commerciale du Cameroun peut devenir temporairement excédentaire ou son déficit se réduire fortement. \textbf{Piège de rédaction au Bac} : ne conclus jamais que « les problèmes sont réglés ». Il s'agit d'un effet de conjoncture (les prix), non d'une transformation des structures : le pays exporte toujours des produits bruts, importe toujours ses biens d'équipement et une partie de sa nourriture, et reste vulnérable dès que les cours retombent. Distingue toujours \emph{conjoncture favorable} et \emph{santé structurelle} de l'économie.
\end{attention}

\courssub{Les pistes de solutions}

Face à ces problèmes, plusieurs voies sont explorées par l'État et les partenaires :

\begin{enumerate}
\item \textbf{Industrialiser et transformer localement} : développer les unités de transformation du cacao (broyage, pâte, beurre, poudre), du bois (meubles, contreplaqué, placages), du coton (filature, teinture, confection), les cimenteries, les raffineries (projet de raffinerie de Kribi). Objectif : exporter des produits à plus forte valeur ajoutée.
\item \textbf{Diversifier les exportations} : ne plus dépendre du trio pétrole--cacao--bois en promouvant de nouvelles filières (agropastoral : ananas, banane, piment ; minerais : fer de Mbalam/Nabeba, cobalt, nickel ; produits transformés ; services ; numérique).
\item \textbf{Réduire la facture alimentaire} : relancer la production de riz (périmètres irrigués de Yagoua, Nanga-Eboko, Ndop), du maïs, de la pêche et de l'aquaculture pour substituer la production locale aux importations.
\item \textbf{Moderniser les infrastructures} : extension du port en eau profonde de Kribi (phase 2), réhabilitation des axes routiers (Douala-Yaoundé, Yaoundé-Ngaoundéré) et ferroviaires (Transcamerounais), amélioration de l'approvisionnement en énergie (barrages de Nachtigal 420 MW, Memve'ele 200 MW, Song Loulou).
\item \textbf{Sécuriser les frontières et formaliser les échanges} : modernisation des douanes (système CAMCIS), corridors officiels, lutte contre la fraude, tout en facilitant le petit commerce frontalier qui fait vivre des millions de personnes (voir Dossier 5).
\item \textbf{Approfondir l'intégration régionale} : la \cle{ZLECAF} (Zone de libre-échange continentale africaine), entrée en vigueur en 2021 (démarrage des échanges préférentiels en 2023), offre un marché de plus d'un milliard de consommateurs ; elle peut stimuler les exportations camerounaises vers l'Afrique, à condition que les entreprises locales soient compétitives.
\end{enumerate}

\begin{aretenir}[L'essentiel de la section]
Les échanges extérieurs du Cameroun souffrent de cinq problèmes majeurs : un \cle{déficit chronique de la balance commerciale} (importations massives de biens d'équipement et de produits alimentaires), une \cle{détérioration des termes de l'échange} (les matières premières exportées achètent de moins en moins de produits manufacturés), une \cle{dépendance} à quelques produits bruts vulnérables aux cours mondiaux, une \cle{faible transformation locale} et des \cle{infrastructures insuffisantes} aggravées par l'informel et la fraude. Les solutions passent par l'industrialisation, la diversification, la souveraineté alimentaire, la modernisation des infrastructures et l'intégration régionale (CEMAC, ZLECAF).
\end{aretenir}

\begin{exercice}[Pour t'entraîner]
\begin{enumerate}
\item Définis les termes de l'échange et explique, avec un exemple chiffré de ton choix, ce qu'est leur détérioration.
\item Le Cameroun exporte 300 000 tonnes de cacao à 2 000 dollars la tonne. Le cours chute de 30 \%. Calcule la perte de recettes et propose deux mesures pour limiter ce type de risque.
\item Construis un croquis simple montrant comment la congestion du port de Douala renchérit le coût des exportations camerounaises (chaîne : producteur $\rightarrow$ route $\rightarrow$ port $\rightarrow$ navire).
\end{enumerate}
\end{exercice}

\begin{corrige}[Exercice]
\textbf{1.} Les termes de l'échange sont le rapport (prix des exportations / prix des importations) $\times$ 100. Ils se détériorent quand les prix à l'exportation progressent moins vite (ou baissent) que les prix à l'importation. Exemple : si 100 sacs de cacao achetaient une moto en 2010 et qu'il en faut 180 en 2024, les termes de l'échange se sont détériorés pour le planteur et pour le pays.

\textbf{2.} Recettes avant : $300\,000 \times 2\,000 = 600$ millions de dollars. Après une chute de 30 \%, le cours vaut $2\,000 \times 0{,}70 = 1\,400$ dollars ; recettes : $300\,000 \times 1\,400 = 420$ millions. Perte : $600 - 420 = 180$ millions de dollars. Mesures : transformer le cacao sur place (vendre de la pâte ou du chocolat, moins sensibles aux cours bruts) et diversifier les exportations pour ne pas dépendre d'un seul produit.

\textbf{3.} Le croquis doit montrer la chaîne producteur $\rightarrow$ route dégradée (coût +) $\rightarrow$ port congestionné (attente, frais +) $\rightarrow$ navire, avec à chaque étape un surcoût qui renchérit le prix final du produit exporté et réduit sa compétitivité. Légende et flèches obligatoires.
\end{corrige}

\coursec{Cadre général et notions clés des échanges extérieurs}

Le Cameroun, par sa position charnière entre l'Afrique de l'Ouest et l'Afrique centrale, sa façade maritime sur le golfe de Guinée (400\,km de côtes) et ses 4\,900\,km de frontières terrestres, est une \cle{plateforme d'échanges} naturelle. Pourtant, sa balance commerciale structurellement déficitaire révèle une intégration inaboutie à la mondialisation. Comprendre ses échanges extérieurs, c'est analyser la tension entre un potentiel géographique réel et une structure économique héritée.

\begin{definition}[Notions fondamentales du commerce extérieur]
\begin{itemize}
    \item \textbf{Exportations (X) :} Biens et services vendus à l'extérieur (recettes en devises). \cle{Produits primaires} (pétrole, bois, cacao, coton, aluminium) dominent.
    \item \textbf{Importations (M) :} Biens et services achetés à l'extérieur (dépenses en devises). \cle{Produits manufacturés} (machines, véhicules, produits chimiques, céréales, carburant raffiné) dominent.
    \item \textbf{Balance commerciale (BC) :} $BC = X - M$. Excédent si $X > M$, déficit si $M > X$. Le Cameroun enregistre un \textbf{déficit chronique} (sauf années de pic pétrolier).
    \item \textbf{Taux de couverture :} $\frac{X}{M} \times 100$. S'il est $< 100\%$, le pays n'autofinance pas ses importations.
    \item \textbf{Termes de l'échange (TE) :} Rapport entre l'indice des prix à l'exportation et l'indice des prix à l'importation ($TE = \frac{P_X}{P_M} \times 100$). Une \textbf{détérioration} ($TE \downarrow$) signifie qu'il faut exporter plus pour importer la même quantité.
    \item \textbf{Partenaires :} Pays ou zones avec lesquels les flux sont intenses. On distingue \textbf{partenaires traditionnels} (UE, France) et \textbf{nouveaux partenaires} (Chine, Inde, pays émergents, Afrique).
\end{itemize}
\end{definition}

\begin{popfigure}
\includegraphics[width=\linewidth,height=10cm,keepaspectratio]{N16/cmr_map.jpg}
\legende{La façade maritime et les frontières terrestres du Cameroun. Au sud-ouest, le golfe de Guinée où s'ouvrent les ports de Douala (95\,\% du trafic) et de Kribi (minerais, vrac) ; au nord, à l'est et au sud, les frontières avec le Nigeria (CEDEAO), le Tchad, la République centrafricaine, le Congo, le Gabon et la Guinée équatoriale (CEMAC). Les routes et la voie ferrée qui mènent aux frontières apparaissent sur la carte. Les produits exportés (pétrole, bois, cacao, coton) et importés (machines, carburant, riz, blé, véhicules) sont ceux du texte ; la carte ne représente pas les flux. \\ Source : Wikimedia Commons, CIA, domaine public.}
\end{popfigure}

\begin{methode}[Calculer et interpréter le Taux de Couverture et les Termes de l'Échange]
\begin{enumerate}
    \item \textbf{Identifier les valeurs :} $X$ (valeur exportations FOB), $M$ (valeur importations CIF) sur une année (source : INS, BEAC, Douanes).
    \item \textbf{Taux de couverture :} $TC = \frac{X}{M} \times 100$. Ex: $X=3\,500$ Mds FCFA, $M=4\,200$ Mds $\to TC \approx 83\%$. \textbf{Interprétation :} Le pays couvre 83\% de ses importations par ses exportations ; le gap (17\%) est financé par emprunts, IDE, transferts diaspora.
    \item \textbf{Termes de l'échange (année de base 100) :} $TE_t = \frac{I_{P_X,t}}{I_{P_M,t}} \times 100$. Si $TE$ passe de 110 à 95 : \textbf{détérioration de 13,6\%}. Le pouvoir d'achat des exportations baisse.
    \item \textbf{Analyse structurelle :} Une détérioration durable des TE + TC $< 100\%$ = \textbf{piège de la spécialisation primaire} (volatilité prix matières premières vs rigidité prix manufacturés).
\end{enumerate}
\end{methode}

\begin{attention}[Pièges de rédaction au Bac]
\begin{itemize}
    \item Ne pas confondre \textbf{Balance commerciale} (biens seulement) et \textbf{Balance courante} (biens + services + revenus + transferts). Le déficit commercial peut être compensé par les transferts (diaspora) ou services (transit).
    \item \textbf{FOB vs CIF :} Exportations valorisées \emph{Free On Board} (à la sortie du port), Importations \emph{Cost Insurance Freight} (à l'arrivée). Cela surévalue mécaniquement le déficit commercial de $\sim 10-15\%$ (coûts fret/assurance payés à des non-résidents).
    \item « Le Cameroun exporte du pétrole » $\to$ Préciser : \textbf{Pétrole brut} (export) vs \textbf{Produits pétroliers raffinés} (import, car la raffinerie SONARA ne couvre pas la demande locale).
\end{itemize}
\end{attention}

\coursec{Les produits échangés : structure et évolution}

La structure des échanges reflète l'absence de transformation locale : on exporte le brut, on importe le transformé.

\courssub{Les exportations : la rente primaire (80-85\% du total)}

\begin{center}
\begin{tabularx}{\linewidth}{|l|c|X|}\hline
\rowcolor{popblueL}
\textbf{Produit} & \textbf{Part \textasciitilde 2022-2023} & \textbf{Caractéristiques et enjeux} \\ \hline
\textbf{Pétrole brut} & \textasciitilde 35-40\% & Exploité par \textbf{SNH} (partenaire : Perenco, TotalEnergies, BW Energy). Production en déclin (\textasciitilde 30-35 millions de barils/an vs 80 M dans les années 80). Prix volatile (Brent). Recettes budgétaires très dépendantes. \\ \hline
\textbf{Bois et articles en bois} & \textasciitilde 15-18\% & Grumes (interdiction export depuis 2017, mais dérogations), sciages, placages, contreplaqué. Principaux clients : \textbf{Chine} (60\%+), UE (VPA-FLEGT), Vietnam. Enjeu : transformation locale (usines de Kribi, Bertoua) vs export brut. \\ \hline
\textbf{Cacao (fèves)} & \textasciitilde 10-12\% & 4\textsuperscript{e}-5\textsuperscript{e} producteur mondial (\textasciitilde 280-300 kt/an). \textbf{90\% exporté en fèves brutes}. Clients : Pays-Bas (port Amsterdam), Belgique, Allemagne, Malaisie. Enjeu : taux de transformation local (\textasciitilde 30\% objectif 50\%), qualité (normes UE déforestation). \\ \hline
\textbf{Coton (fibre)} & \textasciitilde 4-6\% & Géré par \textbf{SODECOTON}. Zone Nord (Adamaoua, Nord, Extrême-Nord). Export vers Asie (Bangladesh, Vietnam, Chine) et UE. Subventions US (Cotton-4) faussent prix mondial. \\ \hline
\textbf{Aluminium (alumine/lingots)} & \textasciitilde 5-7\% & \textbf{ALUCAM} (Édéa). Importe alumine (Guinée, Australie), exporte lingots. Énergivore (barrage de Song Loulou, Nachtigal). \\ \hline
\textbf{Autres (Banane, Café, Caoutchouc, Fruits)} & \textasciitilde 5-8\% & Banane \textbf{Plantations du Haut Penja (PHP)}, \textbf{Boh Plantations} (certifiée Fairtrade). Café robusta (Ouest, Littoral) et arabica (Nord-Ouest). \\ \hline
\end{tabularx}
\end{center}

\begin{exemplebox}[Le piège de la « Dutch Disease » (Maladie hollandaise)]
La rente pétrolière (années 80) et aujourd'hui la rente bois/cacao entraînent une appréciation réelle du taux de change (FCFA arrimé à l'Euro fort). Cela rend les autres secteurs exportateurs (agro-industrie naissante, textile, services) \textbf{moins compétitifs} à l'international. Le Cameroun peine à diversifier son panier d'exportation (Indice de concentration HHI élevé).
\end{exemplebox}

\courssub{Les importations : la facture de la non-transformation}

\begin{center}
\begin{tabularx}{\linewidth}{|l|c|X|}\hline
\rowcolor{poporangeL}
\textbf{Groupe de produits} & \textbf{Part \textasciitilde 2022-2023} & \textbf{Détails et fournisseurs clés} \\ \hline
\textbf{Produits pétroliers raffinés} & \textasciitilde 15-20\% & \textbf{Poste n°1 ou 2}. SONARA (Limbe) raffine \textasciitilde 2,1 Mt/an (capacité 2,5 Mt) mais arrêts fréquents. Import massif (Super, Gasoil, Kérosène) depuis Europe, Moyen-Orient, Nigéria (informel). Subvention étatique lourde (\textasciitilde 1\,000 Mds FCFA/an récemment). \\ \hline
\textbf{Machines, matériel électrique, transport} & \textasciitilde 20-25\% & Équipements chantiers (barrages, routes, ports), groupes électrogènes, véhicules d'occasion (Europe, Japon, USA), téléphones. Fournisseurs : \textbf{Chine} (n°1), France, Inde, Japon. \\ \hline
\textbf{Céréales (Blé, Riz, Maïs, Froment)} & \textasciitilde 10-12\% & \textbf{Blé} (\textasciitilde 900 kt/an) : Russie, France, Ukraine (usines \textbf{SMBC}, \textbf{UCF} à Douala). \textbf{Riz} (\textasciitilde 600 kt/an) : Asie (Thaïlande, Inde, Vietnam, Birmanie), fraude massive via Nigéria. Objectif autosuffisance (PDAR) non atteint. \\ \hline
\textbf{Produits chimiques / Pharmaceutiques} & \textasciitilde 10\% & Engrais (NPK, Urée - importés alors que gaz local existe), médicaments (Inde, Chine, Europe). \\ \hline
\textbf{Autres (Poissons surgelés, Véhicules, Textiles, Boissons)} & \textasciitilde 20\% & Poisson (Chine, Maroc, Mauritanie) pour combler déficit halieutique. Véhicules d'occasion (\"Belgique\", \"France\"). \\ \hline
\end{tabularx}
\end{center}

\begin{aretenir}[Structure type Bac]
\begin{itemize}
    \item \textbf{Export :} 5 produits = 80\% (Pétrole, Bois, Cacao, Coton, Alu). Faible transformation (taux \textasciitilde 15-20\%). Concentration géographique (Chine, UE).
    \item \textbf{Import :} Produits manufacturés + Énergie + Alimentation. Dépendance structurelle.
    \item \textbf{Conséquence :} Vulnérabilité aux chocs externes (prix matières premières, taux de change €/\$, chaînes logistiques).
\end{itemize}
\end{aretenir}

\coursec{Les échanges avec les partenaires africains et la CEMAC}

L'intégration régionale est une priorité affichée (Agenda 2063 UA, ZLECAf), mais la réalité des flux reste faible comparée aux potentialités.

\courssub{La CEMAC : un marché commun inachevé}

\begin{definition}[CEMAC (Communauté Économique et Monétaire de l'Afrique Centrale)]
6 États : \textbf{Cameroun, Tchad, RCA, Gabon, Guinée équatoriale, Congo}. 60 millions d'habitants. Monnaie unique : \textbf{FCFA (XAF)} émis par \textbf{BEAC} (Siège : Yaoundé). Parité fixe : 1 € = 655,957 FCFA. Objectifs : Union douanière (TEC - Tarif Extérieur Commun), libre circulation personnes/biens/capitaaux, marché commun.
\end{definition}

\begin{popfigure}
\includegraphics[width=\linewidth,height=9cm,keepaspectratio]{N16/cemac.png}
\legende{Le marché commun de la CEMAC (six États en vert : Cameroun, Tchad, République centrafricaine, Gabon, Congo, Guinée équatoriale) au sein de l'Afrique. Au centre de cet espace, le Cameroun est le \textbf{fournisseur net} de biens manufacturés (ciment, boissons, savon, farine) et le \textbf{client net} de matières premières (bétail, bois, coton, pétrole brut) ; la monnaie unique supprime le risque de change. La carte localise les membres de la zone mais ne représente pas les flux intra-CEMAC (ceux-ci sont décrits dans le texte). \\ Source : Wikimedia Commons, Sven-Steffen Arndt, AnzeeTM et Hoshie, domaine public.}
\end{popfigure}

\begin{exemplebox}[Le Cameroun, « usine » et « marché » de la CEMAC]
\begin{itemize}
    \item \textbf{Export Cameroun $\to$ CEMAC :} Ciment (\textbf{CIMENCAM}, \textbf{Dangote} - \textasciitilde 2,5 Mt/an), Boissons (\textbf{Brasseries du Cameroun}, \textbf{Guinness}), Savon/Détergents (\textbf{Socaver}, \textbf{Azur}), Farine (\textbf{SMBC}), Médicaments (\textbf{Cinpharm}). Parts de marché souvent $> 70\%$ dans la sous-région.
    \item \textbf{Import CEMAC $\to$ Cameroun :} Bétail sur pied (Tchad, RCA - approvisionne Douala/Yaoundé), Bois grumes/transformés (Gabon, Congo, RCA), Coton (Tchad), Pétrole brut (Tchad, Congo - pour SONARA).
\end{itemize}
\textbf{Chiffre clé :} Le commerce intra-CEMAC représente seulement \textbf{2 à 3\%} du commerce total du Cameroun (vs 15-20\% dans l'UE ou l'ASEAN). Potentiel inexploité énorme.
\end{exemplebox}

\courssub{La CEDEAO et le Nigéria : le géant voisin}

Le Cameroun n'est pas membre de la CEDEAO mais observateur. La frontière avec le \textbf{Nigéria} (200 millions d'habitants, 1\textsuperscript{re} économie Afrique) structure l'Extrême-Nord et le Sud-Ouest.

\begin{propriete}[Asymétrie structurelle Cameroun-Nigéria]
\begin{itemize}
    \item \textbf{Monnaie :} FCFA (stable) vs Naira (NGN, flottant, très volatil). Dévaluation Naira 2023-24 : 1\,000 FCFA $\approx$ 1\,500 NGN (vs 450 NGN en 2020). $\to$ Produits camerounais chers pour Nigérians ; Produits nigérians (carburant, riz, tissus) bon marché pour Camerounais.
    \item \textbf{Flux formels :} Faibles (barrières non tarifaires, paperasse, méfiance).
    \item \textbf{Flux informels :} Massifs (voir Dossier 5). Carburant (Nig$\to$Cmr), Riz (Nig$\to$Cmr), Bétail/Oignons (Cmr$\to$Nig).
    \item \textbf{ZLECAf :} Ratification par le Cameroun (2021). Entrée en vigueur progressive. Doit réduire droits de douane à 0\% sur 90\% des lignes tarifaires sur 10-13 ans. Enjeu : compétitivité entreprises camerounaises face à la concurrence nigériane/ghanéenne.
\end{itemize}
\end{propriete}

\coursec{Les échanges avec le reste du monde : l'UE et la Chine, deux pôles}

\courssub{L'Union Européenne : partenaire historique en mutation}

\begin{center}
\begin{tabularx}{\linewidth}{|l|X|X|}\hline
\rowcolor{popblueL}
\textbf{Dimension} & \textbf{Export Cameroun $\to$ UE} & \textbf{Import UE $\to$ Cameroun} \\ \hline
\textbf{Produits} & Bois (certifié FLEGT), Cacao, Banane, Aluminium, Pétrole, Caoutchouc. & Machines, Véhicules, Pharmaceutique, Produits chimiques, Vins/Alcools, Blé. \\ \hline
\textbf{Accords} & \textbf{APE (Accord de Partenariat Économique)} Cameroun-UE (intérimaire 2009, ratifié 2014). Démantèlement droits de douane sur 15-25 ans. \textbf{VPA-FLEGT} (Bois légal). & Accès préférentiel marché camerounais. \\ \hline
\textbf{Part de marché} & \textasciitilde 40-45\% des exportations (hors pétrole brut souvent vendu hors UE). & \textasciitilde 30-35\% des importations (en baisse face à la Chine). \\ \hline
\textbf{Enjeux actuels} & \textbf{Règlement UE Déforestation (EUDR)} 2024 : traçabilité géolocalisée obligatoire pour cacao, bois, café, caoutchouc, huile de palme. Risque d'exclusion petits producteurs. Normes sanitaires (LMR pesticides cacao). & Transfert technologie, investissements (EIB, AFD), aide au développement. \\ \hline
\end{tabularx}
\end{center}

\courssub{La Chine : premier partenaire commercial (volume)}

\begin{exemplebox}[La relation sino-camerounaise : infrastructures contre ressources]
\begin{itemize}
    \item \textbf{Import depuis Chine (\textasciitilde 20-25\% parts) :} Machines TP, textiles, chaussures, électronique, acier, motos, téléphones (Transsion/Tecno/Infinix/Itel dominent marché). Prix bas, adaptation locale.
    \item \textbf{Export vers Chine (\textasciitilde 15-20\% parts) :} \textbf{Bois (grumes, sciages - 70\% export bois)}, Coton, Pétrole brut, Minerais (Cobalt, Fer - projets Mbalam/Nabeba).
    \item \textbf{Investissements (IDE/BRI) :} Port de Kribi (CHEC), Barrages (Memve'ele - Sinohydro, Nachtigal - CGCOC), Routes, Stade Olembé, Assemblée nationale, Hôpital de Douala (Gynéco-Obstétrique). Financement : Prêts concessionnels / Ressources naturelles (pétrole, minerais) / Budget étatique.
    \item \textbf{Dette :} Chine = premier créancier bilatéral (\textasciitilde 30-35\% dette extérieure publique). Risque de surendettement (FMI alerte).
\end{itemize}
\end{exemplebox}

\courssub{Autres partenaires émergents}
\begin{itemize}
    \item \textbf{Inde :} Pharmaceutique (génériques), véhicules (Tata, Mahindra), riz, coton (achat).
    \item \textbf{Turquie :} Construction (Yenigün, Summa), textiles, engrais, écoles/universités.
    \item \textbf{Corée du Sud / Japon :} Navires, électronique, voitures, projets infrastructurels (JICA, KOICA).
    \item \textbf{Russie / Ukraine (avant 2022) :} Blé, engrais. Guerre = rupture approvisionnement, hausse prix.
\end{itemize}

\begin{aretenir}[Géographie des partenaires - Synthèse Bac]
\begin{enumerate}
    \item \textbf{UE :} Partenaire historique, normes strictes (qualité, durabilité), APE/VPA-FLEGT. Perte de vitesse relative.
    \item \textbf{Chine :} Partenaire n°1 volume, infrastructures lourdes, débouché bois/minerais, dette, produits bon marché.
    \item \textbf{CEMAC :} Marché naturel (monnaie unique), complémentarités fortes, mais freins physiques (routes) et institutionnels (barrières non tarifaires).
    \item \textbf{Nigéria/CEDEAO :} Géographie imposée, informel dominant, enjeu ZLECAf.
    \item \textbf{Diversification :} Inde, Turquie, Golfe (Qatar, EAU - investissements portuaires/hôteliers).
\end{enumerate}
\end{aretenir}

\coursec{Les problèmes des échanges extérieurs du Cameroun}

Au-delà des chiffres, des blocages structurels expliquent la sous-performance.

\courssub{La détérioration structurelle des termes de l'échange (DTE)}

\begin{propriete}[Mécanisme de la DTE (Prébisch-Singer appliqué au Cameroun)]
\begin{itemize}
    \item \textbf{Prix export (Matières premières) :} Volatils, tendance baissière réelle long terme (progrès techniques économie matière, substituts synthétiques, ralentissement Chine).
    \item \textbf{Prix import (Manufacturés) :} Tendance haussière (innovation, valeur ajoutée, propriété intellectuelle, marques).
    \item \textbf{Résultat :} Pour acheter 1 tracteur (import), il faut vendre \textbf{plus de tonnes de cacao/bois/coton} aujourd'hui qu'hier.
    \item \textbf{Preuve :} Indice TE (base 100 en 2000) souvent $< 80$ depuis 2015 (chute prix pétrole, cacao, coton). Pic 2022 (prix matières 1ères guerre Ukraine) mais retombée 2023-24.
\end{itemize}
\end{propriete}

\begin{experience}[TP : Construire une courbe des Termes de l'Échange (Données BEAC/INS)]
\textbf{Matériel :} Série chronologique 2010-2023 : Indice Prix Export ($I_{PX}$), Indice Prix Import ($I_{PM}$).
\textbf{Protocole :}
\begin{enumerate}
    \item Calculer $TE_t = \frac{I_{PX,t}}{I_{PM,t}} \times 100$ pour chaque année.
    \item Tracer la courbe (pgfplots : axe X = années, axe Y = TE).
    \item Identifier phases : Hausse (2010-2014), Chute (2015-2016), Reprise relative (2017-2019), Chute Covid (2020), Pic 2022, Repli 2023.
    \item Corréler avec conjoncture (prix pétrole, cacao, taux de change €/\$).
\end{enumerate}
\textbf{Interprétation attendue :} Volatilité extrême = instabilité recettes budgétaires = difficulté planification investissements publics.
\end{experience}

\courssub{La dépendance à la rente extractive et l'absence de transformation}

\begin{attention}[Erreur fréquente : « Le Cameroun n'a pas d'industrie »]
\textbf{Faux.} Il y a une industrie (Alucam, Sonara, Cimencam, Brasseries, Savonneries, PHP, Usines bois, Montage motos/vélos, Pharmaceutique).
\textbf{Juste :} C'est une \textbf{industrie de première transformation ou d'assemblage}, faiblement intégrée en amont (intrants importés) et en aval (peu de biens d'équipement produits). Taux de transformation locale \textasciitilde 15-20\% (Objectif SND30 : 50\%).
\end{attention}

\courssub{Les coûts de transaction excessifs (Logistique et Gouvernance)}

\begin{center}
\begin{tabularx}{\linewidth}{|l|X|}\hline
\rowcolor{popredL}
\textbf{Facteur} & \textbf{Impact sur la compétitivité} \\ \hline
\textbf{Port de Douala} & Encombrement (délai moyen 15-20 jours vs 3-5 jours ports performants), coûts manutention élevés, vétusté équipements (malgré concession Douala Port Authority / Bolloré). Kribi (KPA) désengorge mais accès routier/ferroviaire perfectible. \\ \hline
\textbf{Corridor Douala-Ndjamena/Bangui} & Route unique (N1/N3), dégradée, barrages multiples (rackets), insécurité (coupeurs de route Est). Coût transport km/tonne parmi les plus élevés Afrique. \\ \hline
\textbf{Procédures douanières} & Paperasse, délais dédouanement (même avec ASYCUDA World), corruption perçue (Indice Transparency International). \\ \hline
\textbf{Énergie} & Coupures fréquentes (ENEO), coût électricité industriel élevé vs concurrence (Maroc, Égypte, Vietnam). \\ \hline
\end{tabularx}
\end{center}

\courssub{L'insécurité : un frein majeur aux flux}

\begin{itemize}
    \item \textbf{Extrême-Nord (Boko Haram/ISWAP) :} Fermeture frontière nigériane 2015-2017, perturbation axes Maroua-Kousseri, destruction marchés (Banki, Fotokol), déplacement populations, chute production oignon/bétail.
    \item \textbf{Nord-Ouest / Sud-Ouest (Crise anglophone) :} Villes fantômes, blocages routes (\"Ghost towns\"), destruction plantations (CDC, Pamol, PHP), insécurité portuaire (Ekondo-Titi, Idenau), fuite investisseurs.
    \item \textbf{Est (RCA) :} Instabilité chronique, groupes armés, coupeurs de route sur axe Garoua-Boulaï-Bertoua, risque sur évacuation bétail/bois.
\end{itemize}

\courssub{Le déficit de données : l'informel invisible}

Comme détaillé dans le Dossier 5, \textbf{50 à 80\% du commerce intra-africain} est informel. Les politiques publiques (SND30, PNDP) sont élaborées sur des statistiques douanières partielles $\to$ mauvais ciblage des appuis (ex: subventions engrais détournées vers Nigéria via informel).

\begin{aretenir}[Problèmes - Check-list Bac]
\begin{enumerate}
    \item \textbf{Structurelle :} Spécialisation primaire (DTE), faible transformation, concentration produits/partenaires.
    \item \textbf{Infrastructurelle :} Port goulot, corridors routiers/ferroviaires saturés/dégradés, énergie chère/instable.
    \item \textbf{Institutionnelle :} Tracasseries, corruption, lenteur administrative, fiscalité complexe.
    \item \textbf{Sécuritaire :} 3 foyers actifs (Extrême-Nord, NOSO, Est) coupant les axes d'export/import.
    \item \textbf{Statistique :} Aveuglement sur l'informel (transfrontalier, urbain) = mauvais diagnostic.
\end{enumerate}
\end{aretenir}

\coursec{Dossier 5 : Les marchés frontaliers du Cameroun}

Après avoir analysé les grands équilibres macroéconomiques des échanges du Cameroun avec le reste du monde — balance commerciale structurellement déficitaire, dépendance aux matières premières, détérioration des termes de l'échange — nous changeons d'échelle. Zoomons maintenant sur les \emph{frontières}, là où la mondialisation se vit au quotidien, loin des statistiques des douanes centrales. Les marchés frontaliers sont le visage concret, parfois chaotique, toujours vital, de l'intégration régionale \emph{par le bas}. Ils illustrent une géographie des flux qui ne respecte pas toujours les cadres institutionnels de la CEMAC ou de la CEDEAO, mais qui structure profondément la vie des populations riveraines.

\courssub{Localisation et présentation des principaux marchés frontaliers}

Le Cameroun partage plus de 4\,900\,km de frontières avec six pays. À chaque poste frontière majeur correspond un marché, nœud d'échanges où se côtoient commerce formel (déclaré, taxé) et commerce informel (toléré ou réprimé). Le programme en retient cinq emblématiques, répartis sur trois régions frontalières distinctes.

\begin{popfigure}
\includegraphics[width=\linewidth,height=10cm,keepaspectratio]{N16/cmr_map.jpg}
\legende{Localisation des marchés frontaliers du programme. Sur la carte : Banki se situe à la frontière nord-ouest avec le Nigeria (Extrême-Nord), Mbaïboum à la frontière orientale avec la République centrafricaine (axe Garoua-Boulaï--Bangui), Kyé-Ossi et Abam-Minko à l'extrême sud, vers le Gabon et la Guinée équatoriale (Ambam est lisible sur la carte), Ekondo-Titi dans le Sud-Ouest, à l'ouest du Mont Cameroun, près du Nigeria. Ces marchés ne sont pas nommés sur la carte : l'élève les place d'après leur position par rapport aux villes et aux frontières. Exemples de flux (hors carte) : bétail et oignons vers le Nigeria, produits manufacturés vers le Gabon. \\ Source : Wikimedia Commons, CIA, domaine public.}
\end{popfigure}

\begin{itemize}
    \item \textbf{Banki (Extrême-Nord, frontière nigériane) :} Situé à quelques kilomètres de Bama (Nigéria), sur l'axe Maroua–Maiduguri. C'est le plus grand marché frontalier du Nord-Cameroun. Il draine les flux du bassin du Lac Tchad (bétail, oignons, mil, maïs) vers le grand marché nigérian de Maiduguri et inversement (tissus, carburant, produits manufacturés asiatiques).
    \item \textbf{Mbaïboum (Est, frontière centrafricaine) :} Sur l'axe Garoua-Boulaï–Bangui. Point de passage obligé pour le bétail centrafricain vers les abattoirs de Douala et Yaoundé, et pour l'approvisionnement de la RCA en produits camerounais (ciment, bière, savon, farine).
    \item \textbf{Abam-Minko (Sud, frontière gabonaise) :} Sur l'axe Ebolowa–Bitam–Libreville. Échange de produits vivriers (plantain, manioc, macabo) du Sud-Cameroun contre produits manufacturés, carburant (historiquement moins cher au Gabon) et bois gabonais.
    \item \textbf{Kyé-Ossi (Sud, triple frontière Gabon / Guinée équatoriale / Cameroun) :} Situation géographique unique. Depuis ce seul point, les commerçants camerounais accèdent aux marchés de Kokokou (Gabon) et d'Ebebiyín (Guinée équatoriale).
    \item \textbf{Ekondo-Titi (Sud-Ouest, frontière nigériane, zone Bakassi) :} Marché fluvial et côtier. Spécialisé dans les produits halieutiques (poisson fumé, crevettes) et le carburant, il relie la zone côtière camerounaise aux marchés de Calabar et Port Harcourt (Nigéria).
\end{itemize}

\begin{savaistu}[Kyé-Ossi : une fenêtre sur l'espace CEMAC]
À Kyé-Ossi (arrondissement de Kyé-Ossi, département de la Ntem, Région du Sud), la frontière est un trait fin sur la carte mais une réalité économique poreuse. Le marché hebdomadaire (souvent le mardi ou le vendredi) voit affluer des acheteurs de \textbf{Kokokou} (Gabon, à 5\,km) et d'\textbf{Ebebiyín} (Guinée équatoriale, à 8\,km). Un atout majeur : les trois pays partagent la \textbf{même monnaie, le FCFA (XAF)}, émise par la BEAC. Contrairement à la frontière nigériane (Banki, Ekondo-Titi) où le change FCFA/Naira est un coût et un risque majeurs, ici \textbf{pas de changeur nécessaire} pour les transactions courantes. Cela fluidifie considérablement les échanges de vivriers (plantain, manioc) contre carburant, bois ou produits manufacturés. La complexité réside davantage dans les formalités douanières (trois régimes différents malgré la CEMAC) et les barrières linguistiques (Français / Espagnol / Fang/Bulu).
\end{savaistu}

\courssub{Nature des flux et organisation spatiale d'un marché frontalier}

Les marchandises qui traversent ces frontières dessinent une géographie de la complémentarité (ressources vs besoins) et de l'arbitrage (prix, taxes, change).

\begin{definition}[Marché frontalier et commerce transfrontalier]
Un \cle{marché frontalier} est un lieu d'échange périodique (hebdomadaire, bi-hebdomadaire) ou permanent, implanté à proximité immédiate d'une frontière internationale, où s'opèrent des transactions entre résidents de pays limitrophes.
Le \cle{commerce transfrontalier} regroupe l'ensemble de ces échanges. On distingue :
\begin{itemize}
    \item Le \textbf{commerce formel} : déclaré en douane, acquittant les droits et taxes, bénéficiant des accords régionaux (ex. exemption de droits dans la CEMAC). Il génère des statistiques officielles.
    \item Le \textbf{commerce informel (ou informel transfrontalier - CIT)} : non déclaré, échappant aux formalités douanières et fiscales, mais souvent toléré socialement. Il concerne des petits volumes (commerce de « valise » ou « sac ») ou des produits interdits/soumis à quotas. Il est \textbf{invisible dans les statistiques officielles} de la balance commerciale.
\end{itemize}
\end{definition}

Les produits échangés varient selon les frontières :
\begin{center}
\begin{tabularx}{\linewidth}{|l|X|X|}\hline
\rowcolor{popblueL}
\textbf{Frontière} & \textbf{Exportations camerounaises (exemples)} & \textbf{Importations (exemples)} \\ \hline
Nigéria (Banki, Ekondo-Titi) & Bétail sur pied, oignons, mil/maïs, poisson fumé (Ekondo), coton & Tissus (pagnes), carburant (essence de contrebande), motos, produits manufacturés asiatiques, riz \\ \hline
RCA (Mbaïboum) & Ciment, boissons, savon, farine, médicaments, produits manufacturés & Bétail sur pied, café, coton, diamant (artisanal), bois \\ \hline
Gabon / Guinée Eq. (Abam-Minko, Kyé-Ossi) & Vivriers (plantain, manioc, macabo, fruits), main-d'œuvre agricole & Carburant (subventionné au Gabon), bois grumes/planches, produits manufacturés, électronique \\ \hline
\end{tabularx}
\end{center}

\begin{popfigure}
\begin{tikzpicture}[scale=0.9]
% Fond : route frontalière
\draw[popline, line width=4pt] (-5,-2) -- (5,-2) node[midway, below, font=\small] {Route nationale / Piste};
% Ligne frontière (symbolique)
\draw[popdark, dashed, thick] (-5,0) -- (5,0) node[midway, above, font=\small\bfseries] {LIGNE FRONTIÈRE};
\node[popdark, font=\tiny] at (-4.5, 0.3) {CAMEROUN};
\node[popdark, font=\tiny] at (-4.5, -0.3) {PAYS VOISIN};

% Poste de douane / Police (Côté Cameroun)
\draw[popblue, fill=popblue!10, thick] (-4, 0.5) rectangle (-2, 1.8);
\node[popblue, font=\small\bfseries, align=center] at (-3, 1.15) {POSTE DE\\ DOUANE / POLICE};
\draw[popblue, ->, thick] (-3, 0.5) -- (-3, 0.1);

% Zone de stockage / Poids lourds (Côté Cameroun)
\draw[poporange, fill=poporange!10, thick] (1, 0.5) rectangle (4, 2);
\node[poporange, font=\small\bfseries, align=center] at (2.5, 1.25) {ZONE DE STOCKAGE\\ POIDS LOURDS (Bétail, Ciment)};
\draw[poporange, ->, thick] (2.5, 0.5) -- (2.5, 0.1);

% Marché principal (Côté Cameroun, près de la barrière)
\draw[popgreen, fill=popgreen!10, thick, rounded corners=5pt] (-1.5, -1.5) rectangle (2.5, 0.3);
\node[popgreen, font=\small\bfseries] at (0.5, -0.6) {MARCHÉ PRINCIPAL (Dalle / Hangars)};
% Étals
\foreach \x in {-1, 0, 1, 2} {
    \draw[popgreen, thick] (\x, -1.5) -- (\x, -0.3);
    \node[font=\tiny] at (\x, -0.9) {Étals};
}

% Zone de change / Services (informel)
\draw[poppink, fill=poppink!10, thick, rounded corners=3pt] (-4, -1.5) rectangle (-1.8, -0.5);
\node[poppink, font=\small\bfseries, align=center] at (-2.9, -1) {CHANGEURS\\ TRANSPORTEURS\\ PORTEURS};

% Flux piétons / petits colis (traversée frontière)
\draw[poppurple, thick, ->, >=Stealth, decorate, decoration={snake, amplitude=1pt, segment length=5pt}] (-0.5, 0) -- (-0.5, -1.2) node[midway, right, font=\tiny, poppurple] {Flux piétons / « sac »};
\draw[poppurple, thick, ->, >=Stealth, decorate, decoration={snake, amplitude=1pt, segment length=5pt}] (0.5, -1.2) -- (0.5, 0) node[midway, left, font=\tiny, poppurple] {Retour acheteurs};

% Flux poids lourds (formel)
\draw[poporange, thick, ->, >=Stealth] (3, 0.5) -- (3, 1.5) node[above, font=\tiny, poporange] {Sortie formelle};
\draw[poporange, thick, ->, >=Stealth] (-3, 1.5) -- (-3, 0.5) node[above, font=\tiny, poporange] {Entrée formelle};

% Légende intégrée au dessin (simplifiée)
\begin{scope}[xshift=6cm, yshift=-1cm]
\draw[popblue, fill=popblue!10] (0,0) rectangle (0.5,0.3); \node[right, font=\small] at (0.6,0.15) {Poste douane/police};
\draw[poporange, fill=poporange!10] (0,-0.5) rectangle (0.5,-0.2); \node[right, font=\small] at (0.6,-0.35) {Zone poids lourds (formel)};
\draw[popgreen, fill=popgreen!10] (0,-1) rectangle (0.5,-0.7); \node[right, font=\small] at (0.6,-0.85) {Marché principal (dalle)};
\draw[poppink, fill=poppink!10] (0,-1.5) rectangle (0.5,-1.2); \node[right, font=\small] at (0.6,-1.35) {Services informels (change, transport)};
\draw[poppurple, decorate, decoration={snake, amplitude=1pt, segment length=5pt}, thick, ->] (0.25,-2) -- (0.25,-2.3); \node[right, font=\small] at (0.6,-2.15) {Flux informels (piétons)};
\draw[popdark, dashed, thick] (0,-2.6) -- (0.5,-2.6); \node[right, font=\small] at (0.6,-2.6) {Ligne frontière};
\end{scope}
\end{tikzpicture}
\legende{Croquis légendé : organisation spatiale type d'un marché frontalier camerounais (ex. Banki ou Mbaïboum). On distingue l'espace formel (douane, zone poids lourds) de l'espace de marché mixte et des services informels (change, transport) collés à la ligne frontière.}
\end{popfigure}

\begin{methode}[Construire un croquis géographique légendé (compétence Bac)]
À partir d'une photo aérienne, d'une description ou d'une observation de terrain :
\begin{enumerate}
    \item \textbf{Analyser le document} : identifier les grands ensembles (route, frontière, bâtiments, zones d'étals, parkings, cours d'eau).
    \item \textbf{Choisir un fond de carte simplifié} : tracer les structures linéaires (route, frontière, rivière) et les limites des zones principales.
    \item \textbf{Sélectionner les figurés} :
        \begin{itemize}
            \item \textbf{Surfaces (zones)} : coloriés / hachurés (ex. marché, zone douane, parking).
            \item \textbf{Lignes (flux)} : flèches proportionnelles (épaisseur = volume), pleines (formel), pointillées/sinuées (informel).
            \item \textbf{Points (équipements)} : symboles (étoile = marché, carré = poste douane, triangle = changeur).
        \end{itemize}
    \item \textbf{Construire la légende ORGANISÉE} : regrouper par thèmes (ex. \textit{Infrastructures}, \textit{Activités commerciales}, \textit{Flux de marchandises}, \textit{Acteurs}). Hiérarchiser (gras pour les titres de thème).
    \item \textbf{Titrer précisément} : « Organisation spatiale du marché frontalier de Banki (Extrême-Nord) », dater, orienter (flèche Nord), échelonner (échelle approximative).
\end{enumerate}
\end{methode}

\courssub{Impacts économiques et sociaux : le moteur de la vie frontalière}

Ces marchés ne sont pas de simples lieux d'échange ; ils sont le poumon économique de zones souvent enclavées.

\begin{exemplebox}[Banki : le grenier et la boucherie de l'Extrême-Nord et du bassin tchadien]
Chaque semaine, ce sont \textbf{2\,000 à 5\,000 têtes de bétail} qui transitent par Banki vers le Nigéria (Maiduguri) et vers le Sud-Cameroun (Douala/Yaoundé via Maroua). Le marché approvisionne en \textbf{oignons} (production de la plaine du Logone, \textasciitilde 150\,000\,t/an) les marchés du Nord-Est nigérian. En retour, le carburant nigérian (essence « zoua-zoua ») alimente à bas prix l'économie de transport de toute la région de Maroua. Les recettes fiscales locales (taxes de marché, patentes) représentent jusqu'à \textbf{30 à 40\% du budget propre} de la commune de Waza/Kousseri. La fermeture de la frontière (2015-2017, crise Boko Haram) a fait chuter le prix du bétail de \textasciitilde 40\% côté camerounais et fait flamber le prix de l'essence de \textasciitilde 60\%, prouvant la dépendance structurelle.
\end{exemplebox}

\begin{aretenir}[Impacts croisés]
\begin{itemize}
    \item \textbf{Économiques :} Revenus directs pour des milliers de commerçants (dont \textasciitilde 60\% de femmes pour le vivrier), transporteurs (motos, camions), changeurs, porteurs (\emph{bayam-sellam}). Recettes communales (taxes). Approvisionnement en vivres des villes frontalières et des capitales (bétail centrafricain à Douala). Intégration régionale \emph{de facto} (CEMAC/CEDEAO vécue).
    \item \textbf{Sociaux :} Brassage ethnique et linguistique (Fulfulde, Haoussa, Arabe, Français, Anglais, Pidgin). Mariages transfrontaliers. Dynamisme urbain : Kousseri, Garoua-Boulaï, Kyé-Ossi, Ekondo-Titi sont nées ou grossissent grâce au marché. Émergence de services (hôtels, restaurants, transferts d'argent Mobile Money transfrontaliers).
\end{itemize}
\end{aretenir}

\courssub{Difficultés et enjeux : la frontière comme obstacle}

Malgré la vitalité de ces échanges, le quotidien des acteurs est semé d'embûches qui renchérissent les coûts de transaction et freinent l'intégration officielle.

\begin{attention}[Le piège statistique : l'iceberg invisible]
\textbf{Erreur fréquente au Bac :} « Le Cameroun exporte peu vers la RCA car les statistiques douanières montrent des flux faibles. »
\textbf{Réalité :} La majeure partie du commerce transfrontalier (bétail, vivres, carburant) est \textbf{informelle}. Elle n'apparaît \textbf{PAS} dans la balance commerciale officielle. Les échanges \emph{réels} sont 3 à 10 fois supérieurs aux échanges \emph{enregistrés} selon les études de la BEAC et de la CEMAC. Ne confondez jamais « statistiques douanières » et « réalité des flux ».
\end{attention}

Les difficultés majeures sont :
\begin{enumerate}
    \item \textbf{Tracasseries douanières, policières, gendarmes :} Multiplicité des barrages (parfois 10 à 15 entre Garoua-Boulaï et Bertoua), rackets, demandes de « pourboires » systématiques. Allongement des délais (un camion met 3 jours pour 300\,km).
    \item \textbf{Fraude et contrebande organisées :} Carburant (Nigéria $\to$ Cameroun), riz (interdit à l'import mais entrant par Banki/Ekondo), produits subventionnés (engrais, ciment) détournés vers le Nigéria ou la RCA. Cela crée une économie de rente pour quelques réseaux, au détriment du Trésor public.
    \item \textbf{Mauvais état des infrastructures :} Pistes impraticables en saison des pluies (axe Garoua-Boulaï–Bertoua, routes d'accès à Ekondo-Titi). Absence de stockage frigorifique pour le poisson/viande.
    \item \textbf{Insécurité :} \textbf{Boko Haram / ISWAP dans l'Extrême-Nord} (Banki, Fotokol, Amchidé) : attaques de marchés, enlèvements, pose d'engins explosifs, interdiction de la pêche sur le Lac Tchad. Dans le Nord-Ouest/Sud-Ouest (crise anglophone), Ekondo-Titi est une zone rouge, marché souvent désert ou déplacé.
    \item \textbf{Fluctuations monétaires :} Taux de change FCFA / Naira (Nigéria) très volatil. Dévaluation du Naira (2023-2024) : 1\,000 FCFA $\approx$ 1\,500 NGN (vs 450 NGN en 2020). Cela rend les produits camerounais très chers pour les Nigérians (baisse exportations) et les produits nigérians (carburant, riz) très bon marché pour les Camerounais (hausse importations informelles).
\end{enumerate}

\courssub{Activité guidée : Étude de cas — Le marché de Mbaïboum (Est-Cameroun)}

\textbf{Contexte :} Vous disposez de deux documents pour réaliser un croquis et un commentaire.
\begin{itemize}
    \item \textbf{Doc 1 (Photographie décrite) :} Vue aérienne du marché de Mbaïboum (jour de grand marché). Au premier plan, la route nationale N1 (bitume dégradé) coupe l'image horizontalement. Au Nord de la route (côté Cameroun), un grand hangar métallique (dalle couverte) dense en étals (vivriers, tissus). Au Sud, une zone découverte, poussiéreuse, avec des centaines de motos-taxis et des camions bennes stationnés. À l'extrême Sud, la barrière frontière (portique jaune/rouge) et deux bâtiments (Douane, Police). De l'autre côté de la barrière (Côté RCA), un enclos à bétail (parc à vaches) et des hangars à coton/café.
    \item \textbf{Doc 2 (Tableau de flux hebdomadaires estimés, source : étude GIZ/CEMAC 2022)} :
\end{itemize}

\begin{center}
\begin{tabular}{|l|c|c|c|}\hline
\rowcolor{popblueL}
\textbf{Produit} & \textbf{Sens Cameroun $\to$ RCA} & \textbf{Sens RCA $\to$ Cameroun} & \textbf{Caractère} \\ \hline
Bétail sur pied & 0 & 800 têtes & Formel (certif. vétérinaire) + Informel \\ \hline
Vivriers (manioc, maïs, banane) & 15 t & 2 t & Majoritairement informel \\ \hline
Carburant (essence/gazole) & 40\,000 L & 0 & Informel (contenance bidons) \\ \hline
Ciment / Matériaux & 200 t & 0 & Formel (facturé TVA) \\ \hline
Produits manufacturés (tissus, savon, motos) & 500 colis & 50 colis & Mixte \\ \hline
Coton / Café / Bois & 0 & 30 t & Formel (déclaré douane) \\ \hline
\end{tabular}
\end{center}

\begin{exoresolu}[Activité guidée : Croquis et commentaire de Mbaïboum]
\textbf{Énoncé :} À partir des documents 1 et 2, construisez un croquis légendé de l'organisation spatiale du marché de Mbaïboum. Rédigez un commentaire géographique structuré (introduction, analyse des flux, acteurs et enjeux, conclusion).

\tcblower
\textbf{Solution détaillée — Croquis (description textuelle pour reproduction) :}

\textit{Fond de carte :} Rectangle horizontal. Ligne médiane épaisse = Route N1. Ligne pointillée épaisse au Sud de la route = Frontière (barrière). Flèche Nord en haut à gauche.

\textit{Figurés de surface (zones) :}
\begin{itemize}
    \item \textcolor{popgreen}{\textbf{Zone Nord (Cameroun) :}} Rectangle hachuré vert = « Marché couvert (dalle/hangar) : vivriers, tissus, quincaillerie ».
    \item \textcolor{poporange}{\textbf{Zone Sud-Ouest (Cameroun) :}} Rectangle pointillé orange = « Zone de transit poids lourds / Parking motos/camions ».
    \item \textcolor{popblue}{\textbf{Point Sud-Est (Frontière) :}} Carré bleu = « Poste Douane / Police / Bureau CEMAC ».
    \item \textcolor{poppurple}{\textbf{Zone Sud (RCA) :}} Rectangle hachuré violet = « Parc à bétail / Hangars coton-café (Export RCA) ».
\end{itemize}

\textit{Figurés linéaires (flux) :}
\begin{itemize}
    \item \textbf{Flèches épaisses pleines (poporange) Nord $\leftarrow$ Sud :} « Bétail (800 t.), Coton, Café, Bois » (label : Flux formels + informels entrants).
    \item \textbf{Flèches moyennes pleines (popblue) Nord $\rightarrow$ Sud :} « Carburant (40\,000 L), Ciment (200 t), Vivriers (15 t), Manufacturés » (label : Flux sortants majoritairement informels sauf ciment).
    \item \textbf{Flèches sinuées pointillées (poppink) traversant la frontière n'importe où :} « Flux piétons / « sac » / Moto-taxis (vivriers, carburant en bidons) ».
\end{itemize}

\textit{Légende organisée (obligatoire) :}
\begin{center}
\begin{tabular}{ll}
\textbf{INFRASTRUCTURES} & \textbf{ACTIVITÉS / FLUX} \\
\textcolor{popblue}{$\square$} Poste douane/police & \textcolor{poporange}{\textbf{$\rightarrow$}} Bétail, coton, café (Entrée) \\
\textcolor{popgreen}{$\square$} Marché couvert (dalle) & \textcolor{popblue}{\textbf{$\rightarrow$}} Carburant, ciment, vivriers (Sortie) \\
\textcolor{poporange}{$\square$} Zone parking/transit & \textcolor{poppink}{\textbf{$\rightsquigarrow$}} Commerce informel piéton/moto \\
\textcolor{poppurple}{$\square$} Parc bétail / Hangars (RCA) & \\
\end{tabular}
\end{center}

\textbf{Commentaire type (structure Bac 16-18/20) :}

\textbf{Introduction :} Localisation (Mbaïboum, Est, axe Garoua-Boulaï–Bangui, frontière RCA), nature du document (croquis issu photo + stats), thèse : Mbaïboum est une interface vitale mais asymétrique, révélatrice des limites de l'intégration CEMAC formelle.

\textbf{I. Une organisation spatiale binaire : formel vs informel}
Le croquis met en évidence une structuration nette de part et d'autre de la frontière (ligne pointillée). Côté Cameroun, l'espace est compartimenté : le marché couvert (formel, taxé par la commune) au Nord, la zone de transit poids lourds (formel, douanier) au Sud-Ouest, le poste de contrôle (pouvoir régalien) à la limite. Côté RCA, l'espace est productif (parc bétail, hangars d'export). Mais la frontière est perméable : les flèches sinuées (informel) la traversent en dehors du poste, reliant directement le parc bétail aux acheteurs camerounais ou le marché couvert aux motos RCA.

\textbf{II. Des flux asymétriques et complémentaires (analyse Doc 2)}
Le tableau révèle une complémentarité structurelle : la RCA (zone d'élevage extensive) exporte du bétail (800 t./sem.) et des matières premières (coton, café, bois) vers le Cameroun (zone de consommation urbaine et portuaire). Inversement, le Cameroun (zone industrielle relative, subventions carburant) exporte massivement du carburant (40\,000 L/sem., \textbf{100\% informel}), du ciment (formel) et des vivriers/manufacturés. \textbf{Asymétrie majeure :} le carburant, produit stratégique, ne circule qu'informellement (différence de prix à la pompe \textasciitilde 300 FCFA/L au Cameroun vs 600+ FCFA en RCA officiellement, mais pénurie en RCA). Le ciment, lui, circule formellement (intégration CEMAC effective sur ce produit).

\textbf{III. Acteurs et enjeux : une intégration par le bas fragilisée}
Acteurs : éleveurs Peuls (transhumance transfrontalière), « bayam-sellam » femmes (vivriers), transporteurs (motos, camions), changeurs (FCFA commun CEMAC = avantage huge, pas de change), douaniers/policiers (rente de situation). Enjeux : la piste Garoua-Boulaï–Bertoua (dégradée) freine l'évacuation du bétail vers Douala (pertes de poids, mortalité). L'insécurité (coupeurs de route, ex-Seleka) sur l'axe RCA renchérit l'assurance transport. La non-enregistrement du carburant prive l'État camerounais de TVA et la BEAC de données sur la masse monétaire.

\textbf{Conclusion :} Mbaïboum illustre le paradoxe cemacquois : une monnaie unique (FCFA) qui devrait fluidifier les échanges, mais des barrières non-tarifaires (pistes, tracasseries, asymétrie des subventions carburant) qui poussent 80\% du commerce dans l'informel. L'intégration réelle se fait \emph{malgré} les États, par les acteurs locaux.
\end{exoresolu}

\begin{exercice}[Entraînement Bac — Sujet type]
À partir de la carte de localisation (figure 1) et de vos connaissances, montrez que les marchés frontaliers du Cameroun sont les « vraies » frontières de l'intégration régionale en Afrique centrale.
\textit{Indications :} Structurez en 3 parties : 1. Une géographie des flux qui ignore les limites administratives (complémentarités). 2. Des acteurs qui construisent l'intégration par le bas (réseaux, change, transport). 3. Des États qui peinent à encadrer (insécurité, tracasseries, stats manquantes). Utilisez les exemples de Banki (bétail/oignons), Mbaïboum (bétail/ciment), Kyé-Ossi (double accès).
\end{exercice}

\begin{corrige}[Exercice Entraînement Bac]
\textbf{Introduction :} Accroche (citation ou chiffre : « 80\% du commerce intra-CEMAC est informel »), définition du marché frontalier, thèse : ils sont le laboratoire réel de l'intégration, bien plus que les textes de la CEMAC.
\textbf{I. Des flux dictés par la géographie physique et humaine, pas par les traités.}
Complémentarités Nord (sahélien) / Sud (forestier) : Banki (bétail/oignons $\leftrightarrow$ carburant/tissus), Mbaïboum (bétail $\leftrightarrow$ ciment/vivriers), Abam/Kyé-Ossi (vivriers $\leftrightarrow$ carburant/bois). Monnaie unique FCFA (zone CEMAC) facilite Kyé-Ossi/Abam/Mbaïboum ; zone CEDEAO (Banki/Ekondo) nécessite change mais volume compense.
\textbf{II. Des acteurs transnationaux : les vrais artisans de l'intégration.}
Réseaux ethniques (Peuls, Haoussas, Bétis/Fangs transfrontaliers). Métiers nés de la frontière : changeurs (Kyé-Ossi double monnaie \textbf{corrigé : change inutile en zone FCFA, mais utile pour importations extra-CEMAC}), transporteurs motos (\"okada\" / \"bensikin\"), courtiers bétail. Villes-frontières devenues pôles urbains (Kousseri 200\,000 hab., Garoua-Boulaï, Kyé-Ossi).
\textbf{III. L'État en retard : régulation vs répression.}
Tracasseries = taxe occulte. Insécurité (Boko Haram, crise NOSO, coupeurs de route RCA) = coût de transaction énorme. Statistiques officielles aveugles = politiques publiques mal calibrées. Nécessité : corridors économiques, guichets uniques, formalisation progressive (ex. carte biométrique commerçants CEMAC).
\textbf{Conclusion :} L'intégration se fait par les marchés ; aux États de la suivre (infrastructures, simplification douanière, sécurité).
\end{corrige}

\begin{aretenir}[Synthèse du Dossier 5 — Les 5 points-clés pour le Bac]
\begin{enumerate}
    \item \textbf{5 marchés, 3 zones :} Banki/Ekondo-Titi (Nigéria/CEDEAO), Mbaïboum (RCA/CEMAC), Abam-Minko/Kyé-Ossi (Gabon/Guinée Eq./CEMAC). Kyé-Ossi = triple frontière unique.
    \item \textbf{Flux complémentaires :} Élevage/Nord $\leftrightarrow$ Vivriers/Industriel/Sud. Carburant (Nigéria/Gabon) $\to$ Cameroun ; Bétail (RCA/Nord) $\to$ Douala/Yaoundé.
    \item \textbf{Double circuit :} Formel (douane, stats, CEMAC) vs Informel (majoritaire, invisible, vital, « sac », motos).
    \item \textbf{Impacts majeurs :} Revenus (femmes, jeunes), approvisionnement villes, recettes communes, urbanisation frontalière, brassage social.
    \item \textbf{Blocages :} Tracasseries (coût transaction), infrastructures (pistes), insécurité (Boko Haram, NOSO, RCA), change (Naira), absence stats réelles $\to$ sous-estimation intégration réelle.
\end{enumerate}
\end{aretenir}

\newpage

\coursec{Activité d'intégration}

\coursec{Les échanges du Cameroun avec l’extérieur}

\courssub{Objectifs de l'activité}

À l'issue de cette activité d'intégration, tu dois être capable de :

\begin{itemize}
\item \cle{inventorier}, \cle{classer} et \cle{interpréter} des documents statistiques et cartographiques portant sur les échanges extérieurs du Cameroun ;
\item \cle{construire} un diagramme en secteurs et \cle{tracer} une courbe d'évolution à partir d'un tableau de données ;
\item \cle{localiser} et \cle{légender} sur une carte muette les principaux partenaires commerciaux et les marchés frontaliers du Cameroun ;
\item \cle{rédiger} un diagnostic argumenté et \cle{formuler} des propositions de solutions réalistes, à la manière d'un rapport d'expert.
\end{itemize}

Cette activité mobilise donc l'ensemble des savoir-faire du programme : observer, localiser, décrire, représenter, schématiser, cartographier et interpréter.

\courssub{Situation}

Tu es membre d'un club de jeunes géographes du Lycée de Biyem-Assi à Yaoundé. À l'occasion de la semaine nationale de l'économie, le ministère du Commerce organise un concours intitulé « Jeunes experts du commerce extérieur camerounais ». Chaque équipe reçoit un dossier documentaire et doit produire un diagnostic complet du commerce extérieur du Cameroun, assorti de propositions concrètes. Ton équipe dispose des quatre documents suivants.

\textbf{Document 1 : structure des exportations du Cameroun (ordres de grandeur, en \% de la valeur totale exportée, moyenne 2018--2022).}

\begin{center}
\begin{tabularx}{\linewidth}{|l|X|}\hline
\rowcolor{popblueL} \textbf{Produit exporté} & \textbf{Part approximative (\%)} \\ \hline
Pétrole brut & 40 \\ \hline
Cacao en fèves & 15 \\ \hline
Bois (grumes et sciages) & 12 \\ \hline
Coton fibre & 7 \\ \hline
Café & 4 \\ \hline
Aluminium & 4 \\ \hline
Bananes & 3 \\ \hline
Caoutchouc naturel & 3 \\ \hline
Autres produits & 12 \\ \hline
\end{tabularx}
\end{center}

\textbf{Document 2 : évolution du solde de la balance commerciale du Cameroun (en milliards de FCFA, valeurs arrondies à visée pédagogique, source : INS / BEAC).}

\begin{center}
\begin{tabularx}{\linewidth}{|l|X|X|X|X|X|X|}\hline
\rowcolor{popblueL} \textbf{Année} & 2018 & 2019 & 2020 & 2021 & 2022 & 2023 \\ \hline
Solde (milliards FCFA) & $-500$ & $-600$ & $-900$ & $-800$ & $-400$ & $-600$ \\ \hline
\end{tabularx}
\end{center}

\textbf{Document 3 : principaux partenaires commerciaux du Cameroun (parts en valeur, ordres de grandeur récents).} À l'exportation : Chine (environ 25 \%), Pays-Bas, Inde, Italie, Espagne. À l'importation : Chine (environ 20 \%), France, Nigeria, Inde, Belgique. Les partenaires de la CEMAC (Tchad, Gabon, Congo, Guinée équatoriale, République centrafricaine) absorbent moins de 5 \% du commerce extérieur camerounais officiel, mais les échanges frontaliers non enregistrés (Nigeria, Gabon, Guinée équatoriale) sont très intenses.

\textbf{Document 4 : photographies décrites.} Photo A : marché de \cle{Banki} (Extrême-Nord, département du Logone-et-Chari, frontière avec le Nigeria) : camions et porteurs chargés de sacs de maïs, de niébé, de pagnes et de produits manufacturés ; monnaies multiples en circulation (FCFA, Naira). Photo B : marché de \cle{Kyé-Ossi} (Région du Sud, tripoint Cameroun--Gabon--Guinée équatoriale) : étals de vivres frais (manioc, plantain, légumes), files de motos-taxis, longue attente au poste de contrôle frontalier.

\courssub{Consignes}

Par groupes de quatre, et en respectant l'ordre des tâches :

\begin{enumerate}
\item À partir du document 1, classe les produits exportés en deux catégories (produits bruts / produits transformés) et déduis-en la nature dominante des exportations camerounaises.
\item Construis un diagramme en secteurs représentant la structure des exportations du document 1 (calcule l'angle de chaque secteur, en degrés, pour au moins les quatre premiers produits).
\item À partir du document 2, trace la courbe d'évolution du solde commercial entre 2018 et 2023, puis commente : le solde est-il excédentaire ou déficitaire ? Quelle année marque la plus forte dégradation et à quelle conjoncture mondiale peut-on la rattacher ?
\item Sur la carte muette du Cameroun et de l'Afrique centrale fournie : (a) localise et nomme les cinq autres États de la CEMAC ; (b) place par des points légendés les cinq marchés frontaliers étudiés (Banki, Mbaïboum, Abam-Minko, Kyé-Ossi, Ekondo-Titi) ; (c) représente par des flèches les grands flux d'exportation vers la Chine, l'Union européenne et le Nigeria. Rédige un titre et une légende complète.
\item À partir des photographies du document 4, décris les produits échangés sur les marchés frontaliers et relève deux difficultés concrètes des échanges transfrontaliers.
\item Rédige en une demi-page (15 à 20 lignes) un \cle{diagnostic argumenté} du commerce extérieur camerounais (forces, faiblesses, problèmes majeurs) suivi de \textbf{deux solutions concrètes} : la transformation locale des matières premières et la sécurisation/organisation des corridors frontaliers.
\end{enumerate}

Restitution : chaque groupe présente ses productions en 5 minutes ; la correction collective est construite au tableau.

\courssub{Ressources à mobiliser}

\begin{aretenir}[Boîte à outils de l'expert]
\begin{itemize}
\item \textbf{Balance commerciale} : différence entre la valeur des exportations et celle des importations (biens uniquement). Si exportations $<$ importations, le solde est \cle{déficitaire}.
\item \textbf{Termes de l'échange} : rapport entre l'indice des prix à l'exportation et l'indice des prix à l'importation. Ils se \cle{détériorent} quand les matières premières exportées se paient moins cher que les produits manufacturés importés.
\item \textbf{CEMAC} : Communauté économique et monétaire de l'Afrique centrale ; 6 membres : Cameroun, Tchad, Gabon, Congo, Guinée équatoriale, Centrafrique. Monnaie commune : le Franc CFA (BEAC).
\item \textbf{Diagramme en secteurs} : angle $= \dfrac{\text{part en \%}}{100}\times 360^{\circ}$.
\item \textbf{Croquis géographique} : titre précis, légende numérotée et classée, orientation (flèche du nord), échelle indicative.
\end{itemize}
\end{aretenir}

\courssub{Démarche et corrigé commenté}

\begin{exoresolu}[Diagnostic du commerce extérieur du Cameroun]
Énoncé : réaliser les six tâches des consignes à partir des quatre documents du dossier.
\tcblower
\textbf{Étape 1 — Classification des produits (consigne 1).} Sur huit grands postes, un seul est un produit industriel transformé : l'aluminium (4 \%, produit par Alucam à Édéa). Le pétrole brut, le cacao en fèves, le bois (majoritairement grumes, sciages peu transformés), le coton fibre, le café, les bananes et le caoutchouc sont des produits \cle{bruts} ou faiblement transformés (première transformation), soit environ 84 \% des exportations. Conclusion : le Cameroun exporte essentiellement des matières premières, ce qui l'expose à la \cle{détérioration des termes de l'échange} et à la volatilité des cours mondiaux.

\textbf{Étape 2 — Diagramme en secteurs (consigne 2).} Calcul des angles (arrondis au degré) :
\begin{align*}
\text{Pétrole} &: \frac{40}{100}\times 360 = 144^{\circ} \\
\text{Cacao} &: \frac{15}{100}\times 360 = 54^{\circ} \\
\text{Bois} &: \frac{12}{100}\times 360 = 43^{\circ} \\
\text{Coton} &: \frac{7}{100}\times 360 = 25^{\circ} \\
\text{Autres (Café 14° + Alu 14° + Bananes 11° + Caoutchouc 11° + Autres 43°)} &: 93^{\circ}
\end{align*}
On reporte ces secteurs au rapporteur, du plus grand au plus petit, dans le sens horaire, avec une couleur distincte par produit et une légende latérale. Le titre précise : « Structure des exportations du Cameroun (valeur, \%, moyenne 2018--2022) ».

\begin{popfigure}
\centering
\begin{tikzpicture}
\fill[popblue] (0,0) -- (0.00:2.3) arc (0.00:144.00:2.3) -- cycle;
\node[white,font=\scriptsize\bfseries] at (72.00:1.55) {40\,\%};
\fill[popblue] (3.0,1.90) rectangle ++(0.28,0.28); \node[anchor=west,font=\scriptsize] at (3.35,2.04) {Pétrole brut (40\,\%)};
\fill[poporange] (0,0) -- (144.00:2.3) arc (144.00:198.00:2.3) -- cycle;
\node[white,font=\scriptsize\bfseries] at (171.00:1.55) {15\,\%};
\fill[poporange] (3.0,1.48) rectangle ++(0.28,0.28); \node[anchor=west,font=\scriptsize] at (3.35,1.62) {Cacao en fèves (15\,\%)};
\fill[popgreen] (0,0) -- (198.00:2.3) arc (198.00:241.20:2.3) -- cycle;
\node[white,font=\scriptsize\bfseries] at (219.60:1.55) {12\,\%};
\fill[popgreen] (3.0,1.06) rectangle ++(0.28,0.28); \node[anchor=west,font=\scriptsize] at (3.35,1.20) {Bois (grumes/sciages) (12\,\%)};
\fill[poppurple] (0,0) -- (241.20:2.3) arc (241.20:266.40:2.3) -- cycle;
\node[white,font=\scriptsize\bfseries] at (253.80:1.55) {7\,\%};
\fill[poppurple] (3.0,0.64) rectangle ++(0.28,0.28); \node[anchor=west,font=\scriptsize] at (3.35,0.78) {Coton fibre (7\,\%)};
\fill[poppink] (0,0) -- (266.40:2.3) arc (266.40:280.80:2.3) -- cycle;
\fill[poppink] (3.0,0.22) rectangle ++(0.28,0.28); \node[anchor=west,font=\scriptsize] at (3.35,0.36) {Café (4\,\%)};
\fill[popgold] (0,0) -- (280.80:2.3) arc (280.80:295.20:2.3) -- cycle;
\fill[popgold] (3.0,-0.20) rectangle ++(0.28,0.28); \node[anchor=west,font=\scriptsize] at (3.35,-0.06) {Aluminium (4\,\%)};
\fill[popblue!50] (0,0) -- (295.20:2.3) arc (295.20:306.00:2.3) -- cycle;
\fill[popblue!50] (3.0,-0.62) rectangle ++(0.28,0.28); \node[anchor=west,font=\scriptsize] at (3.35,-0.48) {Bananes (3\,\%)};
\fill[poporange!50] (0,0) -- (306.00:2.3) arc (306.00:316.80:2.3) -- cycle;
\fill[poporange!50] (3.0,-1.04) rectangle ++(0.28,0.28); \node[anchor=west,font=\scriptsize] at (3.35,-0.90) {Caoutchouc (3\,\%)};
\fill[popgreen!50] (0,0) -- (316.80:2.3) arc (316.80:360.00:2.3) -- cycle;
\node[white,font=\scriptsize\bfseries] at (338.40:1.55) {12\,\%};
\fill[popgreen!50] (3.0,-1.46) rectangle ++(0.28,0.28); \node[anchor=west,font=\scriptsize] at (3.35,-1.32) {Autres (12\,\%)};
\draw[white,thick] (0,0) circle (2.3);
\end{tikzpicture}
\legende{Diagramme en secteurs : structure des exportations camerounaises (Document 1). Dominance du pétrole (40 \%) et des produits bruts.}
\end{popfigure}

\textbf{Étape 3 — Courbe du solde commercial (consigne 3).} On place les années en abscisse et le solde (milliards de FCFA) en ordonnée, avec un axe horizontal à 0. Les six points sont tous \emph{sous} l'axe des abscisses : la balance commerciale est \cle{structurellement déficitaire}. Le creux de 2020 ($-900$ milliards) correspond à la crise sanitaire mondiale (Covid-19) : effondrement des cours pétroliers au printemps 2020 et ralentissement brutal du commerce international. L'amélioration relative de 2022 s'explique par le rebond des prix du pétrole (contexte post-Covid et guerre en Ukraine), avant une nouvelle dégradation en 2023. Le déficit structurel s'explique par l'importation massive de produits manufacturés (machines, véhicules, médicaments), d'hydrocarbures raffinés (le Cameroun importe son essence/gazole malgré son pétrole brut) et de biens alimentaires (riz, poisson congelé, blé).

\begin{popfigure}
\centering
\begin{tikzpicture}
\begin{axis}[
    width=\linewidth, height=8cm,
    xlabel={Année}, ylabel={Solde (milliards FCFA)},
    xtick={2018,2019,2020,2021,2022,2023},
    ytick={-1000,-800,-600,-400,-200,0},
    ymin=-1000, ymax=100,
    axis lines=middle,
    grid=major,
    mark=*,
    popcurve,
    thick
]
\addplot coordinates {
    (2018,-500) (2019,-600) (2020,-900) (2021,-800) (2022,-400) (2023,-600)
};
\addplot[dashed, popmuted] coordinates {(2018,0) (2023,0)};
\end{axis}
\end{tikzpicture}
\legende{Évolution du solde commercial du Cameroun (2018--2023). Déficit structurel, creux majeur en 2020 (Covid-19).}
\end{popfigure}

\textbf{Étape 4 — Croquis : Les échanges extérieurs du Cameroun (consigne 4).}

\begin{popfigure}
\includegraphics[width=\linewidth,height=10cm,keepaspectratio]{N16/cmr_map.jpg}
\legende{Le Cameroun et ses voisins sur une carte routière et ferroviaire. On y lit, au sud-ouest, le golfe de Guinée avec les ports de Douala et de Kribi, la capitale Yaoundé, les routes (traits rouges) et la voie ferrée (Douala--Yaoundé--Ngaoundéré) ; autour du pays, ses partenaires terrestres : le Nigeria à l'ouest, le Tchad au nord-est, la République centrafricaine à l'est, le Congo, le Gabon et la Guinée équatoriale au sud. C'est sur ce fond que l'élève situe les États de la CEMAC, les cinq marchés frontaliers étudiés (Banki, Mbaïboum, Abam-Minko, Kyé-Ossi, Ekondo-Titi, non nommés sur la carte) et trace les flux d'exportation vers l'Europe et la Chine par la façade maritime. \\ Source : Wikimedia Commons, CIA, domaine public.}
\end{popfigure}

\textbf{Étape 5 — Lecture des photographies (consigne 5).} Produits échangés : denrées vivrières locales (maïs, niébé, manioc, plantain, légumes), produits manufacturés importés (pagnes, chaussures, quincaillerie), carburant (souvent de contrebande). Deux difficultés concrètes : 1) Lenteurs, tracasseries et corruption aux postes de contrôle (visible à Kyé-Ossi) qui renchérissent les coûts et découragent la formalisation ; 2) Insécurité physique (zone de Banki touchée par les exactions de Boko Haram) et faiblesse des infrastructures routières (pistes impraticables en saison des pluies), sans oublier la multiplicité des monnaies qui complique les transactions.

\textbf{Étape 6 — Diagnostic rédigé (consigne 6, modèle de réponse).} Le commerce extérieur du Cameroun présente des atouts réels : position géostratégique centrale en Afrique centrale, façade atlantique avec le port autonome de Douala (poumon économique qui dessert aussi le Tchad et la RCA), diversité des ressources exportables (pétrole, cacao, bois, coton, aluminium), marchés frontaliers dynamiques (Banki, Mbaïboum, Abam-Minko, Kyé-Ossi, Ekondo-Titi) qui assurent l'approvisionnement vivrier des zones limitrophes et créent de l'emploi informel. Mais il souffre de trois faiblesses majeures : une structure d'exportation dominée à plus de 80 \% par des matières premières brutes (pétrole, cacao, grumes), une balance commerciale chroniquement déficitaire (déficit 2018--2023 entre 400 et 900 milliards FCFA/an) et une intégration régionale formelle inachevée (commerce intra-CEMAC < 5 \%), masquant une intégration réelle mais informelle. Deux solutions s'imposent : d'une part, développer résolument la \cle{transformation locale} (usines de chocolat type Chococam, scieries et placage pour le bois, filatures CICAM pour le coton, réhabilitation de la SONARA pour le raffinage) afin d'exporter des produits à plus forte valeur ajoutée et d'améliorer les termes de l'échange ; d'autre part, \cle{sécuriser et moderniser les corridors frontaliers} (postes de contrôle uniques « guichet unique », bitumage des axes prioritaires comme Mora--Banki ou Ebolowa--Kyé-Ossi, formalisation du petit commerce transfrontalier via des régimes simplifiés) pour transformer les marchés de Banki ou Kyé-Ossi en véritables leviers d'intégration régionale et de recettes douanières.
\end{exoresolu}

\begin{attention}[Pièges fréquents au Bac]
\begin{itemize}
\item Ne confonds pas \cle{balance commerciale} (biens seuls) et \cle{balance des paiements} (toutes transactions : biens, services, revenus, transferts).
\item Un solde négatif ne signifie pas « aucune exportation » : le Cameroun exporte pour des milliers de milliards, mais importe encore davantage.
\item Sur le croquis, une flèche de flux sans légende, sans titre ni sans flèche du nord vaut zéro point : soigne toujours la forme (titre, légende numérotée, orientation, échelle).
\item Les échanges frontaliers sont largement informels : ne les confonds pas avec le commerce officiel intra-CEMAC, qui reste marginal dans les statistiques.
\item « Détérioration des termes de l'échange » $\neq$ « déficit commercial ». Le premier est un rapport de prix, le second une différence de valeurs.
\end{itemize}
\end{attention}

\courssub{Bilan de l'activité}

Cette activité a permis de mettre en évidence trois idées forces de la leçon. D'abord, la structure des échanges du Cameroun est typique d'une économie de rente : exportation de produits bruts (pétrole, cacao, bois, coton) contre importation de produits manufacturés, d'hydrocarbures raffinés et de denrées alimentaires, d'où une balance commerciale structurellement déficitaire et une détérioration tendancielle des termes de l'échange. Ensuite, l'intégration régionale reste paradoxale : le commerce officiel au sein de la CEMAC est marginal (< 5 \%), alors que les marchés frontaliers (Banki, Mbaïboum, Abam-Minko, Kyé-Ossi, Ekondo-Titi) témoignent d'une intégration réelle, vivrière et sociale, mais freinée par l'insécurité (Boko Haram au Nord), les tracasseries administratives, le mauvais état des corridors et l'informel qui prive l'État de recettes. Enfin, sur le plan méthodologique, tu as mobilisé l'ensemble des savoir-faire du programme : lire et classer des statistiques, calculer des angles pour un diagramme en secteurs, tracer et commenter une courbe d'évolution, construire un croquis géographique légendé et rédiger un raisonnement géographique argumenté — autant de compétences directement exigibles au Baccalauréat.

\newpage

\coursec{Méthodes et exercices résolus}

\coursec{Cadre général et notions clés des échanges extérieurs}

\begin{experience}[Situation d'entrée : le port de Douala, poumon de l'économie]
Chaque jour, des centaines de camions font la queue au \cle{port autonome de Douala} (premier port d'Afrique centrale, \SI{95}{\%} du trafic maritime national) pour charger des grumes, du cacao, du coton ou décharger du riz, du poisson congelé, des véhicules et des machines. Ce ballet incessant illustre l'\cle{extraversion} de l'économie camerounaise : le pays vend à l'extérieur ses matières premières et achète à l'extérieur l'essentiel de ses biens d'équipement et de consommation. Comprendre ces échanges, c'est comprendre la place du Cameroun dans la \cle{mondialisation}.
\end{experience}

\begin{definition}[Vocabulaire de base du commerce extérieur]
\begin{itemize}
\item \textbf{Exportation (X) :} vente de biens ou services à l'étranger (sortie du territoire, entrée de devises).
\item \textbf{Importation (M) :} achat de biens ou services à l'étranger (entrée sur le territoire, sortie de devises).
\item \textbf{Partenaire commercial :} pays avec lequel s'effectuent des échanges (client ou fournisseur).
\item \textbf{Balance commerciale :} différence entre la valeur des exportations et celle des importations sur une période donnée ($X - M$).
\item \textbf{Taux de couverture :} rapport entre exportations et importations en valeur ($\frac{X}{M} \times 100$), exprimé en \%. Il mesure la capacité des exportations à financer les importations.
\item \textbf{Termes de l'échange (TE) :} rapport entre l'indice des prix à l'exportation et l'indice des prix à l'importation ($\frac{I_{px}}{I_{pm}} \times 100$). Ils mesurent le pouvoir d'achat des exportations sur les importations.
\end{itemize}
\end{definition}

\begin{aretenir}[Formules indispensables pour le Bac]
\begin{align*}
\text{Solde commercial} &= X - M \quad (\text{en milliards FCFA}) \\
\text{Taux de couverture (TC)} &= \frac{X}{M} \times 100 \quad (\text{en \%}) \\
\text{Termes de l'échange (TE)} &= \frac{\text{Indice prix exportations}}{\text{Indice prix importations}} \times 100 \quad (\text{base 100 à une année de référence})
\end{align*}
\emph{Interprétation rapide :} Solde $>0$ (excédent), $<0$ (déficit) ; TC $>100$ (excédent), $<100$ (déficit) ; TE $>100$ (favorables), $<100$ (défavorables).
\end{aretenir}

\courssub{Méthode 1 : Lire et exploiter un tableau statistique du commerce extérieur}

\begin{methode}[Observer, décrire et interpréter un tableau de données commerciales]
\begin{enumerate}
\item \textbf{Lire le titre complet} du tableau : quoi (les échanges), où (le Cameroun), quand (l'année ou la période). C'est la première phrase de ta réponse.
\item \textbf{Repérer les unités} (milliards de FCFA, tonnes, \%) et les sources (INS, MINFI, Banque Mondiale, FMI) : ne jamais mélanger des valeurs et des volumes.
\item \textbf{Dégager les faits saillants} : la valeur la plus forte, la plus faible, les évolutions (hausse, baisse, rupture). Cite toujours les chiffres exacts pour appuyer chaque affirmation.
\item \textbf{Calculer les indicateurs simples} : part en pourcentage, taux de variation, rapport exportations/importations.
\item \textbf{Interpréter} : relier chaque fait à une explication géographique (structure de l'économie camerounaise, cours mondiaux des matières premières, partenaires dominants).
\item \textbf{Conclure} en une phrase qui répond au problème posé par le sujet.
\end{enumerate}
\textbf{Formules à retenir :}
\[ \text{Part en \%} = \frac{\text{valeur de la catégorie}}{\text{valeur totale}} \times 100 \qquad ; \qquad \text{Taux de variation} = \frac{V_{t} - V_{0}}{V_{0}} \times 100 \]
\end{methode}

\begin{exoresolu}[Lecture d'un tableau des exportations camerounaises]
\textbf{Énoncé (type Bac).} Le tableau ci-dessous donne la structure approximative des exportations du Cameroun en valeur (ordres de grandeur réalistes, arrondis, année de référence récente) :
\begin{center}
\begin{tabularx}{\linewidth}{|l|X|X|}\hline
\rowcolor{popblueL} \textbf{Produit} & \textbf{Valeur (milliards FCFA)} & \textbf{Part (\%)} \\ \hline
Pétrole brut & 900 & \ldots \\ \hline
Cacao (fèves et produits) & 450 & \ldots \\ \hline
Bois et grumes & 350 & \ldots \\ \hline
Coton (fibre et graines) & 200 & \ldots \\ \hline
Autres produits (aluminium, banane, café, artisanat...) & 600 & \ldots \\ \hline
\textbf{Total} & \textbf{2500} & \textbf{100} \\ \hline
\end{tabularx}
\end{center}
1) Complétez la colonne des pourcentages. 2) Décrivez et interprétez la structure des exportations camerounaises en deux paragraphes.
\tcblower
\textbf{Solution détaillée.}

\textbf{1) Calcul des pourcentages.} On applique la formule $\text{part} = \dfrac{\text{valeur}}{\text{total}} \times 100$ :
\begin{align*}
\text{Pétrole} &= \frac{900}{2500} \times 100 = 36\ \% \\
\text{Cacao} &= \frac{450}{2500} \times 100 = 18\ \% \\
\text{Bois} &= \frac{350}{2500} \times 100 = 14\ \% \\
\text{Coton} &= \frac{200}{2500} \times 100 = 8\ \% \\
\text{Autres} &= \frac{600}{2500} \times 100 = 24\ \%
\end{align*}
Vérification : $36 + 18 + 14 + 8 + 24 = 100\ \%$. Le calcul est cohérent.

\textbf{2) Description et interprétation.}

\emph{Description.} Les exportations du Cameroun sont dominées par le pétrole brut, qui représente à lui seul 36 \% de la valeur totale exportée (900 milliards de FCFA sur 2500). Viennent ensuite le cacao (18 \%), le bois (14 \%) et le coton (8 \%). Les quatre premiers produits, tous des matières premières brutes ou peu transformées, totalisent $36 + 18 + 14 + 8 = 76\ \%$ des exportations. L'aluminium (Alucam/Edéa) et la banane (Plantations du Haut Penja, CDC) entrent dans la catégorie « Autres ».

\emph{Interprétation.} Cette structure révèle une économie \emph{extravertie} et fortement \emph{spécialisée dans les produits primaires} : le Cameroun vend à l'extérieur du pétrole (bassins de la Sanaga maritime et du Rio del Rey, exploités par Perenco, SNH, Addax), du cacao (régions du Centre, Sud-Ouest, Littoral), du bois (Est, Sud, exploités par des compagnies forestières souvent étrangères) et du coton (Nord, filière Sodecoton). Cette spécialisation expose le pays aux fluctuations des cours mondiaux : une baisse du prix du baril ou de la fève de cacao réduit mécaniquement les recettes d'exportation et les rentrées budgétaires, sans que le pays maîtrise les prix (prix \emph{price-taker}).

\emph{Conclusion.} Le tableau montre que les exportations camerounaises reposent sur un petit nombre de matières premières non transformées, ce qui constitue une fragilité majeure du commerce extérieur du pays et freine la diversification économique prônée par la SND30.
\end{exoresolu}

\coursec{Les produits échangés : exportations et importations}

\begin{exemplebox}[Panier d'exportation : des ressources naturelles peu transformées]
\begin{itemize}
\item \textbf{Pétrole brut :} premier produit d'exportation en valeur (\SI{\sim 35-40}{\%}). Bassins : Rio del Rey (Sud-Ouest, offshore) et Logone-Birni / Sanaga maritime (onshore/offshore). Opérateurs : SNH (national), Perenco, Addax, Bowleven. Exporté brut via l'oléoduc Chad-Cameroun (Kribi) et le terminal de Limbé/Kribi.
\item \textbf{Cacao :} 2\textsuperscript{e} produit (\SI{\sim 15-20}{\%}). Premier exportateur africain de cacao brut (après la Côte d'Ivoire et le Ghana). Zones : Centre (Mbalmayo, Nanga Eboko), Sud-Ouest (Kumba, Mamfé), Littoral. Enjeu : transformation locale (usines SIC Cacaos, Chococam, Neo Industry) encore faible (\SI{< 30}{\%} des fèves transformées sur place).
\item \textbf{Bois :} grumes, sciages, placages, contreplaqué (\SI{\sim 10-15}{\%}). Régions de l'Est et du Sud. Certifications FSC/PEFC en progression. Ports : Douala, Kribi.
\item \textbf{Coton :} fibre et graines (\SI{\sim 3-8}{\%} selon années). Zone Nord (Maroua, Garoua). Filière Sodecoton. Usines d'égrenage sur place.
\item \textbf{Autres :} Aluminium (Alucam, énergie hydroélectrique), Banane (PHP, CDC, Boh Plantations), Café (arabica Ouest, robusta Littoral/Sud-Ouest), Caoutchouc, Produits halieutiques (crevettes, poissons).
\end{itemize}
\end{exemplebox}

\begin{exemplebox}[Panier d'importation : biens d'équipement, consommation et énergie]
\begin{itemize}
\item \textbf{Produits pétroliers raffinés (super, gasoil, kérosène) :} poste n°1 des importations (\SI{\sim 15-20}{\%}). Paradoxe : pays producteur de brut, mais raffinerie (SONARA, Limbé) vétuste et souvent à l'arrêt $\rightarrow$ dépendance aux importations de produits finis.
\item \textbf{Biens d'équipement / Machines / Matériel de transport :} (\SI{\sim 25-30}{\%}). Nécessaires aux chantiers d'infrastructures (barrages, routes, ports) et à l'industrialisation naissante.
\item \textbf{Produits alimentaires :} Riz (Asie), Blé (France, Russie, Ukraine), Poisson congelé (Chine, UE, Maroc), Lait, Huiles. Facture alimentaire lourde (\SI{> 1000}{milliards FCFA/an}).
\item \textbf{Produits chimiques / Pharmaceutiques / Plastiques :} (\SI{\sim 10-15}{\%}).
\item \textbf{Véhicules d'occasion :} import massifs (Europe, Asie, USA), impact sur parc automobile vieillissant et pollution.
\end{itemize}
\end{exemplebox}

\courssub{Méthode 2 : Calculer et commenter la balance commerciale et le taux de couverture}

\begin{methode}[Balance commerciale et taux de couverture]
\begin{enumerate}
\item \textbf{Identifier} les exportations (X) et les importations (M) en valeur, pour la même année et la même unité (milliards FCFA ou USD).
\item \textbf{Calculer le solde commercial} : $\text{Solde} = X - M$.
\item \textbf{Qualifier le solde} : si $X > M$, excédent ; si $X < M$, déficit ; si $X = M$, équilibre.
\item \textbf{Calculer le taux de couverture} : $\text{TC} = \dfrac{X}{M} \times 100$. Il se lit : « les exportations couvrent TC \% des importations ».
\item \textbf{Interpréter} : un TC inférieur à 100 \% signale une dépendance extérieure ; un TC supérieur à 100 \% indique que le pays exporte plus qu'il n'importe.
\item \textbf{Expliquer} le résultat par la nature des produits échangés (matières premières exportées contre biens manufacturés et produits alimentaires importés).
\end{enumerate}
\begin{attention}[Piège classique]
Ne confonds jamais \emph{solde} (une différence, en milliards de FCFA) et \emph{taux de couverture} (un rapport, en \%). Un déficit de 500 milliards et un taux de couverture de 80 \% sont deux façons différentes de dire la même réalité ; les deux doivent être donnés dans une copie complète.
\end{attention}
\end{methode}

\begin{exoresolu}[Calcul du solde et du taux de couverture]
\textbf{Énoncé (type Bac).} En 2022 (ordres de grandeur INS/Banque Mondiale), le Cameroun a exporté pour environ 3500 milliards de FCFA et importé pour environ 4800 milliards de FCFA. 1) Calculez le solde commercial et le taux de couverture. 2) Commentez les résultats en cinq à huit lignes.
\tcblower
\textbf{Solution détaillée.}

\textbf{1) Calculs.}
\begin{align*}
\text{Solde commercial} &= X - M = 3500 - 4800 = -1300 \text{ milliards de FCFA} \\
\text{Taux de couverture} &= \frac{X}{M} \times 100 = \frac{3500}{4800} \times 100 \approx 72{,}9\ \%
\end{align*}

\textbf{2) Commentaire.} Le solde commercial est négatif ($-1300$ milliards de FCFA) : la balance commerciale du Cameroun est donc \emph{déficitaire}. Le taux de couverture de 72,9 \% signifie que les recettes d'exportation ne financent qu'environ les trois quarts des dépenses d'importation ; le quart restant doit être financé par l'endettement extérieur, l'aide publique au développement ou les transferts de la diaspora. Ce déficit structurel s'explique par la nature des échanges : le Cameroun exporte des matières premières (pétrole, cacao, bois, coton, aluminium) dont les prix sont fixés sur les marchés mondiaux (volatils, tendance baissière long terme pour les non-énergétiques), et importe des biens manufacturés (machines, véhicules, produits pharmaceutiques), des hydrocarbures raffinés et des produits alimentaires (riz, poisson, blé) à forte valeur ajoutée. \emph{Conclusion :} la structure extravertie de l'économie camerounaise entretient un déficit commercial chronique, qui fragilise la monnaie (FCFA), les réserves de change et les finances publiques.
\end{exoresolu}

\coursec{Les échanges avec les partenaires africains et la CEMAC}

\begin{savaistu}[La CEMAC et le rôle moteur du Cameroun]
La \cle{CEMAC} (Communauté Économique et Monétaire de l'Afrique Centrale) regroupe 6 États : Cameroun, Gabon, Congo, Tchad, RCA, Guinée équatoriale. Monnaie commune : FCFA (garantie par la France, parité fixe 1 EUR = 655,957 FCFA). Marché commun : \SI{60}{millions} d'habitants, PIB \SI{\sim 100}{milliards USD}. Le Cameroun en est la \cle{locomotive économique} : \SI{\sim 45-50}{\%} du PIB de la zone, premier partenaire commercial intra-zone pour la plupart des membres, plaque tournante logistique (port de Douala, corridor routier Douala--N'Djamena, Douala--Bangui, chemin de fer Camrail).
\end{savaistu}

\begin{aretenir}[Le paradoxe camerounais en Afrique centrale]
\begin{itemize}
\item \textbf{Poids économique dominant :} Cameroun = \SI{\sim 50}{\%} PIB CEMAC, \SI{\sim 40}{\%} population.
\item \textbf{Commerce intra-CEMAC marginal :} \textbf{moins de 5 \%} des échanges totaux du Cameroun (ordres de grandeur INS/Douanes). Le Cameroun exporte vers ses voisins : produits agroalimentaires (boissons, savons, huiles, pâtes), ciment (Cimencam), denrées vivrières (oignons du Nord, bananes), carburant (réexport), matériaux de construction. Il importe : bois (Congo, Gabon), produits pétroliers (Gabon, Congo - mais faible), bétail (Tchad, RCA - informel).
\item \textbf{Partenaires africains hors CEMAC :} \textbf{Nigeria} (premier partenaire africain, frontière poreuse, échanges massifs souvent informels : carburant, textiles, véhicules, denrées), Maroc (banques, engrais, ciment, industries), Afrique du Sud, Égypte.
\end{itemize}
\end{aretenir}

\courssub{Méthode 5 : Rédiger un commentaire de document sur les échanges régionaux (CEMAC)}

\begin{methode}[Commenter un document sur l'intégration régionale]
\begin{enumerate}
\item \textbf{Présenter le document} en deux lignes : nature (carte, tableau, texte), auteur ou source, date, thème.
\item \textbf{Définir les termes clés} du sujet : CEMAC (créée 1994, 6 membres), intégration régionale, libre-échange, union douanière, Tarif Extérieur Commun (TEC).
\item \textbf{Décrire les faits} : poids faible du commerce intra-CEMAC (ordre de grandeur : \SI{< 5}{\%} des échanges du Cameroun), domination des partenaires extra-africains (UE, Chine).
\item \textbf{Expliquer cette faiblesse} (causes structurelles) :
    \begin{itemize}
    \item Économies \emph{concurrentes} et non complémentaires (toutes exportent matières premières : pétrole, bois, minerais ; toutes importent manufacturés).
    \item Infrastructures de transport déficientes (routes dégradées, pas de rail inter-États, ports saturés).
    \item Barrières non tarifaires : tracasseries aux frontières (contrôles multiples, taxes illégales), non-application effective du libre-échange communautaire (certificats d'origine, TEC mal appliqué).
    \item Faible industrialisation $\rightarrow$ peu de biens manufacturés à s'échanger.
    \end{itemize}
\item \textbf{Nuancer} : le Cameroun est le « grenier » et la locomotive ; flux réels (formels + informels) sous-estimés par les stats officielles.
\item \textbf{Conclure} par une ouverture : enjeux de la \cle{ZLECAf} (Zone de Libre-Échange Continentale Africaine, phase opérationnelle 2021), nécessité de lever les obstacles logistiques et douaniers.
\end{enumerate}
\end{methode}

\begin{exoresolu}[Commentaire : le commerce intra-CEMAC]
\textbf{Énoncé (type Bac).} Document : « En valeur, les échanges entre les six pays de la CEMAC représentent moins de 5 \% de leur commerce extérieur total, alors que le Cameroun fournit à lui seul près de la moitié du PIB de la sous-région. » À partir de ce document et de vos connaissances, montrez que les échanges du Cameroun avec ses partenaires africains restent marginaux, puis expliquez cette situation (10 à 15 lignes).
\tcblower
\textbf{Solution détaillée.}

\emph{Présentation.} Le document est un extrait statistique portant sur le commerce intra-régional en Afrique centrale ; il oppose la faiblesse des échanges intra-CEMAC (moins de 5 \% du commerce total) au poids économique dominant du Cameroun dans la sous-région (environ la moitié du PIB des six pays).

\emph{Constat.} Les échanges du Cameroun avec ses partenaires africains sont marginaux dans les statistiques officielles : l'essentiel du commerce extérieur camerounais se fait avec l'Union européenne et l'Asie (Chine en tête), tandis que le Nigeria, le Tchad, le Gabon ou le Congo n'absorbent qu'une faible part des exportations officielles (ordre de grandeur : \SI{2-4}{\%} chacun). Le Cameroun vend pourtant à ses voisins de la CEMAC des produits agroalimentaires (boissons, huiles, savonnerie, pâtes alimentaires), du ciment (Cimencam), des denrées vivrières (oignons du Nord, bananes plantains) et du carburant, et joue le rôle de \cle{plaque tournante} logistique : le corridor Douala--N'Djamena et Douala--Bangui achemine l'essentiel des importations du Tchad et de la RCA, pays enclavés.

\emph{Explications.} Cette faiblesse statistique tient d'abord à la \emph{similitude des structures économiques} : les pays de la CEMAC exportent tous des matières premières (pétrole, bois, cacao, coton, minerais) et importent tous des biens manufacturés ; ils n'ont donc que peu de choses à s'échanger dans le cadre d'une division internationale du travail classique. Ensuite, les \emph{infrastructures} sont déficientes : routes dégradées (axe Douala--Bangui souvent impraticable en saison des pluies), absence de liaison ferroviaire inter-États, frontières longues et coûteuses à franchir. Enfin, les \emph{tracasseries administratives et douanières} (multiplication des contrôles, taxes illégales, non-application effective du libre-échange communautaire, lenteur des procédures) découragent les opérateurs formels, ce qui nourrit un vaste \emph{commerce informel transfrontalier} non comptabilisé dans les statistiques douanières.

\emph{Conclusion.} Malgré la CEMAC et l'union douanière, l'intégration commerciale régionale reste un chantier inachevé ; l'entrée en vigueur de la ZLECAf en 2021 ouvre une perspective de relance des échanges intra-africains, à condition de lever résolument les obstacles logistiques, douaniers et sécuritaires.
\end{exoresolu}

\coursec{Les échanges avec le reste du monde : Union européenne, Chine, nouveaux partenaires}

\begin{aretenir}[Géographie des partenaires hors Afrique (ordres de grandeur récents)]
\begin{center}
\begin{tabularx}{\linewidth}{|l|X|X|}\hline
\rowcolor{popblueL} \textbf{Zone / Pays} & \textbf{Rôle principal} & \textbf{Produits échangés (principaux)} \\ \hline
\textbf{Union européenne (UE)} & Partenaire historique, 1\textsuperscript{er} client (export) & \textbf{Export :} Pétrole, Cacao, Bois, Aluminium, Banane, Café. \textbf{Import :} Machines, Véhicules, Pharmaceutique, Produits chimiques, Vins/Alcools. Accords : APE (Accord de Partenariat Économique) intérimaire signé 2009, en application. \\ \hline
\textbf{Chine} & 1\textsuperscript{er} fournisseur (import), 2\textsuperscript{e} client & \textbf{Export :} Bois (grumes, sciages), Pétrole, Coton, Minerais. \textbf{Import :} Machines, Textiles/Habillement, Électronique, Véhicules, Matériaux construction, Riz. Investissements massifs (barrages, routes, port Kribi, stade Olembé). \\ \hline
\textbf{Inde} & Partenaire en forte hausse & \textbf{Export :} Pétrole, Coton, Bois. \textbf{Import :} Pharmaceutique (génériques), Véhicules (Tata, Mahindra), Machines, Riz, Produits chimiques. \\ \hline
\textbf{États-Unis} & Partenaire pétrolier / AGOA & \textbf{Export :} Pétrole (AGOA), Cacao, Bois. \textbf{Import :} Machines, Avions (Boeing), Produits high-tech, Céréales. \\ \hline
\textbf{Autres (Turquie, Russie, Pays du Golfe, Brésil...)} & Diversification & Engrais (Russie, Maroc), Ciment (Turquie, Maroc), Blé (Russie, France, Ukraine), Dattes, Investissements portuaires/hôteliers. \\ \hline
\end{tabularx}
\end{center}
\end{aretenir}

\courssub{Méthode 3 : Calculer et interpréter les termes de l'échange}

\begin{methode}[Les termes de l'échange : calcul et diagnostic]
\begin{enumerate}
\item \textbf{Définir} : les termes de l'échange mesurent le pouvoir d'achat des exportations d'un pays sur ses importations.
\item \textbf{Appliquer la formule} : $\text{TE} = \dfrac{\text{indice des prix à l'exportation}}{\text{indice des prix à l'importation}} \times 100$.
\item \textbf{Interpréter le niveau} : TE $> 100$ : les termes sont favorables ; TE $< 100$ : défavorables.
\item \textbf{Interpréter l'évolution} : si le TE baisse entre deux dates, on parle de \emph{détérioration des termes de l'échange} ; s'il monte, d'amélioration.
\item \textbf{Expliquer} : la détérioration traduit le fait que les prix des matières premières exportées progressent moins vite (ou baissent plus vite) que les prix des biens manufacturés importés (théorie de Prebisch-Singer).
\item \textbf{Conclure} sur les conséquences : pour importer la même quantité de biens, le pays doit exporter de plus en plus de matières premières (effet « course sur le tapis roulant »).
\end{enumerate}
\end{methode}

\begin{exoresolu}[Détérioration des termes de l'échange]
\textbf{Énoncé (type Bac).} Au Cameroun, l'indice des prix à l'exportation (base 100 en 2015) est passé à 92 en 2022, tandis que l'indice des prix à l'importation (même base) est passé à 115 sur la même période (données inspirées des tendances FMI/Banque Mondiale). 1) Calculez les termes de l'échange en 2015 et en 2022. 2) Montrez qu'il y a détérioration et expliquez ses conséquences pour le Cameroun.
\tcblower
\textbf{Solution détaillée.}

\textbf{1) Calculs.}
\begin{align*}
\text{TE}_{2015} &= \frac{100}{100} \times 100 = 100 \\
\text{TE}_{2022} &= \frac{92}{115} \times 100 = 80
\end{align*}

\textbf{2) Interprétation.} Les termes de l'échange sont passés de 100 à 80, soit une baisse de 20 points en sept ans : il y a bien \emph{détérioration des termes de l'échange}. Concrètement, en 2022, une unité de valeur exportée par le Cameroun ne permet d'acheter que 80 \% de ce qu'elle permettait d'importer en 2015. Le Cameroun a perdu 20 \% de son pouvoir d'achat à l'international. Cette dégradation s'explique par la baisse relative des prix des produits exportés (pétrole, cacao, coton, bois : indice passé de 100 à 92, volatilité forte) combinée à la hausse structurelle des prix des produits importés (biens d'équipement, produits alimentaires, pharmaceutiques : indice passé de 100 à 115, inflation importée). \emph{Conséquences :} pour importer la même quantité de machines, de riz ou de médicaments, le Cameroun doit exporter davantage de cacao, de pétrole ou de bois ; ses recettes réelles en devises se contractent, le déficit commercial s'aggrave, la capacité d'investissement de l'État diminue et la dette extérieure s'alourdit. \emph{Conclusion :} la détérioration des termes de l'échange est un mécanisme structurel qui appauvrit les pays exportateurs de matières premières comme le Cameroun au profit des pays exportateurs de produits manufacturés et de services à haute valeur ajoutée.
\end{exoresolu}

\courssub{Méthode 4 : Construire un croquis des échanges extérieurs du Cameroun}

\begin{methode}[Réaliser un croquis de flux commerciaux (savoir-faire cartographique Bac)]
\begin{enumerate}
\item \textbf{Tracer le fond de carte} : contour très simplifié du Cameroun (polygone grossier), flèche du nord (haut de page), échelle indicative (ex. 1 cm $\approx$ 200 km). Place les repères majeurs : Douala (port principal), Kribi (port en eau profonde), Yaoundé (capitale), les 5 frontières (Nigeria, Tchad, RCA, Congo, Gabon, Guinée équatoriale), l'océan Atlantique.
\item \textbf{Choisir des figurés adaptés} (et les mettre direct dans la légende) :
    \begin{itemize}
    \item Flèches d'épaisseur \textbf{proportionnelle} pour les flux : \textbf{vert} = exportations, \textbf{orange/rouge} = importations.
    \item Points (bleus) pour les ports, point (orange) pour la capitale.
    \item Traits fins pour les axes routiers/ferroviaires structurants (Corridors).
    \end{itemize}
\item \textbf{Hiérarchiser les flux} : Flux dominants (UE, Chine) = flèches très épaisses. Flux régionaux (Nigeria, CEMAC) = flèches fines. Flux secondaires = traits de flèches pointillés.
\item \textbf{Rédiger la légende organisée} : Titre du croquis en haut. Légende en bas ou sur le côté, groupée par thème (Flux, Lieux, Axes). Chaque figuré utilisé sur la carte apparaît \emph{exactement pareil} dans la légende.
\item \textbf{Vérifier les 4 éléments obligatoires} : Titre précis, Légende complète, Nord, Noms de lieux (villes, pays, mer). L'absence d'un seul élément coûte des points.
\end{enumerate}
\end{methode}

\begin{exoresolu}[Croquis des flux commerciaux du Cameroun]
\textbf{Énoncé (type Bac).} À l'aide de vos connaissances, construisez un croquis intitulé « Les principaux flux du commerce extérieur du Cameroun » faisant apparaître : les ports de sortie, les principaux partenaires (Union européenne, Chine, Nigeria, CEMAC) et les principaux produits exportés/importés.
\tcblower
\textbf{Solution détaillée (modèle de croquis attendu).}

\begin{popfigure}
\includegraphics[width=\linewidth,height=10cm,keepaspectratio]{N16/cmr_map.jpg}
\legende{Fond de carte pour le croquis des échanges extérieurs du Cameroun. La carte montre les deux ports de sortie sur le golfe de Guinée (Douala, Kribi), la capitale Yaoundé, les routes et la voie ferrée vers Ngaoundéré, et les pays voisins : Nigeria, Tchad, République centrafricaine, Congo, Gabon, Guinée équatoriale. Le croquis attendu y ajoute : 1. les ports (Douala, Kribi) et Yaoundé ; 2. les corridors routiers et ferroviaires vers le Tchad, la RCA et le Nigeria ; 3. les flèches d'exportation (épaisses, vers l'Union européenne et la Chine, depuis les ports) ; 4. les flèches d'importation (en sens inverse, couleur différente) ; 5. les flux régionaux, plus fins, vers le Nigeria et la CEMAC. Les flux ne figurent pas sur la carte d'origine. \\ Source : Wikimedia Commons, CIA, domaine public.}
\end{popfigure}

\textbf{Démarche de construction pas à pas.} (1) Contour + Nord + Échelle. (2) Ports (Douala, Kribi) + Yaoundé + Pays voisins. (3) Axes corridors (traits pointillés). (4) Flèches exports DOMINANTS (UE, Chine) depuis ports vers mer (flèches épaisses vertes). (5) Flèches imports DOMINANTS (UE, Chine, Inde) depuis mer vers ports (flèches épaisses orange). (6) Flux régionaux (Nigeria, CEMAC) flèches fines. (7) Légende groupée + Titre. \emph{Astuce :} Dessine les flèches exports et imports de part et d'autre du pays pour ne pas les superposer.
\end{exoresolu}

\coursec{Les problèmes des échanges extérieurs du Cameroun}

\begin{aretenir}[Diagnostic synthétique : 5 problèmes majeurs]
\begin{enumerate}
\item \textbf{Déficit commercial chronique \& Dépendance extérieure :} Taux de couverture souvent \SI{< 80}{\%}. Nécessité de financer le gap par endettement (Eurobonds, Chine, FMI) ou aide. Fragilité des réserves de change (BEAC).
\item \textbf{Détérioration tendancielle des termes de l'échange :} Prix matières premières (volatils, tendance baissière réelle hors chocs) vs Prix manufacturés (haussiers). Perte de pouvoir d'achat à l'importation. Théorie de Prebisch-Singer confirmée.
\item \textbf{Faible transformation locale (Faible valeur ajoutée) :} Exportation brute (grumes, fèves cacao, coton fibre, pétrole brut). Perte d'emplois industriels, de recettes fiscales, de maîtrise de la chaîne de valeur. Objectif SND30 : \cle{transformation locale} (usines cacao, raffinerie SONARA rénovée, usines bois, coton).
\item \textbf{Infrastructures logistiques insuffisantes et coûts élevés :} Port de Douala saturé (délais, coûts), port de Kribi sous-utilisé (manque d'hinterland connecté), routes corridors dégradées (Douala-Bangui, Douala-N'Djamena), chemin de fer Camrail vétuste (déraillements, lenteur). Coût transport = frein à la compétitivité.
\item \textbf{Informalité massive \& Évasion fiscale :} Commerce frontalier non enregistré (Nigeria, CEMAC), sous-facturation exportations (bois, minerais), sur-facturation importations, contrefaçon. Perte de recettes douanières et statistiques faussées.
\end{enumerate}
\end{aretenir}

\begin{attention}[Les 5 erreurs qui coûtent des points au Bac]
\begin{enumerate}
\item \textbf{Confondre solde commercial et taux de couverture.} Le solde est une différence en milliards de FCFA ($X - M$) ; le taux de couverture est un rapport en \% ($X/M \times 100$). Donner l'un sans l'autre, ou exprimer le taux en FCFA, fait perdre des points de calcul.
\item \textbf{Oublier la formule des termes de l'échange ou l'inverser.} C'est $\text{TE} = \dfrac{\text{prix à l'exportation}}{\text{prix à l'importation}} \times 100$. Inverser numérateur et dénominateur conduit à conclure à une amélioration alors qu'il y a détérioration : raisonnement entièrement faux.
\item \textbf{Affirmer que la CEMAC est le premier partenaire du Cameroun.} Faux : le commerce intra-CEMAC reste marginal (ordre de grandeur : moins de 5 \% des échanges). Les partenaires dominants sont l'Union européenne et la Chine. Citer des chiffres précis inventés est également sanctionné : utilise des ordres de grandeur signalés comme tels (\emph{« environ », « de l'ordre de », « près de »}).
\item \textbf{Présenter un croquis sans titre, sans légende ou sans nord.} Toute carte ou croquis doit comporter un titre précis, une légende complète (chaque figuré expliqué), la flèche du nord et des noms de lieux. Une flèche de flux sans légende ne vaut rien.
\item \textbf{Décrire sans interpréter (ou l'inverse).} Dans un commentaire de document, citer les chiffres sans les expliquer ne vaut que la moitié des points ; mais expliquer sans jamais citer le document est tout aussi pénalisé. Règle d'or : chaque idée = un fait chiffré du document + une explication géographique.
\end{enumerate}
\end{attention}

\coursec{Dossier 5 : Les marchés frontaliers du Cameroun}

\begin{experience}[Situation concrète : Le marché de Banki (Extrême-Nord)]
Chaque mardi, à \textbf{Banki} (frontière Nigeria, axe Mora--Maiduguri), des milliers de commerçants (femmes \emph{bayam-sellam}, transporteurs, changeurs) échangent oignons du Diamaré, poisson du Lac Tchad, carburant nigérian, tissus, bétail. Officiellement, ce marché « n'existe pas » dans les statistiques douanières. Pourtant, il alimente Maiduguri (Nigeria, \SI{2}{millions} d'hab.) et Mora/Kousseri (Cameroun). Depuis 2014, Boko Haram a perturbé ce poumon économique. Cette situation illustre la vitalité, l'informalité et la fragilité des \cle{marchés frontaliers}.
\end{experience}

\courssub{Méthode 6 : Mener une étude de cas sur un marché frontalier (Dossier 5)}

\begin{methode}[Étudier un marché frontalier : démarche en quatre temps (type Bac)]
\begin{enumerate}
\item \textbf{Localiser précisément} : Nom du marché, Région, Frontière concernée, Ville/axe de communication proche, Pays voisins impliqués. (Ex. : Kyé-Ossi, Région du Sud, triple frontière Cameroun--Gabon--Guinée équatoriale, au sud d'Ebolowa, axe Ebolowa--Oyem--Ebebiyin).
\item \textbf{Décrire les flux de marchandises} : Sens (Cameroun $\rightarrow$ Voisin / Voisin $\rightarrow$ Cameroun), Nature (denrées vivrières, carburant, produits manufacturés, bétail, produits agroalimentaires), Acteurs (commerçants grossistes, \emph{bayam-sellam}, transporteurs moto/taxi, douaniers), Périodicité (hebdomadaire, quinzainal).
\item \textbf{Évaluer les impacts} :
    \begin{itemize}
    \item \textbf{Économiques :} Revenus pour populations locales, approvisionnement villes frontalières à bas prix, devises échappant aux stats officielles (informel), dynamisation localités enclavées.
    \item \textbf{Sociaux :} Emplois (vente, transport, change), brassage populations (ethnies transfrontalières : Peuls, Mafa, Fang, Bassa), solidarité, mais aussi tensions (conflits fonciers, insécurité, contrebande armes/médicaments).
    \end{itemize}
\item \textbf{Identifier les difficultés \& Proposer des pistes} : Tracasseries douanières/policières (multiples barrières, taxes illégales), état routes/pistes, insécurité (Boko Haram, coupeurs de route), absence de monnaie commune d'échange fluide (FCFA vs Naira vs Franc CFA BEAC/CEMAC), informalité ; Pistes : Guichets uniques frontières, modernisation postes (Kribi, Garoua-Boulaï, Kyé-Ossi), application effective textes CEMAC/ZLECAf, sécurisation axes.
\end{enumerate}
\textbf{Réflexe de copie :} Une étude de cas gagnante articule toujours \emph{Localisation -- Flux -- Impacts -- Difficultés}, avec au moins \textbf{deux exemples de marchés cités} parmi la liste du programme (Banki, Mbaïboum, Abam-Minko, Kyé-Ossi, Ekondo-Titi).
\end{methode}

\begin{exoresolu}[Étude de cas : les marchés frontaliers du Cameroun]
\textbf{Énoncé (type Bac, Dossier 5).} « De Banki à Kyé-Ossi, les marchés frontaliers drainent chaque semaine des milliers de commerçants et des tonnes de marchandises, pour l'essentiel hors des circuits officiels. » À l'aide de deux exemples de marchés frontaliers, présentez les flux de marchandises, montrez l'impact économique et social de ces marchés, puis dégagez les difficultés des échanges transfrontaliers au Cameroun.
\tcblower
\textbf{Solution détaillée.}

\emph{1. Localisation et exemples choisis.}
\textbf{Banki} (département du Mayo-Sava, Extrême-Nord, frontière Nigeria, sur l'axe Mora--Maiduguri, à \SI{5}{km} de la route nationale) est l'un des plus anciens et vast marchés transfrontaliers du pays. \textbf{Kyé-Ossi} (département de la Vallée-du-Ntem, Région du Sud, triple frontière Cameroun--Gabon--Guinée équatoriale, carrefour routier Ebolowa--Oyem--Ebebiyin) est un marché hebdomadaire très fréquenté. On peut citer aussi \textbf{Mbaïboum} (Est, frontière RCA, axe Garoua-Boulaï--Bangui), \textbf{Abam-Minko} (Sud, frontière Gabon, axe Ebolowa--Bitam) et \textbf{Ekondo-Titi} (Sud-Ouest, zone d'échanges maritimes/fluviaux avec le Nigeria, creek).

\emph{2. Les flux de marchandises (bidirectionnels, largement informels).}
\begin{itemize}
\item \textbf{Du Cameroun vers les voisins :} Denrées vivrières (oignons, niébé, pommes de terre du Nord ; bananes, manioc, plantains du Sud/Centre), Produits agroalimentaires industriels (savons Azur, huiles Dinor, boissons Brasseries/Guinness, pâtes, sucre), Ciment (Cimencam), Tissus, friperie, médicaments génériques, carburant (réexport formel/informel).
\item \textbf{Des voisins vers le Cameroun :} Carburant nigérian (super, gasoil, \emph{« kpayo »}, moins cher que prix pompe camerounais, contrebande massive), Bétail sur pied (Nigeria, Tchad, RCA), Poisson fumé/séché (Lac Tchad, mer), Produits manufacturés réexportés depuis ports nigérians (Lagos, Calabar) ou gabonais (Libreville) : électroménagers, motos, pièces détachées, textiles, riz, huile végétale.
\end{itemize}
Ces flux sont animés par des réseaux de commerçants (grossistes, détaillants \emph{bayam-sellam}), des transporteurs (motos \emph{benskin}, camions, pirogues à Ekondo-Titi) et des changeurs de monnaie (FCFA / Naira / Dollar / Euro).

\emph{3. Impacts économiques et sociaux.}
\textbf{Économiques :} Ces marchés sont le \emph{poumon vital} des villes frontalières (Banki a longtemps approvisionné tout l'Extrême-Nord en carburant et denrées). Ils créent des milliers d'emplois directs/indirects (vente, transport, gardiennage, restauration, change). Ils permettent l'approvisionnement à bas prix des populations (carburant, riz, poisson). Ils génèrent des revenus en devises (Naira, Dollar) pour les opérateurs camerounais. \textbf{Sociaux :} Ils entretiennent le brassage des populations (ethnies à cheval sur les frontières : Peuls, Kanuri, Mafa, Massa, Fang, Bassa) et la solidarité transfrontalière (mariages, parentés, entraide). \textbf{Revers :} L'État perd d'importantes recettes douanières (droits, TVA) ; la contrebande (carburant, médicaments falsifiés, armes légères, produits prohibés) prospère ; l'insécurité (exactions Boko Haram autour de Banki/Kerawa depuis 2014, coupeurs de route sur axes Est/Sud) a durement frappé certains marchés, provoquant déplacements et fermetures temporaires.

\emph{4. Difficultés structurelles des échanges transfrontaliers.}
\begin{itemize}
\item \textbf{Tracasseries institutionnelles :} Multiplication des postes de contrôle (douane, police, gendarmerie, eaux/forêts, santé, brigade mobile), taxes illégales, rackets, lenteur procédures (certificats d'origine, licences).
\item \textbf{Déficit infrastructurel :} Pistes impraticables en saison des pluies (Mbaïboum, Abam-Minko), absence de stockage frigorifique, de zones de déchargement aménagées, de guichets uniques.
\item \textbf{Barrières monétaires et réglementaires :} Difficultés de change FCFA/Naira (taux parallèle), non-reconnaissance mutuelle des certificats phytosanitaires/vétérinaires, application incomplète du TEC et de la libre circulation CEMAC.
\item \textbf{Insécurité :} Terrorisme (Extrême-Nord), banditisme (Est, Sud), piraterie maritime (Ekondo-Titi).
\end{itemize}

\emph{Conclusion.} Les marchés frontaliers (Banki, Kyé-Ossi, Mbaïboum, Abam-Minko, Ekondo-Titi...) sont des lieux d'\textbf{intégration régionale « par le bas »}, dynamiques, résilients, mais largement informels. Leur modernisation (guichets uniques frontaliers, numérisation procédures, sécurisation axes, infrastructures de marché) transformerait une richesse clandestine en moteur officiel du développement local et de l'intégration CEMAC/ZLECAf.
\end{exoresolu}

\begin{exercice}[Entraînement Dossier 5 : Sujet type Bac]
\textbf{Sujet.} À partir de la carte ci-dessous (carte mentale des flux au marché de Kyé-Ossi) et de vos connaissances, présentez l'organisation spatiale de ce marché frontalier et analysez les défis de son intégration dans le commerce formel de la CEMAC. (15 lignes attendues).
\begin{center}
\textit{[Carte mentale à fournir au candidat : Kyé-Ossi au centre. Flèches entrantes : Gabon (bois, carburant, produits manufacturés), Guinée Équatoriale (produits alimentaires, bétail), Cameroun (vivriers, boissons, ciment, friperie). Flèches sortantes vers les deux pays. Symboles : poste frontière, piste, rivière Ntem, zone de change informel.]}
\end{center}
\end{exercice}

\begin{corrige}[Exercice Dossier 5 : Kyé-Ossi]
\emph{Organisation spatiale :} Kyé-Ossi est un marché-carrefour \textbf{tri-national} (Cameroun-Gabon-Guinée Équatoriale), situé sur la rivière Ntem (frontière naturelle), à la convergence de trois axes routiers (Ebolowa-Kyé-Ossi-Oyem/Gabon ; Kyé-Ossi-Ebebiyin/Guinée Eq.). L'espace de marché déborde du poste frontière officiel : il s'étend sur les places publiques, les abords de la piste, et même de l'autre côté de la rivière (pirogues). Une \textbf{zone de change informel} (bureaux de change mobiles) assure la conversion FCFA / Franc CFA BEAC (même monnaie mais billets différents parfois refusés) / Naira / Dollar.

\emph{Flux :} Échanges de complémentarité (vivriers camerounais vs bois/manufacturés gabonais ; bétail/poisson guinéens vs produits agroalimentaires camerounais). Dominance du commerce de proximité, à pied, moto, pirogue.

\emph{Défis de l'intégration formelle :} 1) \textbf{Poste frontière unique absent :} Trois postes distincts, procédures lentes, tracasseries. 2) \textbf{Monnaie :} Bien que zone FCFA, les billets BEAC (Cameroun, Gabon, Guinée Eq.) ne circulent pas librement partout ; le Naira nigérian sert souvent de référence. 3) \textbf{Infrastructures :} Pistes boueuses, pas de halle couverte, pas de froid pour poisson/viande. 4) \textbf{Réglementation :} Certificats d'origine CEMAC peu utilisés, barrières phytosanitaires. 5) \textbf{Sécurité :} Contrôles d'identité multiples, risques de traite personnes.
\emph{Pistes :} Poste de contrôle juxtaposé (PCJ) unique, halle moderne, route bitumée Ebolowa-Kyé-Ossi-Oyem (en cours), sensibilisation opérateurs sur avantages ZLECAf (exonération droits de douane sur produits originaires).
\end{corrige}

\newpage

\coursec{Exercices}

\courssub{Je vérifie mes connaissances}

\begin{exercice}[QCM : les notions clés des échanges extérieurs]
Pour chaque question, entoure la (ou les) bonne(s) réponse(s).
\begin{enumerate}
\item La balance commerciale est :
\begin{enumerate}
\item la différence entre la valeur des exportations et celle des importations ;
\item la différence entre les recettes et les dépenses de l'État ;
\item le rapport entre les prix à l'exportation et les prix à l'importation.
\end{enumerate}
\item Les termes de l'échange se détériorent lorsque :
\begin{enumerate}
\item les prix des exportations augmentent plus vite que ceux des importations ;
\item les prix des importations augmentent plus vite que ceux des exportations ;
\item le volume des exportations diminue.
\end{enumerate}
\item La CEMAC regroupe :
\begin{enumerate}
\item six pays d'Afrique centrale ;
\item huit pays d'Afrique de l'Ouest ;
\item le Cameroun, le Gabon, le Congo, la Guinée équatoriale, le Tchad et la Centrafrique.
\end{enumerate}
\item Le premier produit d'exportation du Cameroun en valeur est généralement :
\begin{enumerate}
\item le cacao ; \item le pétrole brut ; \item le café.
\end{enumerate}
\end{enumerate}
\end{exercice}

\begin{exercice}[Vrai ou faux ? Justifie ta réponse]
Réponds par vrai ou faux, puis justifie chaque réponse en deux ou trois lignes.
\begin{enumerate}
\item Le Cameroun échange davantage avec les pays de la CEMAC qu'avec l'Union européenne et la Chine.
\item La détérioration des termes de l'échange signifie que le Cameroun doit exporter davantage de produits pour importer la même quantité de biens.
\item Le marché frontalier de Kyé-Ossi se situe à la frontière avec le Nigéria.
\item Le Cameroun exporte surtout des produits manufacturés à forte valeur ajoutée.
\end{enumerate}
\end{exercice}

\begin{exercice}[Définitions expresses]
Définis en deux ou trois lignes chacune des notions suivantes : \cle{balance commerciale} ; \cle{termes de l'échange} ; \cle{marché frontalier} ; \cle{CEMAC}.
\end{exercice}

\begin{exercice}[Calculs directs sur la balance commerciale]
En 2021 (ordres de grandeur, source INS), le Cameroun a exporté pour environ \num{2500} milliards de FCFA et importé pour environ \num{4300} milliards de FCFA.
\begin{enumerate}
\item Calcule le solde de la balance commerciale du Cameroun en 2021.
\item La balance commerciale est-elle excédentaire ou déficitaire ? Justifie.
\item Calcule le taux de couverture des importations par les exportations, selon la formule :
\[ \text{Taux de couverture} = \frac{\text{valeur des exportations}}{\text{valeur des importations}} \times 100 \]
\item Interprète ce taux en une phrase.
\end{enumerate}
\end{exercice}

\courssub{Je m'entraîne}

\begin{exercice}[Les produits échangés]
\begin{enumerate}
\item Cite quatre produits exportés par le Cameroun et quatre produits importés.
\item Classe ces produits en trois catégories : produits primaires agricoles, produits énergétiques et miniers, produits manufacturés.
\item Que remarques-tu sur la place des produits manufacturés dans les exportations et dans les importations du Cameroun ?
\end{enumerate}
\end{exercice}

\begin{exercice}[Cartographie : la CEMAC et les marchés frontaliers]
Sur une carte muette de l'Afrique centrale :
\begin{enumerate}
\item Localise et nomme les six États membres de la CEMAC.
\item Place les marchés frontaliers suivants : Banki, Mbaïboum, Abam-Minko, Kyé-Ossi, Ekondo-Titi.
\item À l'aide de flèches, représente deux flux de marchandises transfrontaliers que tu connais (précise le sens du flux et la nature des produits).
\item Rédige une légende complète et donne un titre à ta carte.
\end{enumerate}
\end{exercice}

\begin{exercice}[Diagramme en barres : clients et fournisseurs du Cameroun]
Voici les principaux partenaires commerciaux du Cameroun (ordres de grandeur en \% du total des échanges, données INS/années récentes) : à l'exportation, la Chine et les Pays-Bas environ 15 \% chacun, l'Inde et l'Italie environ 8 \% chacun ; à l'importation, la Chine environ 20 \%, la France environ 10 \%, la Belgique et l'Inde environ 6 \% chacune.
\begin{enumerate}
\item Construis un diagramme en barres comparant les principaux clients et les principaux fournisseurs du Cameroun (n'oublie ni le titre, ni les axes, ni la légende).
\item Montre, à partir de ton diagramme, que les échanges du Cameroun sont \cle{extravertis}, c'est-à-dire tournés vers l'extérieur de l'Afrique.
\item Calcule la part cumulée de la Chine dans les exportations et les importations du Cameroun. Commentaire ?
\end{enumerate}
\end{exercice}

\begin{exercice}[Commentaire de photographie : le marché de Kyé-Ossi]
On te donne la description suivante : « À Kyé-Ossi, à la triple frontière Cameroun--Gabon--Guinée équatoriale, des centaines de commerçants exposent vivres, tissus, quincaillerie et produits manufacturés ; des taxis-motos et des pick-up chargés à ras bord franchissent les postes de contrôle ; des files d'attente se forment devant les guichets des douanes et de la police. »
\begin{enumerate}
\item Décris les flux observés sur ce marché frontalier (nature des produits, acteurs, directions).
\item Dégage deux impacts économiques et deux impacts sociaux du commerce transfrontalier pour les populations riveraines.
\item Relève dans le texte un indice des difficultés des échanges transfrontaliers et explique-le.
\end{enumerate}
\end{exercice}

\courssub{Je me prépare au Bac}

\begin{exercice}[Type Bac : exploitation de données chiffrées sur la balance commerciale]
Le tableau ci-dessous présente les exportations et les importations du Cameroun sur cinq ans (valeurs fictives mais réalistes, en milliards de FCFA).

\begin{center}
\begin{tabularx}{\linewidth}{|l|X|X|X|X|X|}
\hline
\rowcolor{popblueL} \textbf{Année} & \textbf{Année 1} & \textbf{Année 2} & \textbf{Année 3} & \textbf{Année 4} & \textbf{Année 5} \\
\hline
Exportations & \num{2100} & \num{2300} & \num{2000} & \num{2500} & \num{2600} \\
\hline
Importations & \num{2800} & \num{3100} & \num{3300} & \num{3800} & \num{4300} \\
\hline
\end{tabularx}
\end{center}

\begin{enumerate}
\item Définis : balance commerciale ; solde commercial.
\item Calcule le solde commercial du Cameroun pour chacune des cinq années et présente tes résultats dans un tableau.
\item Calcule le taux de couverture des importations par les exportations pour l'année 1 et pour l'année 5.
\item Construis un graphique (courbes ou barres groupées) représentant l'évolution des exportations, des importations et du solde commercial sur la période.
\item Rédige un commentaire de dix à quinze lignes : tu montreras l'évolution du déficit, tu en proposeras deux explications (structure des produits échangés, détérioration des termes de l'échange) et tu suggéreras une solution.
\end{enumerate}
\end{exercice}

\begin{exercice}[Type Bac : question de cours et étude de document sur les marchés frontaliers]
\textbf{Document.} « Aux frontières du Cameroun, des marchés comme Banki (frontière nigériane), Mbaïboum (frontière centrafricaine), Abam-Minko et Kyé-Ossi (frontières gabonaise et équato-guinéenne) ou Ekondo-Titi (frontière nigériane) drainent chaque semaine des milliers de commerçants. On y échange denrées alimentaires, bétail, carburant, tissus et produits manufacturés importés. Mais les tracasseries douanières et policières, la multiplicité des contrôles, la faiblesse des infrastructures routières et l'insécurité sur certains axes freinent ces échanges, pourtant essentiels au ravitaillement des zones frontalières et à l'intégration régionale. »

\begin{enumerate}
\item Question de cours : définis \cle{marché frontalier} et cite quatre marchés frontaliers du Cameroun en précisant pour chacun le pays voisin concerné.
\item À partir du document, décris les flux de marchandises qui transitent par les marchés frontaliers (nature des produits, sens des échanges).
\item Dégage du document trois difficultés des échanges transfrontaliers et explique chacune en une ou deux phrases.
\item Analyse en une dizaine de lignes les impacts économiques et sociaux des marchés frontaliers sur les populations riveraines et sur l'intégration régionale en Afrique centrale.
\end{enumerate}
\end{exercice}

\begin{exercice}[Type Bac : dissertation guidée]
\textbf{Sujet :} « Les échanges extérieurs du Cameroun : atouts, faiblesses et perspectives. »

Tu rédigeras une dissertation géographique complète en respectant la méthode suivante :
\begin{enumerate}
\item Une introduction comportant une phrase d'accroche (situation concrète : un produit exporté du port de Douala ou un produit importé dans un marché camerounais), la définition des termes clés du sujet, une problématique claire et l'annonce du plan.
\item Un développement en deux grandes parties argumentées :
\begin{itemize}
\item première partie : les atouts et les points forts des échanges extérieurs du Cameroun (diversité des produits exportés, position de carrefour en Afrique centrale, port de Douala, appartenance à la CEMAC, marchés frontaliers dynamiques) ;
\item deuxième partie : les faiblesses et les problèmes (déficit de la balance commerciale, détérioration des termes de l'échange, faible part des produits manufacturés, extraversion des échanges, tracasseries aux frontières), suivie des perspectives d'amélioration (transformation locale, intégration régionale, ZLECAf).
\end{itemize}
\item Une conclusion qui répond à la problématique et s'ouvre sur l'enjeu de l'industrialisation du Cameroun à l'horizon 2035.
\end{enumerate}
\end{exercice}

\newpage

\coursec{Corrigés des exercices}

\coursec{Exercices d'application et corrigés détaillés : Les échanges du Cameroun avec l'extérieur}

\begin{corrige}[Exercice 1 — QCM : les notions clés des échanges extérieurs]
\begin{enumerate}
\item \textbf{Bonne réponse : a.} La balance commerciale est la différence entre la valeur des exportations et celle des importations de biens d'un pays sur une période donnée. La réponse b désigne le solde budgétaire de l'État, et la réponse c correspond aux termes de l'échange : attention à ne pas confondre ces trois notions.
\item \textbf{Bonne réponse : b.} Les termes de l'échange se détériorent lorsque les prix des importations augmentent plus vite que ceux des exportations : le pays doit alors vendre davantage pour acheter la même quantité de biens. La réponse a décrit une \emph{amélioration} des termes de l'échange ; la réponse c concerne les volumes, pas les prix relatifs.
\item \textbf{Bonnes réponses : a et c.} La CEMAC (Communauté économique et monétaire de l'Afrique centrale) regroupe six pays : le Cameroun, le Gabon, le Congo, la Guinée équatoriale, le Tchad et la Centrafrique. La réponse b évoque la CEDEAO (Afrique de l'Ouest), à ne pas confondre.
\item \textbf{Bonne réponse : b.} Le pétrole brut est généralement le premier produit d'exportation du Cameroun en valeur (avec le bois et le cacao selon les années). Le cacao et le café sont des produits de rente importants, mais leur valeur reste inférieure à celle du pétrole.
\end{enumerate}
\textbf{Points de vigilance :} au Bac, une confusion entre balance commerciale et termes de l'échange coûte cher ; retiens : balance = \emph{différence de valeurs}, termes de l'échange = \emph{rapport de prix}.
\end{corrige}

\begin{corrige}[Exercice 2 — Vrai ou faux ? Justifie ta réponse]
\begin{enumerate}
\item \textbf{Faux.} Les échanges du Cameroun sont \cle{extravertis} : l'essentiel du commerce se fait avec l'Union européenne et la Chine, qui figurent parmi ses premiers clients et fournisseurs. Les échanges intra-CEMAC restent faibles (de l'ordre de quelques pourcents du total), malgré l'appartenance du Cameroun à cette communauté.
\item \textbf{Vrai.} La détérioration des termes de l'échange signifie que les prix des exportations progressent moins vite que ceux des importations : pour acheter la même quantité de biens importés (machines, carburant, médicaments), le Cameroun doit exporter un volume croissant de matières premières (cacao, pétrole, bois).
\item \textbf{Faux.} Kyé-Ossi se situe dans la région du Sud, à la triple frontière Cameroun--Gabon--Guinée équatoriale. Le grand marché frontalier avec le Nigéria est Banki (Extrême-Nord), ainsi qu'Ekondo-Titi (Sud-Ouest).
\item \textbf{Faux.} Le Cameroun exporte surtout des produits primaires bruts ou peu transformés (pétrole brut, bois en grumes, cacao, café, coton, bananes). Ce sont au contraire les \emph{importations} qui sont dominées par les produits manufacturés (machines, véhicules, produits pharmaceutiques), ce qui explique le déficit commercial.
\end{enumerate}
\textbf{Points de vigilance :} une justification doit toujours mobiliser un fait précis (un partenaire, un lieu, un produit) ; « c'est faux parce que c'est faux » ne vaut aucun point.
\end{corrige}

\begin{corrige}[Exercice 3 — Définitions expresses]
\begin{itemize}
\item \textbf{Balance commerciale :} différence entre la valeur des exportations et celle des importations de biens d'un pays sur une période donnée. Elle est excédentaire si les exportations dépassent les importations, déficitaire dans le cas contraire.
\item \textbf{Termes de l'échange :} rapport entre l'indice des prix des exportations et l'indice des prix des importations, multiplié par 100. Ils se détériorent quand les prix des importations augmentent plus vite que ceux des exportations.
\item \textbf{Marché frontalier :} marché situé à proximité d'une frontière internationale, où s'échangent des marchandises entre commerçants et consommateurs de deux ou trois pays voisins (ex. : Banki, Kyé-Ossi). Il joue un rôle de ravitaillement et d'intégration régionale.
\item \textbf{CEMAC :} Communauté économique et monétaire de l'Afrique centrale, organisation d'intégration régionale regroupant six États (Cameroun, Gabon, Congo, Guinée équatoriale, Tchad, Centrafrique), dotée d'une monnaie commune, le franc CFA (BEAC).
\end{itemize}
\textbf{Points de vigilance :} une définition au Bac comporte l'idée générale + un élément de précision (formule, exemple ou membre) ; évite les définitions circulaires (« la balance commerciale est la balance des échanges commerciaux »).
\end{corrige}

\begin{corrige}[Exercice 4 — Calculs directs sur la balance commerciale]
\begin{enumerate}
\item Solde commercial = exportations $-$ importations :
\[ \num{2500} - \num{4300} = \num{-1800} \text{ milliards de FCFA.} \]
\item La balance commerciale est \textbf{déficitaire}, car la valeur des importations (\num{4300} milliards) est supérieure à celle des exportations (\num{2500} milliards) : le solde est négatif.
\item Taux de couverture :
\[ \frac{\num{2500}}{\num{4300}} \times 100 \approx 58,1\,\%. \]
\item Interprétation : en 2021, les exportations du Cameroun ne couvrent qu'environ 58 \% de la valeur des importations ; autrement dit, le pays doit financer par d'autres moyens (endettement, aide, envois de la diaspora) près de 42 \% de ses achats à l'étranger.
\end{enumerate}
\textbf{Points de vigilance :} n'oublie pas le signe « $-$ » du solde et le symbole \% du taux ; arrondis proprement et précise toujours l'unité (milliards de FCFA).
\end{corrige}

\begin{corrige}[Exercice 5 — Les produits échangés]
\begin{enumerate}
\item \textbf{Exportations (exemples) :} pétrole brut, cacao, bois, café (aussi : coton, bananes, aluminium, caoutchouc). \textbf{Importations (exemples) :} carburants raffinés, machines et véhicules, produits pharmaceutiques, riz et blé (aussi : matériel électrique, ciment, produits chimiques).
\item \textbf{Classement :}
\begin{center}
\begin{tabularx}{\linewidth}{|l|>{\centering\arraybackslash}X|>{\centering\arraybackslash}X|}
\hline
\rowcolor{popblueL} \textbf{Catégorie} & \textbf{Exportations} & \textbf{Importations} \\
\hline
Produits primaires agricoles & cacao, café, coton, bananes & riz, blé, poisson congelé \\
\hline
Produits énergétiques et miniers & pétrole brut, aluminium & carburants raffinés \\
\hline
Produits manufacturés & (très peu : savons, bois scié) & machines, véhicules, médicaments \\
\hline
\end{tabularx}
\end{center}
\item \textbf{Remarque :} les produits manufacturés sont quasi absents des exportations mais dominent les importations. Le Cameroun vend des produits bruts à faible valeur ajoutée et achète des produits transformés chers : c'est la structure typique d'un commerce \cle{extraverti} et déséquilibré, source du déficit commercial et de la détérioration des termes de l'échange.
\end{enumerate}
\textbf{Points de vigilance :} le paradoxe pétrole brut exporté / carburant raffiné importé (faute de raffinage suffisant, la SONARA étant à capacité limitée) est un excellent argument à mobiliser au Bac.
\end{corrige}

\begin{corrige}[Exercice 6 — Cartographie : la CEMAC et les marchés frontaliers]
\textbf{Corrigé de localisation (à reporter sur la carte muette) :}
\begin{enumerate}
\item \textbf{Les six États de la CEMAC :} Cameroun (à l'ouest, sur le golfe de Guinée), Gabon et Guinée équatoriale (au sud-ouest, sur la côte), Congo (au sud-est), Tchad (au nord-est), Centrafrique (à l'est).
\item \textbf{Marchés frontaliers :} Banki (Extrême-Nord, frontière avec le Nigéria, près du lac Tchad) ; Mbaïboum (Est, frontière avec la Centrafrique) ; Abam-Minko et Kyé-Ossi (Sud, frontières avec le Gabon et la Guinée équatoriale) ; Ekondo-Titi (Sud-Ouest, frontière avec le Nigéria).
\item \textbf{Flux à représenter (exemples) :} flèche Nigéria $\rightarrow$ Cameroun à Banki (bétail, carburant, produits manufacturés) ; flèche Cameroun $\rightarrow$ Gabon/Guinée équatoriale à Kyé-Ossi (denrées alimentaires : plantain, manioc, légumes) ; flèche Cameroun $\rightarrow$ Centrafrique (produits manufacturés et vivres via Garoua-Boulaï et Mbaïboum).
\item \textbf{Titre possible :} « La CEMAC et les principaux marchés frontaliers du Cameroun ». \textbf{Légende :} 1. États membres de la CEMAC (zone colorée) ; 2. marché frontalier (symbole) ; 3. flux de marchandises (flèche, sens et nature précisés) ; 4. capitale (point nommé). Ajoute l'orientation (flèche du nord) et une échelle indicative.
\end{enumerate}
\textbf{Points de vigilance :} une carte sans titre, sans légende numérotée ou sans nord est sanctionnée ; chaque élément du croquis doit figurer dans la légende, et rien dans la légende qui ne soit sur la carte.
\end{corrige}

\begin{corrige}[Exercice 7 — Diagramme en barres : clients et fournisseurs du Cameroun]
\begin{enumerate}
\item \textbf{Construction :} axe horizontal = partenaires (Chine, Pays-Bas, Inde, Italie, France, Belgique) ; axe vertical = part en \% du total des échanges (de 0 à 25 \%). Deux séries de barres groupées : « clients » (exportations) et « fournisseurs » (importations), avec deux couleurs distinctes et une légende. Titre : « Principaux clients et fournisseurs du Cameroun (en \% du total des échanges, ordres de grandeur) ».

\begin{popfigure}
\begin{center}
\begin{tikzpicture}
\begin{axis}[
  ybar, bar width=8pt,
  width=0.95\linewidth, height=6.5cm,
  ymin=0, ymax=25,
  ylabel={Part en \%},
  symbolic x coords={Chine,Pays-Bas,Inde,Italie,France,Belgique},
  xtick=data, x tick label style={font=\small},
  legend style={at={(0.5,1.02)},anchor=south,legend columns=2,font=\small},
  nodes near coords, every node near coord/.append style={font=\scriptsize},
]
\addplot[fill=popblue] coordinates {(Chine,15) (Pays-Bas,15) (Inde,8) (Italie,8) (France,0) (Belgique,0)};
\addplot[fill=poporange] coordinates {(Chine,20) (Pays-Bas,0) (Inde,6) (Italie,0) (France,10) (Belgique,6)};
\legend{Clients (exportations), Fournisseurs (importations)}
\end{axis}
\end{tikzpicture}
\end{center}
\legende{Principaux clients et fournisseurs du Cameroun, en \% du total des échanges (ordres de grandeur, source INS). Les barres à 0 signifient que le partenaire n'est pas cité parmi les principaux pour ce flux.}
\end{popfigure}

\item \textbf{Extraversion :} tous les partenaires représentés sont extra-africains (Asie : Chine, Inde ; Europe : Pays-Bas, Italie, France, Belgique). Aucun pays africain, et notamment aucun membre de la CEMAC, n'apparaît parmi les principaux clients ou fournisseurs : les échanges du Cameroun sont donc tournés vers l'extérieur du continent, ce qui définit un commerce \cle{extraverti}.
\item \textbf{Part cumulée de la Chine :} $15 + 20 = 35\,\%$. La Chine est à la fois le premier client (environ 15 \% des exportations) et le premier fournisseur (environ 20 \% des importations) du Cameroun : elle pèse à elle seule environ un tiers des échanges cités ici, ce qui traduit une dépendance croissante vis-à-vis d'un partenaire unique.
\end{enumerate}
\textbf{Points de vigilance :} un diagramme sans titre, sans axes gradués ou sans légende est incomplet ; dans le commentaire, cite toujours des chiffres précis pour appuyer le qualificatif « extraverti ».
\end{corrige}

\begin{corrige}[Exercice 8 — Commentaire de photographie : le marché de Kyé-Ossi]
\begin{enumerate}
\item \textbf{Flux observés :} les produits échangés sont variés : vivres et denrées alimentaires (en provenance surtout du Cameroun), tissus, quincaillerie et produits manufacturés (souvent réexportés vers le Gabon et la Guinée équatoriale). Les acteurs sont des centaines de commerçants (détaillants, grossistes, transporteurs) ; les flux sont bidirectionnels entre les trois pays de la triple frontière, comme le montrent les taxis-motos et pick-up chargés qui franchissent les postes de contrôle dans les deux sens.
\item \textbf{Impacts économiques :} (1) création de revenus et d'emplois (commerce, transport, manutention) pour les populations riveraines ; (2) ravitaillement des zones frontalières en produits variés à prix accessibles et recettes douanières pour l'État. \textbf{Impacts sociaux :} (1) brassage des populations et renforcement des liens de voisinage entre les trois pays (intégration régionale vécue au quotidien) ; (2) dynamisation et urbanisation des localités frontalières, mais aussi parfois congestion, insalubrité et tensions autour des contrôles.
\item \textbf{Indice de difficulté :} « des files d'attente se forment devant les guichets des douanes et de la police ». Cet indice traduit la multiplicité des contrôles et les tracasseries administratives : perte de temps pour les commerçants, coûts supplémentaires (droits, parfois racket), qui renchérissent les marchandises et freinent les échanges.
\end{enumerate}
\textbf{Points de vigilance :} dans un commentaire de photographie, chaque affirmation doit s'appuyer sur un élément visible (ou décrit) du document ; distingue bien impacts économiques (revenus, emplois, prix) et impacts sociaux (populations, modes de vie, relations).
\end{corrige}

\begin{corrige}[Exercice 9 — Type Bac : exploitation de données chiffrées sur la balance commerciale]
\begin{enumerate}
\item \textbf{Définitions :} la \cle{balance commerciale} est le document comptable qui enregistre la valeur des exportations et des importations de biens d'un pays ; le \cle{solde commercial} est la différence entre la valeur des exportations et celle des importations (excédent si positif, déficit si négatif).
\item \textbf{Soldes commerciaux} (solde = exportations $-$ importations) :
\begin{center}
\begin{tabularx}{\linewidth}{|l|>{\centering\arraybackslash}X|>{\centering\arraybackslash}X|>{\centering\arraybackslash}X|>{\centering\arraybackslash}X|>{\centering\arraybackslash}X|}
\hline
\rowcolor{popblueL} \textbf{Année} & \textbf{1} & \textbf{2} & \textbf{3} & \textbf{4} & \textbf{5} \\
\hline
Exportations & \num{2100} & \num{2300} & \num{2000} & \num{2500} & \num{2600} \\
\hline
Importations & \num{2800} & \num{3100} & \num{3300} & \num{3800} & \num{4300} \\
\hline
Solde & \num{-700} & \num{-800} & \num{-1300} & \num{-1300} & \num{-1700} \\
\hline
\end{tabularx}
\end{center}
Vérification : $2100-2800=-700$ ; $2300-3100=-800$ ; $2000-3300=-1300$ ; $2500-3800=-1300$ ; $2600-4300=-1700$ (milliards de FCFA).
\item \textbf{Taux de couverture :}
\[ \text{Année 1 : } \frac{2100}{2800}\times 100 = 75\,\% \qquad ; \qquad \text{Année 5 : } \frac{2600}{4300}\times 100 \approx 60,5\,\%. \]
\item \textbf{Graphique :} barres groupées (exportations en bleu, importations en orange) et courbe du solde en rouge sous l'axe, ou trois courbes. Titre : « Évolution des échanges extérieurs du Cameroun (années 1 à 5) » ; axe vertical gradué en milliards de FCFA, légende complète.
\item \textbf{Commentaire (10 à 15 lignes) :} sur la période, le commerce extérieur du Cameroun reste structurellement déficitaire et le déficit s'aggrave : il passe de $-700$ milliards de FCFA en année 1 à $-\num{1700}$ milliards en année 5, soit une multiplication par près de 2,4. Les exportations progressent faiblement et irrégulièrement (de \num{2100} à \num{2600} milliards, avec un recul en année 3), tandis que les importations augmentent fortement et continûment (de \num{2800} à \num{4300} milliards, soit $+54\,\%$). Le taux de couverture se dégrade en conséquence, de 75 \% à environ 60 \%. Deux explications : d'une part, la structure des produits échangés — le Cameroun exporte des matières premières brutes à faible valeur ajoutée (pétrole, cacao, bois) et importe des produits manufacturés coûteux (machines, véhicules, carburant raffiné) ; d'autre part, la détérioration des termes de l'échange — les prix des produits importés augmentent plus vite que ceux des matières premières exportées, ce qui oblige à exporter davantage pour importer autant. Pour inverser la tendance, le Cameroun doit transformer localement ses matières premières (chocolaterie, sciage, raffinage) afin d'exporter des produits à plus forte valeur ajoutée et de réduire ses importations de biens manufacturés.
\end{enumerate}
\textbf{Barème indicatif (sur 20) :} définitions 2 pts ; calculs des soldes et tableau 4 pts ; taux de couverture 3 pts ; graphique (titre, axes, légende, exactitude) 4 pts ; commentaire (évolution 2 pts, explications 3 pts, solution 2 pts) 7 pts.
\textbf{Points de vigilance :} cite les chiffres calculés dans le commentaire ; un commentaire qui paraphrase le tableau sans expliquer ni proposer de solution plafonne à la moitié des points.
\end{corrige}

\begin{corrige}[Exercice 10 — Type Bac : question de cours et étude de document sur les marchés frontaliers]
\begin{enumerate}
\item \textbf{Question de cours :} un \cle{marché frontalier} est un marché situé à proximité d'une frontière internationale, où s'échangent des marchandises entre les populations de deux ou trois pays voisins. Quatre exemples : Banki (Nigéria), Mbaïboum (Centrafrique), Abam-Minko (Gabon), Kyé-Ossi (Gabon et Guinée équatoriale) ; on peut ajouter Ekondo-Titi (Nigéria).
\item \textbf{Flux de marchandises :} le document mentionne des denrées alimentaires et du bétail (qui circulent notamment du Nigéria et du Cameroun vers les pays voisins), du carburant (souvent du Nigéria vers le Cameroun), des tissus et des produits manufacturés importés (réexportés du Cameroun vers la Centrafrique, le Gabon et la Guinée équatoriale). Les échanges sont donc bidirectionnels et drainent chaque semaine des milliers de commerçants.
\item \textbf{Trois difficultés :} (1) les tracasseries douanières et policières et la multiplicité des contrôles : elles allongent les délais, multiplient les paiements informels et renchérissent les produits ; (2) la faiblesse des infrastructures routières : les pistes dégradées ralentissent le transport, abîment les marchandises et isolent les marchés pendant la saison des pluies ; (3) l'insécurité sur certains axes (coupeurs de route, conflits, groupes armés dans l'Extrême-Nord) : elle décourage commerçants et transporteurs et réduit les flux.
\item \textbf{Analyse des impacts (une dizaine de lignes) :} les marchés frontaliers jouent un rôle économique majeur : ils assurent le ravitaillement des zones frontalières en vivres et produits manufacturés, créent des emplois (commerce, transport, change, restauration) et des revenus pour des milliers de ménages, et procurent des recettes douanières à l'État. Sur le plan social, ils favorisent le brassage des populations, les mariages et les solidarités transfrontalières, et dynamisent des localités autrefois marginales. À l'échelle régionale, ces échanges, même informels, concrétisent l'intégration prônée par la CEMAC : ils font circuler biens, monnaies et personnes entre pays voisins. Cependant, leur caractère largement informel échappe aux statistiques et prive l'État de recettes, tandis que les tracasseries et l'insécurité limitent leur potentiel. Sécuriser les axes, simplifier les contrôles et moderniser les infrastructures renforceraient leur contribution à l'intégration de l'Afrique centrale.
\end{enumerate}
\textbf{Barème indicatif (sur 20) :} question de cours 4 pts ; description des flux 4 pts ; difficultés expliquées 6 pts ; analyse des impacts 6 pts.
\textbf{Points de vigilance :} distingue clairement ce qui est tiré du document (cite-le) et tes connaissances de cours ; chaque difficulté relevée doit être \emph{expliquée}, pas seulement recopiée.
\end{corrige}

\begin{corrige}[Exercice 11 — Type Bac : dissertation guidée]
\textbf{Proposition de copie (plan détaillé rédigé).}

\textbf{Introduction.} \emph{Accroche :} chaque semaine, des dizaines de conteneurs de cacao et de bois quittent le port de Douala vers l'Europe et l'Asie, tandis que débarquent sur les mêmes quais voitures, machines et médicaments qui inonderont les marchés de Mokolo ou de Mvog-Mbi. \emph{Définitions :} les échanges extérieurs désignent l'ensemble des exportations et importations de biens (et de services) entre le Cameroun et le reste du monde ; leurs atouts sont les points forts qui les soutiennent, leurs faiblesses les facteurs qui les limitent. \emph{Problématique :} dans quelle mesure les échanges extérieurs du Cameroun constituent-ils un levier de développement, et quels obstacles en limitent les bénéfices ? \emph{Annonce du plan :} nous verrons d'abord les atouts réels du commerce extérieur camerounais, puis ses faiblesses structurelles et les perspectives pour les surmonter.

\textbf{Première partie : des atouts indéniables.} Le Cameroun dispose d'une gamme diversifiée de produits d'exportation : pétrole brut (premier produit en valeur), cacao (parmi les premiers producteurs mondiaux), bois, café, coton, bananes, aluminium. Sa position de carrefour en Afrique centrale en fait la porte d'entrée des pays enclavés : le port de Douala, par lequel transite l'essentiel du commerce extérieur national, dessert aussi le Tchad et la Centrafrique via le corridor Douala--N'Djamena et la voie ferrée Transcamerounaise. L'appartenance à la CEMAC (six États, monnaie commune, le franc CFA) offre un cadre d'intégration régionale, et les marchés frontaliers (Banki, Kyé-Ossi, Abam-Minko, Mbaïboum, Ekondo-Titi) témoignent d'un commerce de proximité dynamique. Enfin, le pays dispose de partenaires solides et diversifiés : Chine, Union européenne, Inde.

\textbf{Deuxième partie : des faiblesses structurelles, mais des perspectives.} Ces échanges souffrent d'abord d'un déficit chronique de la balance commerciale (environ $-\num{1800}$ milliards de FCFA en 2021, ordre de grandeur) : les importations dépassent largement les exportations, avec un taux de couverture voisin de 58 \%. Ensuite, la détérioration des termes de l'échange oblige le pays à exporter toujours plus de matières premières pour importer la même quantité de biens manufacturés. Le commerce reste extraverti (tourné vers l'Asie et l'Europe, très peu vers l'Afrique et la CEMAC), dominé par des produits bruts à faible valeur ajoutée, et freiné par les tracasseries douanières, la faiblesse des infrastructures et l'insécurité sur certains axes. Les perspectives existent pourtant : la transformation locale des matières premières (chocolaterie, sciage, raffinage, coton-textile), l'approfondissement de l'intégration régionale et l'entrée en vigueur de la ZLECAf (Zone de libre-échange continentale africaine), ainsi que la modernisation du port de Kribi, peuvent rééquilibrer les échanges.

\textbf{Conclusion.} Les échanges extérieurs du Cameroun reposent sur des atouts réels — ressources variées, position de carrefour, port de Douala — mais demeurent prisonniers d'une structure extravertie et déficitaire. Répondre à la problématique revient donc à transformer avant d'exporter : c'est tout l'enjeu de l'industrialisation du Cameroun, que le pays s'est fixée à l'horizon 2035 dans sa vision d'émergence.

\textbf{Barème indicatif (sur 20) :} introduction (accroche, définitions, problématique, plan) 4 pts ; première partie argumentée et chiffrée 5 pts ; deuxième partie (faiblesses + perspectives) 7 pts ; conclusion et ouverture 2 pts ; qualité de l'expression et organisation 2 pts.
\textbf{Points de vigilance :} la problématique doit être une vraie question ; chaque partie doit contenir au moins trois arguments illustrés d'exemples camerounais précis (produits, lieux, chiffres) ; évite le catalogue sans liens logiques et soigne les transitions entre les parties.
\end{corrige}

\newpage

\coursec{Fiche bilan}

\begin{popfigure}
\begin{center}
\begin{tikzpicture}[mindmap, grow cyclic,
  every node/.style={concept, font=\small},
  root concept/.append style={concept color=popblue, font=\bfseries\small, minimum size=2.6cm},
  level 1 concept/.append style={sibling angle=60, level distance=3.6cm, minimum size=2.1cm},
  level 2 concept/.append style={sibling angle=45, level distance=2.4cm, minimum size=1.7cm, font=\scriptsize}]
\node[root concept]{Échanges du Cameroun avec l'extérieur}
  child[concept color=poporange]{ node {Notions clés}
    child[concept color=poporangeL]{ node {Balance commerciale} }
    child[concept color=poporangeL]{ node {Termes de l'échange} }
  }
  child[concept color=popgreen]{ node {Produits échangés}
    child[concept color=popgreenL]{ node {Export : pétrole, cacao, bois} }
    child[concept color=popgreenL]{ node {Import : biens d'équipement} }
  }
  child[concept color=popblue]{ node {Partenaires africains}
    child[concept color=popblueL]{ node {CEMAC : 6 États} }
    child[concept color=popblueL]{ node {Nigeria voisin} }
  }
  child[concept color=poppurple]{ node {Reste du monde}
    child[concept color=poppurpleL]{ node {UE, Chine} }
    child[concept color=poppurpleL]{ node {APE, AGOA} }
  }
  child[concept color=poppink]{ node {Problèmes}
    child[concept color=poppinkL]{ node {Déficit, dépendance} }
    child[concept color=poppinkL]{ node {Détérioration des termes} }
  }
  child[concept color=popgold]{ node {Marchés frontaliers}
    child[concept color=popgoldL]{ node {Banki, Kyé-Ossi} }
    child[concept color=popgoldL]{ node {Flux, impacts, obstacles} }
  };
\end{tikzpicture}
\end{center}
\legende{Carte mentale de la leçon : les six branches à maîtriser pour le Bac (notions, produits, partenaires africains, reste du monde, problèmes, marchés frontaliers).}
\end{popfigure}

\begin{bilan}
\textbf{Définitions clés à connaître par cœur}
\begin{itemize}
  \item \cle{Balance commerciale} : différence entre la valeur des exportations et celle des importations d'un pays sur une période donnée. Elle est \emph{excédentaire} si exportations $>$ importations, \emph{déficitaire} dans le cas contraire.
  \item \cle{Termes de l'échange} : rapport entre l'indice des prix à l'exportation et l'indice des prix à l'importation. Leur \emph{détérioration} signifie que les exportations rapportent de moins en moins par rapport aux importations.
  \item \cle{CEMAC} : Communauté Économique et Monétaire de l'Afrique Centrale, regroupant 6 États (Cameroun, Tchad, Centrafrique, Gabon, Congo, Guinée équatoriale), avec une monnaie commune, le franc CFA (BEAC).
  \item \cle{Marché frontalier} : lieu d'échanges commerciaux situé à proximité d'une frontière, où circulent des marchandises entre pays voisins, souvent en partie informelles.
\end{itemize}

\textbf{Repères chiffrés et formules}
\begin{center}
\begin{tabularx}{\linewidth}{|l|X|}\hline
\rowcolor{popblueL} \textbf{Grandeur} & \textbf{Formule ou ordre de grandeur} \\ \hline
Solde commercial & Exportations $-$ Importations \\ \hline
Termes de l'échange & (Indice des prix à l'exportation / Indice des prix à l'importation) $\times$ 100 \\ \hline
Taux de couverture & (Exportations / Importations) $\times$ 100, en \% \\ \hline
Part d'un partenaire & (Échanges avec le partenaire / Échanges totaux) $\times$ 100 \\ \hline
1\textsuperscript{er} produit d'exportation & Pétrole brut (environ 40 \% des recettes, ordre de grandeur) \\ \hline
Principaux clients & Union européenne, Chine, Inde, Nigeria \\ \hline
\end{tabularx}
\end{center}

\textbf{Méthodes express}
\begin{itemize}
  \item \emph{Commenter un tableau d'échanges} : repérer les 3 premiers partenaires, calculer les parts en \%, comparer exportations et importations, conclure sur le solde.
  \item \emph{Construire un croquis de flux} : contour simplifié du Cameroun, flèches d'épaisseur proportionnelle aux flux, légende numérotée, titre, orientation.
  \item \emph{Étudier un marché frontalier} : localiser, inventorier les produits échangés, évaluer les impacts (revenus, emplois, approvisionnement) puis les difficultés (fraude, tracasseries, insécurité).
\end{itemize}
\end{bilan}

\begin{aretenir}[Checklist avant le Bac]
\begin{itemize}
  \item $\square$ Je sais définir la balance commerciale et distinguer excédent et déficit.
  \item $\square$ Je sais calculer et interpréter les termes de l'échange et le taux de couverture.
  \item $\square$ Je connais les principaux produits exportés (pétrole, cacao, café, bois, coton, banane, aluminium) et importés (biens d'équipement, produits alimentaires, hydrocarbures raffinés).
  \item $\square$ Je cite les 6 États membres de la CEMAC et son siège (Bangui).
  \item $\square$ Je connais les grands partenaires du Cameroun hors Afrique (Union européenne, Chine).
  \item $\square$ Je sais expliquer la détérioration des termes de l'échange et ses causes (baisse des cours des matières premières, hausse des prix des produits manufacturés importés).
  \item $\square$ Je sais localiser au moins trois marchés frontaliers : Banki, Mbaïboum, Abam-Minko, Kyé-Ossi, Ekondo-Titi.
  \item $\square$ Je sais décrire les flux de marchandises et les impacts économiques et sociaux d'un marché frontalier.
  \item $\square$ Je sais citer les difficultés des échanges transfrontaliers (fraude, contrebande, tracasseries, état des routes, insécurité).
  \item $\square$ Je sais construire un croquis de flux avec titre, légende numérotée et flèche du nord.
  \item $\square$ Je sais calculer une part en pourcentage et un taux de variation à partir d'un tableau statistique.
  \item $\square$ Je rédige un commentaire de document structuré : présentation, analyse, conclusion.
\end{itemize}
\end{aretenir}

\findecours

\end{document}
